% 本文件由 LaTeX 排版 Agent 根据有效 translation.json 自动生成。
% author=Leo the Great, Pope (c. 395-461); work_id=3603; title=圣大良讲道集

\chapter{圣大良讲道集}

% structure: p0001 source=h1 target=omitted reason=page-title-already-rendered-by-parent

讲道1\newline{}讲道2\newline{}讲道3\newline{}讲道9\newline{}讲道10\newline{}讲道12\newline{}讲道16\newline{}讲道17\newline{}讲道19\newline{}讲道21\newline{}讲道22\newline{}讲道23\newline{}讲道24\newline{}讲道26\newline{}讲道27\newline{}讲道28\newline{}讲道31\newline{}讲道33\newline{}讲道34\newline{}讲道36\newline{}讲道39\newline{}讲道40\newline{}讲道42\newline{}讲道46\newline{}讲道49\newline{}讲道51\newline{}讲道54\newline{}讲道55\newline{}讲道58\newline{}讲道59\newline{}讲道62\newline{}讲道63\newline{}讲道67\newline{}讲道68\newline{}讲道71\newline{}讲道72\newline{}讲道73\newline{}讲道74\newline{}讲道75\newline{}讲道77\newline{}讲道78\newline{}讲道82\newline{}讲道84\newline{}讲道85\newline{}讲道88\newline{}讲道90\newline{}讲道91\newline{}讲道95\par

\SourceRecord{Translated by Charles Lett Feltoe.；Nicene and Post-Nicene Fathers, Second Series；Vol. 12.；Edited by Philip Schaff and Henry Wace.；Buffalo, NY: Christian Literature Publishing Co.,；1895.；<http://www.newadvent.org/fathers/3603.htm>.}

% structure: p0001 source=h1 target=section reason=page-title

\section{讲道一}

\Emph{在他的诞辰或晋铎日讲道。}\par

% structure: p0003 source=h2 target=subsection reason=relative-heading-rank; base=h2

\subsection{当他在缺席中被选，他感谢善意并热切请求他的教会祈祷。}

\Quote{愿我的口述说对主的赞美，} 愿我的气息和心灵、我的肉体和舌头赞美祂的圣名。因为对天主的恩惠保持沉默，不是谦逊而是忘恩的标志：并且我们作为祝圣的教宗，以对主的赞美之祭开始我们的职责，是十分合宜的。因为 \Quote{在我们的卑微中} 主 \Quote{顾念了我们}，并祝福了我们：因为 \Quote{唯独祂为我行了伟大的奇事，} 以致你们对我的圣爱把我算作在场，尽管漫长的旅程迫使我缺席。因此，我感谢并永远感谢我们的天主，为祂所赏赐我的一切。我也公开承认你们对我的好感，向你们致以我应得的感谢，并以此表明我明白，你们出于爱的热忱，能够为我这个以牧者焦虑渴望你们灵魂安全、在你们身上作出了如此良心的判断的人，付出多少尊重、爱和忠诚。因此，我因主的慈悲恳求你们，以你们的祈祷帮助那个你们用恳求寻找出来的他，好使恩宠之神常驻于我，并使你们的判断不致改变。愿那感动你们如此一致旨意的那位，赐给我们众人共同的平安祝福：如此，在我一生中，随时准备好为全能的天主服务，并尽我向你们的职责，我就可以满怀信心地向主恳求：\BibleQuote{圣父，求祢因祢的名保全那些祢所赐给我的人\BibleRef{若 17:11}：} 当你们不断迈向救恩时，愿 \BibleQuote{我的灵魂颂扬上主\BibleRef{路 1:46}，} 在将来审判的报应中，愿我的司铎职的账目如此向公义的审判者呈报，以致通过你们的善行，你们成为我的喜乐和冠冕，你们以善意在此生给了我一个热切的见证。\par

\SourceRecord{Translated by Charles Lett Feltoe.；Nicene and Post-Nicene Fathers, Second Series；Vol. 12.；Edited by Philip Schaff and Henry Wace.；Buffalo, NY: Christian Literature Publishing Co.,；1895.；<http://www.newadvent.org/fathers/360301.htm>.}

% structure: p0001 source=h1 target=section reason=page-title

\section{讲道集 第二篇}

\Emph{论他的诞辰，第二篇：在晋牧周年纪念日宣讲。}\par

% structure: p0003 source=h2 target=subsection reason=relative-heading-rank; base=h2

\subsection{一、主扶起软弱者，并按照他的需要赐予恩宠}

天主的俯就使这一天为我成为可敬的日子，因为他藉着将我这卑微之人提拔至最高品位，显示了他绝不轻视自己的任何一个人。因此，尽管人必须对自己的功德存有戒惧，却仍当尽本分因这恩赐而喜乐，因为那加给人重担的，正是那帮助人完成这重担的；并且，为使软弱的领受者不至于因恩宠的伟大而跌倒，那赐予尊位者也将赐予能力。因此，当主所预定我应开始主教职务的日期按时返回时，我确有理由为天主的荣耀而喜乐，因为他为使我多多爱他，而多多宽恕了我；为彰显他的恩宠的神奇，将他的恩赐赐予了在他眼中毫无功德可言之我。藉着这工作，主向我们的心提示并嘱咐的，不就是叫谁也不可依仗自己的正义，也不可怀疑天主的仁慈吗？这仁慈正是在罪人被圣化、卑微者被高举时格外彰显的。因为天上恩赐的尺度并不取决于我们行为的好坏，并且在这个\Quote{整个人生都是考验，}的世界中，各人并非按他的功德受赏，因为主若细察人的过犯，谁能在他的审判台前站立得住呢？\par

% structure: p0005 source=h2 target=subsection reason=relative-heading-rank; base=h2

\subsection{二、主教的盛大集会证明人们藉着伯多禄的不肖继任者而忠诚地接受了伯多禄}

因此，亲爱的诸位，\BibleQuote{请你们与我一同赞扬主，让我们齐声颂扬他的名号}，好使今日集会的全部原因归于那成就这事者的赞美。因为就我个人的感受而言，我承认我最因你们众人的热忱而喜乐；当我瞻仰这由我可敬的弟兄——司铎们所组成的辉煌集会时，我感到，在许多圣人聚集的地方，天使们也就在我们中间。我也不怀疑，我们今天蒙受了更丰沛的神性临在的倾注，因为这么多天主的华丽帐幕、基督身体的这么多杰出肢体聚集在一处，发出同一光芒。而且，我有把握，最蒙福的宗徒伯多禄的慈爱俯就和真诚之爱并未缺席于这会众：他没有舍弃你们的虔诚，你们正是为尊敬他而聚集。因此，他也乐于接受你们对他的情感，并欢迎你们对主自己的制度所表示的尊敬——这尊敬是向着与他分享尊位者而表达的，他称赞整个教会的井然有序的爱，这爱常是在伯多禄的座位上找到伯多禄，并且因对这位伟大牧者的爱，即使对我这样一个如此卑下的继任者也不冷淡。因此，亲爱的诸位，为使你们一致对我这卑微之人所显示的忠诚获得热忱的果实，请屈膝恳求我们天主的仁慈良善：愿他在我们的日子里驱逐那些攻击我们的人，坚固信德，增加爱德，广赐平安；并愿他屈尊使我——他可怜的仆人，他为了显示他恩宠的富饶，愿意我站在教会的舵轮前——足以承担如此伟大的工作，对建造你们有所助益；为此目的，愿他延长我们服务的时日，使我们蒙他赏赐的年岁能用以光荣他，藉着我们的主基督。阿们。\par

\SourceRecord{Translated by Charles Lett Feltoe.；Nicene and Post-Nicene Fathers, Second Series；Vol. 12.；Edited by Philip Schaff and Henry Wace.；Buffalo, NY: Christian Literature Publishing Co.,；1895.；<http://www.newadvent.org/fathers/360302.htm>.}

% structure: p0001 source=h1 target=section reason=page-title

\section{讲道三}

\Emph{论他的生日，第三篇：在他被提升至教宗职位的周年纪念日上所讲。}\par

% structure: p0003 source=h2 target=subsection reason=relative-heading-rank; base=h2

\subsection{一、被提升至主教职位的荣耀必须完全归于教会的神圣元首}

每当天主的仁慈恩准带给我们祂恩赐的日子时，亲爱的诸位，只要我们受任此职归于赐予者的赞美，就有正当合理的理由欢乐。因为虽然这种对天主的承认可能出现在所有祂的司铎身上，但我认为它对我特别有约束力，我考虑到自己极其微不足道和所担任职务的伟大，自己也应当发出先知的那声惊呼：\Quote{上主，我听到祢的言语就害怕；我观察祢的作为就惊慌。}因为还有什么比让软弱者劳作、卑微者升高、不配者获尊更不寻常、更令人惊慌的呢？然而我们并不绝望，也不灰心，因为我们不靠自己，而靠那在我们内工作的祂。因此，亲爱的诸位，我们也以和谐的声音咏唱了\Person{达味}的圣咏，不是为赞扬我们自己，而是为光荣主基督。因为关于祂，先知性记载说：\Quote{祢是依照\Person{默基瑟德}的品位永为司祭。}也就是说，不是依照\Person{亚郎}的品位，\Person{亚郎}的司祭职按自己后裔的世系相传，是一种暂时的职务，并随旧约的法律而终止，而是依照\Person{默基瑟德}的品位，在他身上预表了永恒的至高司祭。并且也没有提及他的出身，因为在他身上，我们理解所描绘的是那位，祂的根源无可言说。最后，既然这神圣司祭职的奥秘已下传至人的代理，它不按血统传承，也不选择那由血肉所创造的，而是不顾父权及继承的特权，那些由圣神所预备的人被教会接纳为她的统治者：以致在天主所收养的子民中（其整个身体是司祭的和王家的），获得傅油的不是世俗出身的特权，而是天主的恩宠的屈尊就卑造就了主教。\par

% structure: p0005 source=h2 target=subsection reason=relative-heading-rank; base=h2

\subsection{二、从基督并经由圣伯多禄，司铎职务永久传递}

因此，亲爱的诸位，尽管我们在履行职务时发现既软弱又怠惰，因为无论我们想做什么虔诚而有力的行动，都因我们自身状况的脆弱而受阻碍；然而，由于我们拥有全能和永恒司祭的不断的赎罪祭，祂像我们一样，却与父同等，将祂的天主性降至人性的事物，并将祂的人性升至天主性的事物，我们便值得并虔诚地因祂的安排而欢乐，因为祂虽将照顾祂的羊群委托给许多牧人，却没有放弃对祂所爱羊群的守护。并且从祂至高无上的永恒保护中，我们也获得了宗徒助佑的支持，这助佑肯定不会停止运作；而教会整个上层建筑所建立的根基的坚固性，并未因建在其上的殿宇的重量而减弱。因为在宗徒之长中所受到赞扬的那信德的坚固性乃是永恒的：正如那\Person{伯多禄}在\Person{基督}内所相信的保持不变一样，那\Person{基督}在\Person{伯多禄}内所建立的也保持不变。因为，正如在福音课中所读到的，当主在各种不同的意见中问门徒们他们认为祂是谁，而真福伯多禄回答说：\BibleQuote{祢是基督，永生天主之子。}主就说：\BibleQuote{西满巴尔约纳，你是有福的，因为不是肉和血启示了你，而是我在天之父。我也告诉你，你是伯多禄，在这磐石上，我要建立我的教会，阴间的门不能战胜她。我要将天国的钥匙交给你；凡你在地上束缚的，在天上也要被束缚；凡你在地上释放的，在天上也要被释放。}\BibleRef{玛 16:16}\par

% structure: p0007 source=h2 target=subsection reason=relative-heading-rank; base=h2

\subsection{三、圣伯多禄的工作仍由他的继承者继续}

真理的安排因此常存，而真福伯多禄，持守他所领受的磐石力量，没有放弃他所承担的教会舵轮。因为他以这样的方式被立在众人之前：从他被称作磐石，从他被宣布为基础，从他被立为天国的看门人，从他被设置为束缚和释放的裁判者（他的判决将在天堂保持效力），从所有这些奥秘的称号，我们可以知道他与\Person{基督}结合的性质。并且今天，他更充分、更有效地执行所交托给他的事，并在那光荣了他的那一位内和与那一位一起，履行他职责和使命的每一部分。因此，若有什么由我们正确做成和正确决定的事，若有什么通过我们每日的恳求从天主的仁慈中获得的事，那是他的工作和功劳的结果，因为他的能力在其宗座上生活，他的权威在其宗座上盛行。因为，亲爱的诸位，这是由那个宣认所获得的，那个宣认由天主父启示在宗徒心中，超越了人间意见的一切不确定性，并被赋予磐石的坚固，任何攻击都不能动摇它。因为在整个教会中，伯多禄每天都说：\BibleQuote{祢是基督，永生天主之子。}而每个承认主的舌头都接受他声音所教导的教训。这信德战胜魔鬼，打破其囚犯的锁链。它将我们从地土拔出，栽种于天上，阴间的门不能战胜它。因为它是被天主赋予如此坚固，以致异端的邪恶不能损害它，外邦人的不信不能克服它。\par

% structure: p0009 source=h2 target=subsection reason=relative-heading-rank; base=h2

\subsection{四、因此，这个庆节是为光荣圣伯多禄，而他的羊群的进步也归于他的光荣}

所以，亲爱的诸君，我们以合理的顺服，藉此方式庆祝今日的庆节，好使在我卑微之人身上，他得以被认出并受尊崇——他内住着对所有牧者的关怀，连同托付给他的群羊的重任，而他的尊荣即使在一个如此不堪的继承者身上也未曾减损。因此，我所如此渴望和珍视的可敬的弟兄及同僚司铎们的在场，将更为神圣和宝贵，如果他们愿意将这服务——他们已屈尊参与——的首席荣耀，转移到他们所深知的那一位身上，他不仅是此教区的主保，更是所有主教的首席。因此，当我们向你们的耳朵讲出劝言时，圣洁的弟兄们，请相信是他在说话，因为我们是他的代表：因为我们所给予的是他的警告，我们所宣讲的唯有他的教导，恳请你们\BibleQuote{束起你们心的腰}\BibleRef{伯前 1:13}，并在敬畏天主中度贞洁和清醒的生活，不要让你的心忘记他的至高无上并屈服于肉体的情欲。此世欢愉的快乐短暂而飞逝，它们试图使那些蒙召得永生的人偏离生命的道路。因此，忠实而虔诚的心灵必须渴求天上的事物，热切盼望天主的恩许，提升自己臻于对不朽之善的爱德和对真光的望德。但是亲爱的诸君，请确信，你们抵抗恶习、对抗肉体情欲的劳苦，在天主眼中是悦纳而珍贵的，并因天主的仁慈，不仅对你们自己有益，也对我有益，因为热心的牧者以主的群羊的进步为夸耀。\BibleQuote{你们就是我的冠冕和喜乐}\BibleRef{得前 2:20}，正如宗徒所说；只要你们的信德——这信德自福音开端就已在普世传扬——能坚持在爱德和圣德之中。因为虽然普世教会——它遍及天下——理应在一切美德上富足，然而你们尤其，超越万民，理应卓越于虔诚善行，因为你们被建立在宗徒磐石的堡垒之上，我们的主耶稣基督不仅与所有人一同救赎了你们，而且有福的宗徒伯多禄远超众人地教导了你们。借着同一基督我们的主。\par

\SourceRecord{Translated by Charles Lett Feltoe.；Nicene and Post-Nicene Fathers, Second Series；Vol. 12.；Edited by Philip Schaff and Henry Wace.；Buffalo, NY: Christian Literature Publishing Co.,；1895.；<http://www.newadvent.org/fathers/360303.htm>.}

% structure: p0001 source=h1 target=section reason=page-title

\section{讲道第九}

\Emph{论募捐，第四}\par

% structure: p0003 source=h2 target=subsection reason=relative-heading-rank; base=h2

\subsection{一、魔鬼的邪恶在引人误入歧途方面，如今被救赎之工所抵消，使其回归真理}

极亲爱的诸位，天主的仁慈与公义，藉我们的主耶稣基督的教导，以慈爱向我们启示了祂报应的方式，正如从创世之初所预定的，好让我们接受事实的意义，将我们所信将要发生的事，视为好像已经发生一样。因为我们的救赎主和救主知道魔鬼的欺骗在世上散布了何等大的谬误，以及它通过多少迷信使人类的大部分屈从于自己。但为了使按天主的肖像受造之物不再因对真理的无知而被驱向永死的悬崖，祂将祂审判的性质插入福音篇章中，好使每个人都能从狡猾仇敌的陷阱中获救；因为如今所有人都将知道善人可期望什么赏报，恶人必须畏惧什么惩罚。因为罪的煽动者和创始者，他先因骄傲而堕落，然后因嫉妒而伤害我们，因为他\Quote{他不在真理中\BibleRef{若 8:44}}，他将所有力量用于说谎，并从他狡猾的有毒源头产生各种欺骗，为要断绝人对那幸福的虔诚希望，那幸福是他因自己的抬高而失去的，并将人拖入与他同受定罪，而他本人无法获得和好。因此，人当中凡因自己的邪恶而得罪天主的，都是被他的诡计引入歧途，被他的恶行所败坏。因为他轻易地将那些在宗教上被他欺骗的人驱入一切恶行。但他知道天主不仅被言语否认，也被行为否认，对于许多他不能夺去信德的人，他夺去了他们的爱德，并用贪婪的杂草堵塞他们心田的土壤，夺去了他们善工之果，尽管他不能夺去他们口唇的宣认。\par

% structure: p0005 source=h2 target=subsection reason=relative-heading-rank; base=h2

\subsection{二、天主公义的审判宣布反对罪恶，为使我们通过慈悲与爱德的行为避免它}

因此，极亲爱的诸位，因着我们远古仇敌的这些狡猾计谋，基督不可言说的良善愿意我们知晓，在报应之日关于全人类将要如何判决，好使在此生中，当合法的治愈药物被提供，当忏悔者不被拒绝修复，那些长久荒芜者终能结果实之时，正义所决定的判决可以被预先阻止，并且天主前来审判世界的景象永不离开心灵的眼目。因为主将如祂自己预言的那样，以祂荣耀的威严来临，并有无数光辉灿烂的天使军团与祂同在。在祂权能的宝座前，世上万民都将被聚集；所有世代、遍及全地所生的人，都要站在审判者面前。那时，义人与不义者、无罪者与有罪者将被分开；当虔诚之子，在审查了他们的慈悲工作之后，接受了为他们预备的国时，不义者将因他们的全然荒芜而被斥责，左边的人与右边的人毫无共同之处，将受全能审判者的定罪，被投入为折磨魔鬼及其天使所预备的火中，与牠同受惩罚，因为他们选择了遵行牠的意愿。那么，有谁不会在这永恒折磨的判决前战栗呢？有谁不会恐惧那永不止息的灾祸呢？但既然这严厉被宣布仅是为使我们寻求仁慈，我们在此生也必须展示如此慷慨的仁慈，以致在危险的疏忽之后，重返虔诚的工作，我们可能得以从这判决中获得释放。因为这就是审判者权能和救主恩慈的目的：使不义者离弃他的道路，罪人放弃他的邪恶习惯。愿那些希望基督宽恕他们的人，怜悯穷人；愿他们慷慨施舍，喂养那些渴望达到真福者团体的人。愿没有人认为自己的同僚卑贱，也不要轻视任何人身上那造物主使之成为祂自己的本性。因为哪个劳苦的人能否认基督将这劳苦视为做在祂自己身上？你的同仆因此得到帮助，但酬报的是主。周济穷人是天国的赎价，慷慨分施现世之物的人成了永恒之物的继承人。但如此微小的花费何以配得如此高的估价，除非因为我们的工作是在爱的天平上衡量的，并且当一个人爱天主所爱之时，他当得起被提升到祂的国中，因为爱的属性已部分成为他的？\par

% structure: p0007 source=h2 target=subsection reason=relative-heading-rank; base=h2

\subsection{三、我们在祂穷人的身上服事基督自己}

为此，亲爱的诸位，宗徒制度所确立的日子邀请我们履行这行善的虔诚本分。在这日子上，教父们明智而有益地规定了首次收集我们神圣的奉献；目的是：昔日外邦人曾在此季节迷信地供奉魔鬼，如今我们却以施舍的神圣奉献来对抗不虔敬者的邪恶祭献。由于这对教会的成长极为有益，故决定使之成为永久的制度。因此，我们劝勉你们，圣洁的弟兄们，在各地区的教会中，于下周三按各自的能力和意愿，为爱德的目的献出你们的财物，好使你们能获得那福乐——\Quote{关心穷人和弱小者}的人将在其中永享喜乐。若我们要\Quote{关心}他，亲爱的诸位，就必须以爱心的关怀和警醒，好能找到那被谦逊遮掩、被羞耻心所阻拦的人。因为有些人羞于公开乞求所需，宁愿默默忍受贫困，也不愿因公开求助而蒙羞。这些人正是我们应当\Quote{关心}的，并应解除他们隐藏的困境，好使他们因自己的羞怯和贫困同时得到照顾而更加喜乐。我们在穷人和弱小者身上正确地认出我们的主耶稣基督本人，正如蒙福的宗徒所说：\Quote{祂本是富有的，却成了贫穷的，好使你们因祂的贫穷而成为富有的。}\BibleRef{格后 8:9}为使祂的同在永不离开我们，祂如此实现了祂的谦卑与荣耀的神妙结合：我们既在圣父的尊威中朝拜祂为王为主，也在祂的穷人身上喂养祂；为此，在邪恶的日子我们必得脱永罚，并因对穷人的细心照料而与整个天乡的团体结合。\par

% structure: p0009 source=h2 target=subsection reason=relative-heading-rank; base=h2

\subsection{四、为使他们的奉献完全蒙天主悦纳，不可忽略检举城中的摩尼教徒}

为使你们的虔诚，亲爱的诸位，在一切事上蒙天主悦纳，我们也劝勉你们，要热心地向你们的司铎告发隐藏在任何地方的\OriginalTerm{摩尼教徒}{Manichees}。因为揭发恶人的藏身处，并在他们身上推翻他们所侍奉的魔鬼，这无非是虔敬之举。亲爱的诸位，对抗他们，确实全世界的整个教会都当处处穿上信德的盔甲；但你们的虔诚当在这方面走在前列，因为你们在先辈身上，直接从最蒙福的宗徒伯多禄和保禄的口中，学习了基督十字架的福音。决不能让那些人隐藏起来，他们不相信藉梅瑟所颁布的法律（其中显示天主是宇宙的创造者）应当被接受；他们出言反对先知和圣神，竟敢以可咒诅的亵渎拒绝那在普世教会中备受恭敬的达味圣咏，否认主基督按肉身的诞生，声称祂的苦难与复活是虚构而非真实的，并剥夺重生圣洗圣事作为恩宠媒介的一切能力。在他们那里，没有什么是神圣的，没有什么是完整的，没有什么是真实的。他们当被远离，以免伤害任何人；他们当被交出，以免在我们城中任何地方定居。亲爱的诸位，我们如此吩咐、如此请求，在主审判台前，这将是你们的收益。因为理当将这行为的胜利也与我们的施舍奉献相结合，主耶稣基督在一切事上帮助我们，祂永生永王，万世无疆。阿们。\par

\SourceRecord{Translated by Charles Lett Feltoe.；Nicene and Post-Nicene Fathers, Second Series；Vol. 12.；Edited by Philip Schaff and Henry Wace.；Buffalo, NY: Christian Literature Publishing Co.,；1895.；<http://www.newadvent.org/fathers/360309.htm>.}

% structure: p0001 source=h1 target=section reason=page-title

\section{讲道第十篇}

\Emph{论募捐，第五篇}\par

% structure: p0003 source=h2 target=subsection reason=relative-heading-rank; base=h2

\subsection{一、我们的财物并非私有，而是为天主服务之用}

亲爱的诸位，我们遵循宗徒传统的规诫，如同警醒的牧者劝勉你们：要以虔敬的实践热忱庆祝那一日——他们既已清除邪恶的迷信，并将之奉献于仁慈的事工，由此表明教父的权威仍活在我们中间，我们亦顺从地坚守他们的教导。既然这等实践的神圣益处不仅影响过去，也惠及我们这时代，使昔日助他们破除虚妄者，如今助我们增长德行。还有什么比以下之事更适合信德、更契合虔敬呢？即：救助穷困者的匮乏，照顾弱者的需要，扶助弟兄们的急难，并在他人的劳苦中记念自身的处境。在此事业中，唯有那位洞悉各人所领受之物的主，能正确分辨一个人能做什么、不能做什么。因为不仅属灵的财富和天上的恩赐来自天主，连地上的、物质的产业也出自祂的慷慨，为叫祂有理由要求我们交账——这些事物，祂与其说是赐予我们为私有，不如说是委托我们管理。因此，我们必须正当而明智地使用天主的恩赐，免得行善的素材反成为犯罪的机缘。因为财富，就其种类并作为手段而言，原是好的，且对人类社会大有裨益，只要它掌握在仁慈慷慨之人手中，不被奢侈者挥霍，也不被吝啬者囤积；因为无论存放不当或使用愚昧，财富同样丧失。\par

% structure: p0005 source=h2 target=subsection reason=relative-heading-rank; base=h2

\subsection{二、若只为私己目的，滥用财富比无用更糟}

虽然逃离无节制、避免卑劣享乐的浪费是可赞许的，并且许多人在其豪阔中不屑隐藏财富，因其财物丰裕而轻视卑鄙小气的悭吝，然而这样的人其慷慨并不幸福，其节俭也不值得称赞——如果他们的财富只对自己有益，如果没有穷人因他们的财物得到帮助，没有病人得到滋养，如果从其丰裕的大量财富中，俘虏得不到赎金，旅客得不到安慰，流亡者得不到救济。这类富人比所有穷人都更贫穷。因为他们失去了可能永久拥有的回报，当他们陶醉于对自己财物短暂且并非总是自由的享用时，他们没有被正义之粮或仁慈之甘饴喂养：外表辉煌，内心无光；拥有丰富的时间之物，却全然欠缺永恒之物：因为他们使自己的灵魂忍饥挨饿，并带它们陷入羞耻和赤贫，不肯将任何放入他们地上仓库的东西用于天上宝库。\par

% structure: p0007 source=h2 target=subsection reason=relative-heading-rank; base=h2

\subsection{三、仁慈之责超越所有其他德行}

但或许有些富人，虽然不惯于以丰厚的捐赠帮助教会的穷人，却遵守天主的其他诫命，并以为在他们许多有功劳的信德与正直行为中，会因缺乏这一德行而得到宽恕。然而这一德行如此重要，以至于即使其他德行存在，没有它也无济于事。因为纵然一个人充满信德、贞洁、清醒，并饰有其他更大的美德，但若他不仁慈，就不能得到仁慈：因为主说：\Quote{仁慈的人是有福的，因为天主要仁慈地对待他们。\BibleRef{玛 5:7}}当人子在自己的光荣中来临，坐在祂光荣的宝座上，万民聚集，善人与恶人分别时，站在右边的人将因什么受赞颂？岂不正是因仁爱的工作和爱德的行为，而耶稣基督将这些视同做在祂自己身上？因为祂既取了人性，便丝毫不与人的卑微分离。至于左边的人，将受什么责难？岂不正是因他们忽略爱德，无人性的冷酷，拒绝向穷人施以仁慈？仿佛右边的人没有别的德行，左边的人没有别的过失。但在那伟大而最终的审判之日，慷慨的宽大与不敬的吝啬将被视为如此重要，以致压倒所有其他德行和所有其他缺失，以至于前者使人进入天国，后者使人被投入永火。\par

% structure: p0009 source=h2 target=subsection reason=relative-heading-rank; base=h2

\subsection{四、其效能，如圣经所证，不可估量}

所以，亲爱的诸位，谁也不要因任何善功而自满，若缺乏爱德的行为，也不要信赖自己身体的洁净，若未因施舍的净化而得洁除。因为\Quote{施舍消除罪恶，}杀死死亡，并扑灭永火的刑罚。但那在此方面没有结果实的人，将不会从伟大的报偿者那里得到宽恕，正如撒罗满所说：\Quote{掩耳不听弱小哀求的，他呼唤上主，也没有人应允\BibleRef{箴 21:13}。}同样，多俾亚也劝诫他的儿子敬虔的规诫，说：\Quote{拿出你的财产施舍，不要转面不顾任何穷人；这样，天主的圣容也不会转离你。}这个美德使一切美德有益；因为它的临在赋予那信德本身生命，\Quote{义人因信德而生活\BibleRef{哈 2:4}，}而这信德被称为\Quote{没有行为是死的\BibleRef{雅 2:26}：}因为正如行为的缘由在于信德，同样，信德的力量在于行为。所以\Quote{我们一有机会，}如同宗徒所说，\Quote{就该向众人行善，尤其应向有同样信德的人。}\Quote{我们行善不要厌倦；因为到了时候，我们要收割。}因此，今生是播种的时候，报应的日子是收割的时候，那时各人将按自己播种的多少收割所种之果。谁也不会在那次收割的产物中失望，因为所计量的不是花费的数目，而是心中的意向。小力量的小数目与大力量的大数目产生同样效果。所以，亲爱的诸位，让我们实践这项宗徒的制度。既然第一次募捐将在下个主日，愿众人都预备自己乐意奉献，使每人按自己的能力参与这最神圣的献礼。你们的施舍和那些由你们的礼物得到帮助的人，将为你们转求，使你们常准备好行一切善事，因我们的主耶稣基督，祂永生永王，至于无穷之世。阿们。\par

\SourceRecord{Translated by Charles Lett Feltoe.；Nicene and Post-Nicene Fathers, Second Series；Vol. 12.；Edited by Philip Schaff and Henry Wace.；Buffalo, NY: Christian Literature Publishing Co.,；1895.；<http://www.newadvent.org/fathers/360310.htm>.}

\SourceRecord{Translated by Charles Lett Feltoe.；Nicene and Post-Nicene Fathers, Second Series；Vol. 12.；Edited by Philip Schaff and Henry Wace.；Buffalo, NY: Christian Literature Publishing Co.,；1895.；<http://www.newadvent.org/fathers/360312.htm>.}

% structure: p0001 source=h1 target=section reason=page-title

\section{讲道词 十六}

\Emph{论第十个月的斋戒。}\par

% structure: p0003 source=h2 target=subsection reason=relative-heading-rank; base=h2

\subsection{一、富裕者必须通过慷慨施予穷人和困苦者，向天主表达感恩}

天主恩宠的超越大能，亲爱的诸位，确实每天都在基督徒心中促使我们将一切愿望从属地事物转向属天事物。但今世的生活也藉造物主的助佑而度过，并赖祂的眷顾而维系，因为那应许永恒之事的也是供应暂时之事的主。因此，我们理当因着对未来幸福的希望（我们凭信德向此希望奔跑）而感谢天主，因祂提升我们体认那为我们存留的幸福；同样，对那些我们在每年过程中领受的事物，也应尊崇并赞美天主，祂从起初就赐给大地丰产，并为每种籽实和种子规定了结果实的律则，绝不会离弃祂自己的法令，而是作为造物主，永远继续祂对所造之物仁慈的管理。因此，凡田亩、葡萄园和橄榄园为人的用途所出产的，这一切都是天主在祂慷慨的良善中赐予的：因为在元素千变万化的条件下，祂仁慈地协助了农夫们不确定的辛劳，使风、雨、寒冷、炎热、白昼和黑夜都为我们所需而效力。因为人的方法不足以使他们的工作产生效果，除非天主在他们惯常的种植和浇灌上赐予增长。因此，虔敬而公义的是，我们也应把天父仁慈赐予我们的，用来帮助他人。因为有许多人没有田地、没有葡萄园、没有橄榄园，我们必须用天主所赐的财富来供给他们的需求，使他们也能与我们一同因大地的丰饶而赞美天主，并因拥有者所领受的（他们也与穷人和外方人分享了）而欢欣。那仓廪是有福的、最当称赞的，内中各种果实能更加增多，因为饥饿者和软弱者的需要从中得到满足，陌生人的缺乏从中得到缓解，病人的愿望从中获得满足。因为这些人，天主公义地允许他们遭受各种磨难，为要既因忍耐而加冕受苦者，也因仁慈而加冕怜悯者。\par

% structure: p0005 source=h2 target=subsection reason=relative-heading-rank; base=h2

\subsection{二、施舍和斋戒是祈祷最重要的辅助}

而当所有季节都适合此项责任时，亲爱的诸位，目前的季节尤其适宜且恰当，在此时期，我们神圣的祖先蒙受神圣灵感，确立了第十个月的斋戒，为使我们当所有庄稼的收获完成后，能向天主奉献我们合理的节制之祭，并且每个人都要记得如此使用他的富足，即在自己身上更节制，对穷人更慷慨。因为罪过的赦免是藉施舍和斋戒最有效地祈求的，而由这些辅助翼助的恳求迅速上达天主的耳中：因为正如经上所载：\BibleQuote{仁慈的人善待自己的灵魂\BibleRef{箴 11:17}，} 而没有什么比用在邻人身上的更属于一个人自己。因为他用来施予穷人的那部分物质财富，转化为永恒的财富，并且这样的财富出自这种慷慨，永不会减少，也绝不会以任何方式被毁坏，因为\BibleQuote{怜悯人的人是有福的，因为他们要蒙受天主的怜悯\BibleRef{玛 5:7}，} 而祂自己将是他们最大的赏报，因为祂是自己命令的典范。\par

% structure: p0007 source=h2 target=subsection reason=relative-heading-rank; base=h2

\subsection{三、基督徒的虔敬活动如此激怒了撒殚，以致他增多异端来加害他们}

但在这一切敬虔行为（它们使我们越来越蒙天主悦纳）面前，亲爱的诸位，我们的仇敌——他是如此急切且精通于加害我们——无疑被更尖锐的仇恨之刺所激怒，为要藉着伪装的基督徒之名来败坏那些他不能以公开的流血迫害所攻击的人，为此目的，他役使了异端者，这些异端者被他从公教信仰中引入歧途，使之臣服于自己，并在各种错误之下强迫他们为他的阵营服务。正如他为欺骗原始人而使用了蛇的服务，同样，为误导正直之人的心智，他用自己谎言的毒液武装了这些人的口舌。但这些奸诈的阴谋，亲爱的诸位，我们将以牧者的关怀，并尽主所赐予的帮助，予以击败。并且，为谨慎起见，免得圣洁羊群中的任何一只丧亡，我们以父亲的警告劝诫你们远离\Quote{说谎的嘴唇}和\Quote{欺诈的舌头}，先知求主使他的灵魂脱离这些；因为\Quote{他们的言语，}正如蒙福的宗徒所说，\BibleQuote{如同毒疮一样蔓延\BibleRef{弟后 2:17}。} 它们谦卑地爬行，轻柔地捕捉，温和地束缚，隐秘地杀害。因为\Quote{他们来到，}正如救主所预言的，\BibleQuote{披着羊皮，里面却是抢夺的狼\BibleRef{玛 7:15}；} 因为他们若不以基督的名字掩盖其兽性的愤怒，就无法欺骗真实而单纯的羊。但在他们众人身上，那一位在运行，他虽真是光明的仇敌，却\BibleQuote{装作光明的天使\BibleRef{格后 11:14}。} 他的诡计激发了巴西里德；他的巧计在玛尔西翁身上运作；他是撒伯里乌行动的领袖；他是福提诺跌倒的制造者；亚略和欧诺米所事奉的是他的权威和灵；总之，在他的命令和权威之下，这群野兽的整个群体已从教会的合一中分离，并与真理断绝了联系。\par

% structure: p0009 source=h2 target=subsection reason=relative-heading-rank; base=h2

\subsection{四、在所有异端中，摩尼教是最坏、最污秽的}

然而，当他（撒殚）对所有异端保持着这种变化无常的统治权时，他却将他的城堡建在摩尼教人的疯狂之上，并在他们那里找到了最宽敞的庭院，好让他耀武扬威、自我夸耀：因为他在那里不仅拥有一种形式的错误信仰，而是所有错误和不敬神的大杂烩。因为外邦人所有偶像崇拜的、肉体犹太人所有盲目的、魔术技艺秘密中所有不合法的、最后所有异端中所有亵渎神明的东西，都连同各种污秽聚在这些人的身上，如同在一个污水池中一样。因此，要描述他们所有的不敬神之处，实在太过冗长：因为指控他们的数目超过了我的言辞供应。只需指出几个例子就够了，使你们能从所听到的，推断出我们因谦逊而略去的部分。然而，关于他们的仪式，这些仪式在道德上和在宗教上一样不体面，我们不能对主乐意启示给我们调查之事保持沉默，免得有人以为我们在这件事上相信了模糊的谣言和不确定的意见。因此，当主教和司铎坐在我旁边，基督徒贵族也聚集在同一地方时，我们命令他们的选民——男女——被带到我们面前。当他们对自己的歪曲教义和举行节期的方式作了许多披露后，也揭露了这个极其堕落的故事，我羞于描述，但这个故事已被如此仔细地调查，以致怀疑或吹毛求疵者没有留下任何怀疑的余地。因为所有犯下这不可言说的罪行的人都在场，即一个最多十岁的女孩，以及两个抚养她并为这次暴行准备她的妇女。那个侵犯她的少年也在场，还有安排他们可怕罪行的主教。所有这些人都作了同一招供，并披露了如此污秽的秘仪故事，我们的耳朵几乎无法忍受。为了避免更直白的言语冒犯贞洁的耳朵，只需说明这些诉讼记录就够了，其中充分表明，在那个教派中，既找不到任何谦逊，也找不到任何荣誉感，更找不到任何贞洁：因为他们的法律是谎言，他们的宗教是魔鬼，他们的祭献是无耻。\par

% structure: p0011 source=h2 target=subsection reason=relative-heading-rank; base=h2

\subsection{五、每个人都应弃绝这样的人，并将其所知的一切信息禀报当局}

所以，亲爱的诸位，要弃绝与这些完全可憎和有害之人的一切友谊；其他地区的骚乱已将他们大批带进此城。尤其是你们妇女，要避免与这样的人相识和来往，免得你们的耳朵在不知不觉中被他们荒诞的故事所迷惑，落入魔鬼的圈套。魔鬼知道他用女人的口诱惑了第一个人，并藉着女性的轻信将所有的人从乐园的幸福中驱逐出去，所以他仍以更自信的诡计窥伺你们的性别，以便他能用他虚假的仆从所缠住的人，既夺取他们的信德，又夺取他们的贞洁。亲爱的诸位，我也恳切地劝告并忠告你们，如果你们中有人知道他们住在哪里、在哪里教导、常去谁的家、与谁同住，要忠实地告知我们；因为如果有人藉圣神的保护自己没有被他们抓住，但明知别人正在被抓住却毫不作为，那对他几乎没有益处。为对付共同的敌人，为共同的安全，所有人都应同样保持警惕，免得因一个成员受伤，其他成员也受损伤；那些认为不应将这些人交出来的人，尽管他们没有因赞同而被玷污，但他们的缄默在基督的审判中将被判为有罪。\par

% structure: p0013 source=h2 target=subsection reason=relative-heading-rank; base=h2

\subsection{六、铲除异端的热情将使其他虔诚善工更蒙悦纳}

因此，要表现出神圣的宗教警惕热情，让所有信友团结一致起来对抗这些灵魂的残暴敌人。因为仁慈的天主已将一部分有害的敌人交在我们手中，为藉着揭示危险唤起最大的谨慎。不要满足于已经做过的事，而要我们坚持搜查他们；藉天主的助佑，结果不仅是那些仍站立的人继续安全，而且许多被魔鬼诱惑所欺骗的人也将从错误中恢复。你们向仁慈天主献上的祈祷、施舍和斋戒，将因这本身而更加神圣，因为这件信德的行为也加在你们所有其他虔诚善工中。因此，在星期三和星期五，让我们守斋；在星期六，让我们在最蒙福的宗徒伯多禄面前守夜；我们经验并知道，他如同牧人一样不停看守主托付给他的羊群，并将在他的恳求中获胜，使天主的教会——由他的宣讲所建立的——藉我们的主基督脱离一切错误。阿们。\par

\SourceRecord{Translated by Charles Lett Feltoe.；Nicene and Post-Nicene Fathers, Second Series；Vol. 12.；Edited by Philip Schaff and Henry Wace.；Buffalo, NY: Christian Literature Publishing Co.,；1895.；<http://www.newadvent.org/fathers/360316.htm>.}

% structure: p0001 source=h1 target=section reason=page-title

\section{讲道集 第十七篇}

\Emph{论第十个月的斋戒 六}\par

% structure: p0003 source=h2 target=subsection reason=relative-heading-rank; base=h2

\subsection{一、斋戒的职责基于旧约和新约，并与祈祷和施舍的职责密切相关}

法律的教导，亲爱的诸位，为福音的规诫赋予了极大的权威，因为某些事物从旧的规诫转移到新的规诫，并且借着教会本身的虔诚行为表明，主耶稣基督\Quote{不是来废除，而是来成全法律的}\BibleRef{玛 5:17}。既然宣告我们救主来临的标记已经终止，在真理本身面前的预象已被废除，那么我们的宗教为规范风俗或为单纯敬拜天主所设立的那些事物，仍以起初的形式与我们同在，而与两约相符的并未因任何改变而更改。其中也包括第十个月的隆重斋戒，我们现在应按年例守斋，因为向至高的眷顾引导下大地为人类使用所出产的庄稼，向天主的恩惠表示感谢，这是完全正义和虔诚的。为表示我们以乐意的心去做这事，我们不仅应实行斋戒的节制，还应勤于施舍，以便从我们心底也能萌发正义的萌芽和爱德的果实，并因怜悯天主的穷人而配得天主的仁慈。因为由虔诚行为支持的恳求，在天主面前最为有效，因为谁不把心从穷人那里转开，他很快就使自己转向聆听主，如主所说：\Quote{你们应当慈悲，如同你们的父那样慈悲……你们要释放，就必被释放。}有什么比这正义更仁慈？有什么比这报应更慈悲，因为审判者的判决取决于被审判者的权力？他说：\Quote{你们给，就必给你们。}\par

% structure: p0005 source=h2 target=subsection reason=relative-heading-rank; base=h2

\subsection{二、借给主比借给人更合算}

基督徒的施与者，要坚定：给你所能收取的，播下你所能收割的，分散你所能聚集的。不要害怕花费，不要对收益的不确定性叹息。你的财产在明智地支配时就会增长。把你的心放在因怜悯而得的利润上，交易永久的收益。你的酬报者希望你慷慨，那赐予你使你拥有的那一位，命令你花费，说：\Quote{你们给，就必给你们，} \Quote{你们用什么量器量，也要用什么量器量给你们}\BibleRef{路 6:38}。你必须欣然接受这诺言的条件。因为你所拥有的一切没有一样不是你领受的，但是你却不能没有你所施与的。因此，那爱钱财并想通过过分的利润增加财富的人，更应实践这神圣的高利贷，通过这样的借贷致富，不是为了捕捉困于困境中的人，用奸诈的援助使他们陷入永远无法偿还的债务，而是成为他的债权人和放贷者，他说：\Quote{你们给，就必给你们}，\Quote{你们用什么量器量，也要用什么量器量给你们}\BibleRef{路 6:38}。但那不想要永远拥有他所视为渴望之物的人，对自己也是不忠和不公的。任凭他积累多少，任凭他囤积储存多少，他将空手匮乏地离开这个世界，如先知达味所说：\Quote{因为当他死时，什么也不能带去，他的荣耀也不会随他而下。}如果他顾念自己的灵魂，他会把他的财物托付给那一位，他是穷人的适当担保者和慷慨的贷款偿还者。但不义而可耻的贪婪，承诺做某种善行却逃避它，不信任天主的诺言——天主的诺言永不落空，却信任人——人做出如此仓促的交易；而他认为当前比未来更确定，常常理当发现，他不义之财的贪婪成为绝无不义的损失的原因。\par

% structure: p0007 source=h2 target=subsection reason=relative-heading-rank; base=h2

\subsection{三、高利贷在任何方面都是邪恶的}

因此，无论结果如何，放贷者的行业总是邪恶的，因为减少或增加金额都是罪：如果他失去了他所借出的，他是可怜的；如果他拿走的比他借出的多，他就更可怜。必须绝对弃绝高利贷的不义，并躲避缺乏任何仁爱的收益。人的财产确实通过这些不义而可怜的手段增加了，但心灵的财富却减少了，因为金钱的高利贷是灵魂的死亡。至圣先知达味清楚表明了天主对这样的人的看法，当他问：\Quote{上主，谁能住在你的帐幕里？谁能安居在你的圣山上？}时，他从神圣的答复中得知，那遵守其他神圣生活规则而\Quote{没有将金钱放高利贷}的人方能获得永恒的安息。因此，那通过放高利贷获取欺诈收益的人被证明是与天主的帐幕隔绝、并从他的圣山被逐出的人，他试图以别人的损失来致富，理当受到永恒的匮乏的惩罚。\par

% structure: p0009 source=h2 target=subsection reason=relative-heading-rank; base=h2

\subsection{四、让我们远避贪婪，与他人分享天主的恩惠}

因此，亲爱的诸位，你们全心信赖主许诺的人，要逃开这不洁的贪婪癞病，并虔诚明智地使用天主的恩赐。既然你们因他的恩惠而喜乐，要当心使别人也能分享你们的喜乐。因为许多人缺乏你们所富有的，一些人的匮乏为你们提供了效法天主良善的机会，以便通过你们，天主的恩惠也可传递给别人，并且通过明智地管理你们的暂时财物，你们能获得永久的财富。因此，在下周三和周五，让我们斋戒，并在周六与最蒙福的宗徒伯多禄一同守夜，藉着他的祈祷，我们能在一切事上因我们的主耶稣基督而获得天主的保护。阿们。\par

\SourceRecord{Translated by Charles Lett Feltoe.；Nicene and Post-Nicene Fathers, Second Series；Vol. 12.；Edited by Philip Schaff and Henry Wace.；Buffalo, NY: Christian Literature Publishing Co.,；1895.；<http://www.newadvent.org/fathers/360317.htm>.}

% structure: p0001 source=h1 target=section reason=page-title

\section{第十九篇讲道}

\Emph{论第十个月的斋戒（八）}\par

% structure: p0003 source=h2 target=subsection reason=relative-heading-rank; base=h2

\subsection{一、克己导致更高享受}

当救主教导祂的门徒关于天主之国的来临和世界末日的时期，并以宗徒的身份教导整个教会时，祂说：\Quote{你们应当谨慎，免得你们的心因宴饮沉醉及世事的挂虑而变得沉重。\BibleRef{路 21:34}}确实，亲爱的诸位，我们承认这条诫命更特别适用于我们，因为那日子虽然隐藏，却无疑临近了。因为人人都必须预备自己迎接祂的来临，免得祂发现我们沉溺于贪食或纠缠于世事。亲爱的诸位，日常经验证明，肉体的过度放纵会削弱心灵的锐气，饮食过量会迟钝内心的精力，以致饮食的快乐甚至与身体的健康相悖，除非合理的节制抗拒诱惑，并考虑未来的不适而远离快乐。因为尽管肉体没有灵魂便一无所欲，并从同一源头接受感觉和运动，但同一个灵魂的职责是拒绝某些事情，来顺从于它的身体，并凭内在的判断力约束外在部分远离不合时宜的事物，以便更频繁地摆脱肉欲，在心灵的宫殿中享有神圣智慧的闲暇，远离一切尘世喧嚣，静静享受神圣的默想和永恒的喜乐。虽然这在今生难以持守，但可以常常重新努力，以便我们更频繁、更长久地专注于灵性而非肉体的挂虑；通过将越来越多的时间用于更高尚的挂虑，甚至我们现世的行为也能最终获得不朽的财富。\par

% structure: p0005 source=h2 target=subsection reason=relative-heading-rank; base=h2

\subsection{二、一年四期斋戒的教导表明，灵性的克己与肉身的克己同样必要}

亲爱的诸位，这种有益的遵守尤其规定于教会的斋戒中，这些斋戒按照圣神的教导如此分布于全年，使守斋的法则时刻呈现在我们面前。因此，我们在四旬期守春季斋戒，在圣神降临期守夏季斋戒，在第七个月守秋季斋戒，在现在的第十个月守冬季斋戒，深知没有什么事物不与天主的命令相关，一切受造物都为天主圣言服务，以教导我们，使我们从世界运行的枢轴本身，犹如从四部福音，学会不知疲倦地宣讲和实行。因为当先知说：\Quote{诸天述说天主的光荣，穹苍宣扬祂手的化工：日复一日发出言语，夜复夜传出知识，\BibleRef{咏 19:2-3}}还有什么真理不向我们说话呢？日日夜夜，祂的声音被听见，唯一天主手创之物的美永不停止将理性的教导注入我们心灵的耳朵，以致\Quote{天主不可见的事，可通过所造之物被感知和看见，\BibleRef{罗 1:20}}使人侍奉万有的创造者，而非受造之物。既然一切恶习都被克己摧毁，无论贪婪所渴求的、骄傲所追求的、情欲所贪恋的，都被这美德坚固的力量所战胜，谁不明白斋戒给予我们的帮助呢？因为斋戒中，我们不仅要在饮食上克制自己，也要在一切肉欲上克制。否则，忍受饥饿却不戒除不正当的愿望，节食折磨自己却不逃避罪恶的念头，便是徒劳的。那是肉体的斋戒，而非灵性的斋戒，其中身体被削弱，却坚持那些比一切快乐更有害的事物。灵魂在外表上装作主人，在内里却成为俘虏和奴隶，对肢体发号施令却失去自身自由的权利，这有什么益处呢？那灵魂通常（并且理应）在其仆人身上遇到反叛，因为它没有向天主尽应尽的服务。因此，当身体在食物上斋戒时，让心灵从恶习中斋戒，并按照其君王的法律审判一切尘世的挂虑和欲望。\par

% structure: p0007 source=h2 target=subsection reason=relative-heading-rank; base=h2

\subsection{三、因此，在心灵和身体上都要斋戒，并慷慨施舍，我们必将赢得天主最高的恩宠}

让我们记住，我们首先应爱天主，其次爱近人，所有情感必须如此调整，以免使我们远离对天主的敬拜或对同仆的益处。但我们如何敬拜天主，除非祂所喜悦的也为我们所喜悦？因为，如果我们的意愿就是祂的意愿，我们的软弱将从祂那里获得力量，意愿本身也是从祂而来；\Quote{因为是天主，}正如宗徒所说，\Quote{在你们内工作，使你们愿意，并实行祂的善意。\BibleRef{斐 2:13}}因此，人若将天主所赐的恩赐用于祂的光荣，并克制自己的倾向远离那些他知道会损害自身的事物，就不会因骄傲而膨胀，也不会因绝望而沮丧。因为藉着戒除恶毒的嫉妒、奢侈放荡的生活、愤怒的激动、报复的欲望，他因真正的斋戒而变得纯洁圣善，并以不朽喜乐的满足为食物；这样，他就学会如何以属灵的方式使用现世的财富，将其转变为天上的宝藏，不是为自己囤积所领受的，而是因施予而获得百倍的回报。因此，我们以慈父般的情怀劝勉你们，亲爱的诸位，通过慷慨施舍使这冬季的斋戒为你们结出果实，因主藉着你们喂养和衣著祂的穷人而喜乐；祂本可将赐予你们的财富赐给他们，但祂出于不可言喻的仁慈，愿意使他们因忍耐劳苦而成义，使你们因爱德工作而成义。因此，让我们在星期三和星期五守斋，在星期六与最可祝福的宗徒伯多禄一同守夜，他必以自己的祈祷助佑我们的祈求、斋戒和施舍，我们的主耶稣基督将呈献这些，祂与圣父及圣神永生永王，万世无疆。阿们。\par

\SourceRecord{Translated by Charles Lett Feltoe.；Nicene and Post-Nicene Fathers, Second Series；Vol. 12.；Edited by Philip Schaff and Henry Wace.；Buffalo, NY: Christian Literature Publishing Co.,；1895.；<http://www.newadvent.org/fathers/360319.htm>.}

% structure: p0001 source=h1 target=section reason=page-title

\section{讲道集 21}

\Emph{圣诞节讲道（一）}\par

% structure: p0003 source=h2 target=subsection reason=relative-heading-rank; base=h2

\subsection{一、人人共享圣诞的喜乐}

我们的救主，亲爱的诸位，今天诞生了：让我们欢欣吧。因为当我们庆祝那消灭死亡恐惧、带给我们预许永生喜乐的生命的诞辰时，没有悲伤的余地。没有人被排除在分享这喜乐之外。所有人都有一个共同的喜乐尺度，因为我们的主——罪恶与死亡的毁灭者——既然发现没有人无罪，便来释放我们所有人。让圣人因接近胜利而欢跃；让罪人因被邀请获得宽恕而欢喜；让外邦人因被召唤获得生命而鼓起勇气。因为天主子在神性计划不可测度的深渊所决定的时期满全时，取了人性，为使人性与它的创造者和好，以便死亡的创造者——魔鬼——能通过他（魔鬼）曾征服的那种本性而被战胜。在这为我们承担的争战中，战斗是以伟大而奇妙的公平原则进行的：全能的主不是以祂自己的威严，而是以我们的谦卑与祂凶残的仇敌交锋，用同样的形象和同样的本性来对抗祂，这本性确实分享我们的有死性，却完全无罪。这诞生确实异于我们所读到的所有其他诞生：\BibleQuote{没有一人是洁净的，连在地上只活了一天的婴儿也不例外。}\BibleRef{约 19:4}因此，情欲的冲动丝毫未进入那无与伦比的诞生，罪恶的律法丝毫未侵入。从达味的后裔中拣选了一位王室童贞女，使她受孕神圣的种子，并首先在心灵中、然后在身体上怀上神人结合的后裔。为免她因不知晓天上的计划而畏惧如此奇异的结果，她通过与天使的谈话得知，在她身上成就的事将来自圣神。她并不认为即将成为天主之母有损尊严。因为既然至高者的能力已应许成就这事，她为何要对这新奇受孕感到绝望呢？她的笃信也因一个前驱奇迹的见证而得到确认，且依撒伯尔获得了意想不到的生育能力：为使人们毫不怀疑，那位使不生育者怀孕的，也必能使童贞女怀孕。\par

% structure: p0005 source=h2 target=subsection reason=relative-heading-rank; base=h2

\subsection{二、降生成人的奥秘是天使和人类同乐的适宜主题}

因此，天主的圣言，祂本身就是天主，天主子——\BibleQuote{在起初就与天主同在}，藉着祂\BibleQuote{万有是藉着祂而造成的}，没有祂\BibleQuote{什么也没有造成}\BibleRef{若 1:1}——为将人类从永死中拯救出来，成了人：如此屈尊俯就，取了我们的谦卑，却不减少祂自己的威严，以致祂仍是祂所是的，并取了祂所不是的，使祂能将真正的仆人形象与祂与天主圣父相等的形象结合起来，并以如此紧密的联结将两种本性合一，以至较低的本性不因被高举而消失，较高的本性也不因新的结合而受损。因此，在不损害两者各自特性的情况下——这两者当时在一个位格中结合——威严取了谦卑，能力取了软弱，永恒取了有死性。为偿还我们状况所负的债务，不可侵犯的本性与可受苦的本性结合，真天主和真人结合成为一主，使得，按我们案件的需要，那同一位天主与人之间的中保、基督耶稣、真人，既能与一方同死，又能与另一方同复活。\par

因此，我们的救恩的诞生并未对童贞女的纯洁造成任何玷污，这是理所当然的，因为真理的诞生就是尊荣的持守。亲爱的诸位，这样的诞生正是天主的德能和天主的智慧——亦即基督——所当有的，藉此祂在人性上与我们合一，在天主性上超越我们。因为除非祂是真天主，祂就不能带给我们救治；除非祂是真真人，祂就不能给我们表率。所以，主诞生时，欢跃的天使歌唱道：\BibleQuote{天主受享光荣于高天}，他们的信息是：\BibleQuote{主爱的人在世享平安}\BibleRef{路 2:14}。因为他们看见天上的耶路撒冷正从世界万民中建立起来：对于这神性之爱难以言表的工程，卑微的人类应当如何欢欣，既然崇高天使的喜乐如此巨大？\par

% structure: p0008 source=h2 target=subsection reason=relative-heading-rank; base=h2

\subsection{三、基督徒因此必须活出与他们的头基督相称的生活}

因此，亲爱的诸位，让我们藉着祂的圣子、在圣神内感谢天主圣父，祂\BibleQuote{因祂爱我们的那大爱}，怜悯了我们，\BibleQuote{当我们死于过犯时，使我们与基督一同复活}\BibleRef{弗 2:4}，为使我们能在祂内成为新的受造物和新的产物。让我们脱去旧人及其行为；既已获得基督降生的分享，就让我们弃绝肉性的事业。基督徒，请认清你的尊严，并因成为神性本性的分享者，拒绝因堕落的行为回到旧日的卑贱。记住头和你是其肢体的身体。记念你曾被从黑暗的权势中救出，被带入天主的光明和国度。藉着圣洗圣事的奥迹，你成了圣神的宫殿：不要以卑劣的行为将这居者从你内赶走，并再次把自己置于魔鬼的奴役之下：因为你的赎价是基督的血，因为祂要在真理中审判你，祂曾以慈悲赎回了你，祂与圣父和圣神一起，永生永王。阿们。\par

\SourceRecord{Translated by Charles Lett Feltoe.；Nicene and Post-Nicene Fathers, Second Series；Vol. 12.；Edited by Philip Schaff and Henry Wace.；Buffalo, NY: Christian Literature Publishing Co.,；1895.；<http://www.newadvent.org/fathers/360321.htm>.}

% structure: p0001 source=h1 target=section reason=page-title

\section{第二十二篇讲道}

\Emph{论圣诞庆节，第二篇}\par

% structure: p0003 source=h2 target=subsection reason=relative-heading-rank; base=h2

\subsection{一、降生成人的奥迹要求我们喜乐}

亲爱的诸位，我们要在主内喜乐，以属神的喜乐欢欣，因为为我们已经出现了永远常新的救赎之日、远古预备之日、永恒真福之日。随着年岁的循环，我们救恩的纪念再次来临，这救恩从起初就预许，在时期一满时实现，并将永远长存。我们应当举心向上，朝拜这神圣奥迹，好使天主大恩的成效，能藉教会的大喜乐来庆祝。因为全能而仁慈的天主，祂的本性就是良善，祂的旨意就是权能，祂的作为就是慈悲：魔鬼的恶毒一经以仇恨的毒药杀害了我们，祂就在世界的元始预先宣布了祂的慈爱为我们必朽的人所预备的救药：向蛇宣告：女人的后裔要来，以祂的能力击碎那毒蛇昂起的头，这无疑是指示基督将降生成人，既是天主又是人，由童贞女所生，以祂无玷的诞生定人类仇敌的罪。这样，真天主在真实而完全的人性中诞生，圆满地拥有祂自己的，也圆满地拥有我们的。我们将造物主从起初在我们内所塑造的，以及祂所承担要修复的，称为\Quote{我们的}。因为欺骗者所引入、受骗者所接纳的，在救主身上毫无踪影。祂虽分担了人的软弱，但并未因此分享我们的过犯。祂取了奴仆的形体，毫无罪的玷污，增加了人性，并未减少天主性：因为那\Quote{空虚自己}，使不可见者成为可见，使万物的创造者与主宰甘愿成为可朽的，是怜悯的俯就，并非权能的缺失。\par

% structure: p0005 source=h2 target=subsection reason=relative-heading-rank; base=h2

\subsection{二、基督诞生的新特性解释}

因此，亲爱的诸位，当人类救赎所预定的时期来到时，天主子降入这世界的低下之处，从祂天上的宝座降下，却未离弃圣父的荣耀，以新的秩序、新的诞生而诞生。以新的秩序，因为祂在自己的本性内是不可见的，却在我们的本性内成为可见的；那不能受任何限制的，甘愿受限制；祂在万世之前就已存在，却在时间中开始存在；万有之主隐藏了祂不可测量的尊威，取了奴仆的形体；身为不能受苦的天主，祂不轻视成为能受苦的人，虽然祂是不朽的，却服从死亡的法律。祂以新的诞生而诞生，由童贞女怀孕，由童贞女生子，没有父性的欲望，没有损伤母亲的贞洁：因为这样的诞生没有沾染人性的污秽，正合乎那将要成为人类救主的一位，同时自身又拥有人性本质的自然。因为当天主降生成肉躯时，天主自己就是圣父，正如总领天使向童贞玛利亚作证说：\Quote{因为圣神要临于你，至高者的能力要庇荫你，因此，那要诞生的圣者，将称为天主子。\BibleRef{路 1:35}} 根源不同，但本性相似：不是藉男女交合，而是藉天主的能力得以成就；因为童贞女怀孕，童贞女生产，童贞女仍保持童贞。这里要考虑的不是生产者的状况，而是诞生者的意愿；因为祂按自己的意愿和能力降生为人。你若追究祂本性的真实，你必须承认是人的实质；你若探究祂诞生的方式，你必须承认是天主的权能。因为主耶稣基督来是为消除我们的污秽，而不是忍受它们；不是屈服于我们的过犯，而是医治它们。祂来是为治愈我们腐败的一切软弱和受玷污灵魂的所有创伤；因此，祂应以新的秩序诞生，将无玷纯洁的新恩赐带给人的身体。因为那诞生者的无玷本性必须保护母亲的原始童贞，而圣神灌注的能力必须保持祂为自己所选之圣所的无瑕和圣洁；就是说，圣神已决意要扶起跌倒者，恢复破碎者，并藉战胜肉体的诱惑，赏赐我们丰富贞洁能力，为使那在他人身上不能保留在生育中的童贞，能在他们重生时获得。\par

% structure: p0007 source=h2 target=subsection reason=relative-heading-rank; base=h2

\subsection{三、正义要求撒殚应由降生成人的天主击败}

亲爱的诸位，基督选择由童贞女诞生，这事实本身岂不正显示出最深远的安排吗？我是说，为使魔鬼不致察觉救恩是为人类而生，并因这属灵受孕的隐密，当他看见祂与常人无异时，便以为祂的诞生也与常人无异。因为他见祂的本性与众人相似，便以为祂与众人有同样的起源；他既未发现祂在必死性的软弱上异于常人，就不明白祂是无罪束缚的。虽然天主的真慈爱有无数方法可用来恢复人类，却选择了那一种计划，即运用正义的法则而非能力的效能来摧毁魔鬼的作为。因为古仇的骄傲并非不当地对所有人行使了其暴虐权，并以毫无道理的统治压迫那些自愿背离天主的命令、成为其意志奴隶的人。因此，他若不被他所征服的那位所征服，就不会公正地失去对人类亘古的奴役。为此，基督由童贞女因圣神受孕，并非藉男精，她不是因人的交合，而是因圣神而受孕。而在所有母亲中，受孕无不带罪污，这一位却从她受孕的根源获得了净化。因为未发生交合之处，便无罪的瑕疵渗入。她无瑕的童贞在供应物质时未知情欲。主从祂的母亲取了我们的本性，而非我们的过犯。奴仆的形状被造而无奴仆的身份，因为新人与旧人如此结合，既取了我们的真实人性，又除去了其古老的瑕疵。\par

% structure: p0009 source=h2 target=subsection reason=relative-heading-rank; base=h2

\subsection{四、降生成人欺骗了魔鬼，并使他打破了束缚人类的锁链}

因此，当仁慈全能的救主如此安排祂人性旅程的起始，以至将祂与人性不可分割的天主性之能力隐藏在我们软弱的帷幔之下时，狡猾的仇敌被出其不意，以为那为人类救恩而生的婴孩的诞生，与其他人的诞生一样受他控制。因为他见祂哭泣流泪，见祂裹在襁褓中，受割损，献上律法所要求的祭献。然后他察觉祂有通常的童年成长，毫不怀疑祂会按自然步骤达到成年。同时，他施加侮辱，倍增伤害，使用诅咒、凌辱、亵渎、辱骂，总之将他的全部怒气倾倒在他身上，用尽了各种试探：他深知自己如何毒害了人的本性，却未想到那一位既无分于原初的过犯——他已通过如此多证据证实了祂的必死性。那无耻的窃贼和贪婪的强盗持续攻击那本无所有的一位，并在执行原罪的普遍判决时，超出了他所依据的契约，向那一位要求罪恶的惩罚——祂身上并未发现任何过犯。于是，那致命契约的恶意条款被废除，因超额的不义，全部债务被勾销。那强人被自己的锁链束缚，邪恶者的一切计谋都回落到他自己头上。当世界之首被束缚时，他所囚禁的一切都被释放。我们的本性从旧染的污秽中被洁净，恢复其尊贵的地位，死亡被死亡消灭，诞生被诞生恢复：因为救赎一举消除了奴役，重生改变了我们的起源，信德使罪人成义。\par

% structure: p0011 source=h2 target=subsection reason=relative-heading-rank; base=h2

\subsection{五、劝勉基督徒分享降生成人的恩惠}

因此，无论你是谁，你虔诚忠信地以基督徒的名号自夸，就应正确估量这赎价。因为对你来说，你本是被遗弃者，从乐园之境被放逐，因疲乏的流徙而濒死，化为尘土灰烬，再无生活的希望，但由于圣言的降生成人，你获得了从远方回归造物主、认明你的根源、从奴役中获得自由、从弃儿升为儿子的能力：这样，你由可朽的肉身所生，却可因天主之神重生，藉恩宠获得你本性所无的，并且，你若藉着义子的精神承认自己是天主的儿子，就敢称天主为父。从良心的控告中得释放，你当向往天国，靠天主的助佑承行天主的旨意，在地上效法天使，以不朽养生的力量为食，无畏地为虔敬与敌对的诱惑争战，你若在天上战役中保持忠诚，就毫无疑问将在永生之王的得胜营中因你的胜利而获冠冕，因为那为信徒预备的复活将把你提升到分沾天国。\par

% structure: p0013 source=h2 target=subsection reason=relative-heading-rank; base=h2

\subsection{六、这节日与太阳崇拜毫无关联，如某些人所主张的}

因此，亲爱的诸位，既然怀着如此确信的望德，就应坚守你们所建基其上的信德：免得那同样的诱惑者——\Person{基督}已摧毁了它对你们的统治——再用它的任何诡计将你们夺回，并以它的欺骗伎俩玷污甚至当前节日的喜乐，用致死的观念误导较单纯的灵魂：有些人认为，我们这隆重的庆节所获得的尊荣，并非主要来自\Person{基督}的誕生，而是，按照他们的说法，来自新太阳的升起。这些人的心灵完全笼罩在黑暗之中，对真光毫无渐增的领悟；因为他们仍被异教的愚妄谬误所引，无法将心灵的眼目提升到他们肉眼所见之上，竟对服事世界的发光体施以神圣的尊崇。愿基督徒的灵魂绝不接纳任何这等邪恶的迷信和怪异的谎言。暂时的事物与永恒的事物、有形的事物与无形的事物、被治理的事物与治理者之间，有着无限的差距。因为尽管它们拥有奇妙的美丽，但并没有可受人崇拜的神性。因此，当受敬拜的是那从无中创造宇宙、凭其全能的方式按祂所意愿的形态与尺寸形成了天地实体的\Person{基督}的大能、智慧与威严。太阳、月亮和星辰对我们可能极为有用，极为悦目；但只有当我们为它们而感谢它们的创造者，并崇拜造它们的天主、而非那服事祂的受造物时，才是正当的。所以，亲爱的诸位，要在祂的一切作为与判断中赞美天主。要毫不怀疑地相信\Person{童贞玛利亚}纯洁的怀孕。要以神圣而真诚的侍奉，尊崇人类恢复的神圣而属神的奥迹。要怀抱降生在我们肉躯内的\Person{基督}，好使你们也能堪当看见祂作为荣耀的天主，以祂的威严为王；祂与圣父及圣神，在永远的天主性合一中共存，世世无穷。阿们。\par

\SourceRecord{Translated by Charles Lett Feltoe.；Nicene and Post-Nicene Fathers, Second Series；Vol. 12.；Edited by Philip Schaff and Henry Wace.；Buffalo, NY: Christian Literature Publishing Co.,；1895.；<http://www.newadvent.org/fathers/360322.htm>.}

% structure: p0001 source=h1 target=section reason=page-title

\section{讲道集 第二十三篇}

\Emph{论主诞辰节，第三篇}\par

% structure: p0003 source=h2 target=subsection reason=relative-heading-rank; base=h2

\subsection{一、降生成人的真理反复宣讲不会受损}

亲爱的诸位，今日隆重庆节之奥秘所关联的事，你们已熟知，也常听闻；但正如那可见之光使无疾之眼愉快，同样，救主的诞辰给健全之心带来永恒的喜乐；我们不可缄默不言，尽管我们不能按其所当论说。因为我们相信，依撒意亚所说的\BibleQuote{谁能述说他的世代？\BibleRef{依 53:8}}不仅适用于天主子与圣父同永恒的那奥秘，也适用于\BibleQuote{圣言成了血肉}的这次诞生。因此，天主、天主子，与圣父同等同体，出于圣父并与圣父同在，宇宙的创造者和主，祂完全临在于万有，又完全超越万有，按祂自己安排的恰当时间，为自己选择了这一天，要从荣福童贞玛利亚诞生，为拯救世界，而不损及母亲的尊荣。因为她的童贞在受孕和诞生时都没有受损：正如福音作者所说，\Quote{这一切的发生，是为应验主藉依撒意亚先知所说的话：\BibleQuote{看，童贞女将怀孕生子，人要称祂的名字为厄玛奴耳，即\InnerQuote{天主与我们同在}}。\BibleRef{玛 1:22}}因为圣洁童贞女这奇妙的生育，使她的后裔成为一位既是真人又是真天主的位格，因为每一个本性都保持其特性，以致二者之间不会有人位的分裂；受造物与造物主的结合也不是一位在内居住、另一位作为居所，而是两种本性彼此融合。虽然被取的本性是一种，取的本性是另一种，但这两个不同的本性如此紧密地结合，以至于同一位圣子既作为真人说\Quote{祂比父小}，又作为真天主说\Quote{祂与父同等}。\par

% structure: p0005 source=h2 target=subsection reason=relative-heading-rank; base=h2

\subsection{二、亚略派人士不能理解天主与人的结合}

结合，亲爱的诸位，造物主与受造物结合，亚略派的盲目不能用理智的眼看到，他们不信天主子与圣父有同样的光荣和同体，便说天主的性体是次等的，从那些应归于\Quote{奴仆的形体}的话中提出论据；就这形体而言，为了表明它不属于祂内的另一个或不同的位格，同一圣子以同一形体说：\BibleQuote{父比我大\BibleRef{若 14:28}}，正如祂以同一形体说：\BibleQuote{我与父原是一体}。在\Quote{奴仆的形体}中（祂在时日的末期为我们复兴取了这形体），祂小于父；但在\Quote{天主的形体}中（祂在万古之前即具有的形体），祂与父同等。在祂的人性卑微中，祂\BibleQuote{生于女人，生于法律之下\BibleRef{迦 4:4}}；在祂的天主性尊威中，祂常是圣言，\BibleQuote{万物是藉着祂而造的\BibleRef{若 1:3}}。因此，那位在天主形体中造了人的，在奴仆的形体中成了人。因为两种本性都保持其特性而不失：正如天主的形体没有取消奴仆的形体，同样奴仆的形体也没有损害天主的形体。所以，软弱与能力的奥秘，就同一人性而言，允许圣子被称为小于父；但是，在圣父、圣子、圣神三位一体中的天主性，排除一切不平等的观念。因为三位一体的永恒没有任何暂时的、在本性上不相似的事物；祂的意愿同一，实体同一，能力同等，却不是三位神，而是一位天主，因为那里有真正的、不可分的统一，不能有任何差异。因此，在真实人的整个完全的本性中，真天主诞生了，在祂自己的方面完全，在我们的方面也完全。所谓\Quote{我们的}，是指造物主从起初在我们内所造的，以及祂所承担修复的。因为欺骗者所引入的、被欺骗的人所犯的罪，在救主身上没有痕迹；祂既分担了人的软弱，并不因此分享我们的过犯。祂取了奴仆的形体而无罪的玷污，增进了人性而没有减损天主性：因为那\Quote{自我空虚}，使不可见者成为可见的，是怜悯的俯就，不是能力的衰竭。\par

% structure: p0007 source=h2 target=subsection reason=relative-heading-rank; base=h2

\subsection{三、降生成人为除去罪过是必要的}

因此，为使我们从原始的束缚和现世的错误中被召至永恒的福乐，祂亲自降临到我们这里，而我们却不能上升到祂那里；因为，许多人虽有爱真理之心，但变幻不定的意见受到误导恶魔的诡计所欺骗，人的无知被虚假所谓科学拖入各种对立的概念。但为消除这种对人心灵的嘲弄（使人被俘去服事骄傲的魔鬼），法律的教导并不足够，单凭先知的劝勉也不能恢复我们的本性；而必须在道德训诫之上加上救赎的实际，并使根本上败坏的我们的起源得以重新诞生。必须献上一只为我们的补赎而牺牲的祭品，这牺牲既是我们的族类，又没有我们的玷污，使天主这计划——祂乐意在耶稣基督的诞生和苦难中除去世界的罪——能够达于万代；而且我们不应为这些因时代性质而变化的奥迹感到困扰，反而应得坚固，因为使我们生活的信德从未改变。\par

% structure: p0009 source=h2 target=subsection reason=relative-heading-rank; base=h2

\subsection{四、降生成人的恩惠既回溯过去也指向将来}

因此，让那些以不忠的怨言反对天主安排、并唠叨说主的诞生太迟，好像这在世界末期中完成的事与过去的时代毫无关系的人停止他们的抱怨吧。因为圣言的降生成人只是促成了所成就的事：人类救恩的奥迹在远古时代从未停滞不前。宗徒们所预言的，正是先知们所宣告的：那一直为人们所相信的事，实现得并不算迟。相反，天主的智慧与良善推迟了这带来救恩的工作，使我们更易于接受祂的召唤：这样，经过如此多时代、由许多征兆、许多言语和许多奥迹所预言的事，在福音的今日就不致令人怀疑；救主的诞生，既超越一切奇迹和人类知识的尺度，就应在我们心中产生更为坚定的信德，因为它被预言的年代久远且一再重复。因此，天主为人所筹划的，并非新的计谋，也非迟来的怜悯：而是自世界的创立之初，祂便为众人预定了同一个救恩的缘由。因为天主的恩宠，即历代圣人之整体得以成义的恩宠，在基督诞生时是得以增强，而非开始：这伟大天主之爱的奥迹，如今充满普世，曾被如此有效地预先预示，以至于那些相信这应许的人所得的不亚于实际领受者。\par

% structure: p0011 source=h2 target=subsection reason=relative-heading-rank; base=h2

\subsection{V. 基督在我们血肉中的来临与成为祂身体的肢体相对应}

因此，亲爱的诸位，既然那将天主良善的一切财富倾注于我们的慈爱是如此显明——我们蒙召得永生，不仅得益于前人的有益榜样，也得益于真理本身可见有形的显现——我们就不应以懈怠或肉身的欢乐来庆祝主的诞生之日。若我们记得我们是哪一身体的肢体，我们与哪一首领相连，以免任何人如同不合榫的关节不与神圣建筑的其他部分相连，那么我们就必能相称且彻底地庆祝这一日。亲爱的诸位，请思量，并在圣神的光照下用心记念：那接纳了我们的是谁，而我们所接纳的又是谁；因为正如主耶稣因诞生而成为我们的血肉，我们亦因重生而成为祂的身体。因此，我们既是基督的肢体，也是圣神的宫殿。为此，蒙福的宗徒说：\Quote{你们要在你们身上光荣并背负天主：}因为主在提示我们祂自己的温良与谦卑的尺度时，也以祂救赎我们的德能充满我们，正如主亲自许诺说：\BibleQuote{凡劳苦和负重担的，你们都到我这里来，我要使你们安息。你们背起我的轭，跟我学罢，因为我是良善心谦的：这样你们必要找到你们灵魂的安息。} \BibleRef{玛 11:28} 那么，让我们背起统治我们的真理之轭——这不沉重也不辛苦——并效法祂的谦卑，因为我们愿意肖似祂的光荣。愿祂亲自帮助我们并引领我们到达祂的应许，祂能按祂的丰盛仁慈消除我们的罪恶，并在我们内完善祂的恩赐，我们的主耶稣基督，祂永生永王，直到永远。亚孟。\par

\SourceRecord{Translated by Charles Lett Feltoe.；Nicene and Post-Nicene Fathers, Second Series；Vol. 12.；Edited by Philip Schaff and Henry Wace.；Buffalo, NY: Christian Literature Publishing Co.,；1895.；<http://www.newadvent.org/fathers/360323.htm>.}

% structure: p0001 source=h1 target=section reason=page-title

\section{讲道 24}

\Emph{论圣诞节，第四篇。}\par

% structure: p0003 source=h2 target=subsection reason=relative-heading-rank; base=h2

\subsection{一、降生成人实现了它所有的预像和应许}

亲爱的诸位，天主的良善确曾以多种方式、在许多部分常常眷顾人类，并仁慈地将祂眷顾的许多恩赐赐予古时的世代；但在这些末后的时期，祂超越了祂惯常仁慈的一切丰盛，因为在基督内，仁慈亲自降临到罪人身上，真理亲自降临到迷途者身上，生命亲自降临到死者身上：为使那与圣父同等永恒、同等的圣言，能拿我们卑微的本性与祂的天主性结合，由天主而生为天主，也由人而生为人。这确是从创世之初就应许的，并一直通过许多事件和言语的预兆得到预言：但是，若基督的来临没有实现那些长久而隐秘的应许，若那在当时只惠及少数相信前景的人的事，如今没有在实现中惠及无数信友，那么这些预像和预表的神秘又能拯救多少人类呢？如今我们不再通过预兆和预像来相信，而是藉着福音的记载被确认，我们敬拜那我们所相信已经成就的事；先知的知识帮助我们认识，使我们毫不怀疑那藉着如此确切的预言所预告的事。因此，主对亚巴郎说：\Quote{地上万民都要因你的后裔蒙受祝福\BibleRef{创 22:18}}；因此，达味在预言的精神中歌唱说：\Quote{上主向达味起了誓，决不变更，我要使你的亲生儿子，登上你的宝座；}；因此，主又藉着依撒意亚说：\Quote{看，一位贞女将怀孕生子，人将称祂的名字为厄玛奴耳，意思是：天主与我们同在\BibleRef{依 7:14}}，又说：\Quote{由叶瑟的根茎将生出一根枝条，由他的根上将开出一朵花。}在这枝条中，无疑预言了蒙福的童贞玛利亚，她出自叶瑟和达味的族系，由圣神受孕，生出了人类肉身的一朵新花，成为一位童贞母亲。\par

% structure: p0005 source=h2 target=subsection reason=relative-heading-rank; base=h2

\subsection{二、降生成人是堕落唯一有效的补救}

愿义人在主内喜乐，愿信徒的心转向赞美天主，愿人子宣扬祂的奇事；因为在这件天主的工作中，我们卑微的境况特别认识到它的造物主多么看重它：在祂按自己的肖像造了我们，赐予人类伟大的恩赐之后，祂更丰盛地贡献了我们的恢复，当时主亲自取了\Quote{奴仆的形体}。因为虽然造物主为祂的受造物所付出的一切，都是同一父爱的一部分，但人升到天主的事物，不如天主下降到人性那么奇妙。但是，除非全能的天主屈尊这样做，没有任何义德、任何智慧的形式能将任何人从魔鬼的奴役和永死的深渊中拯救出来。因为那随着罪从一人传到众人的定罪仍将存留，而我们那被致命创伤腐蚀的本性，找不到任何补救，因为它无法凭自身的力量改变其状态。因为第一个人从地上得了肉身的实体，并因造物主的嘘气而获得理性精神，以便他按照造物主的肖像和模样生活，能手如明亮的镜子保持天主良善和公义的样式。如果他藉着遵守所赐的法律，坚持维护他本性这崇高的尊严，他那未腐化的心智甚至会将他属地的身体提升到天上的荣耀。但是因为他在不幸的鲁莽中，相信了那嫉妒的欺骗者，同意了他的狂妄建议，宁愿抢先取得那为他存留的尊荣的增加，而不是去赢得它，于是不仅那一个人，而且在他内所有后来的人都听到了判决：\Quote{你是土，你要归于土\BibleRef{创 3:19}}；因此，\Quote{那属土的怎样}，因此，\Quote{凡属土的也怎样\BibleRef{格前 15:48}}，没有人是不朽的，因为没有人是属天的。\par

% structure: p0007 source=h2 target=subsection reason=relative-heading-rank; base=h2

\subsection{三、藉着圣洗的重生，我们都成为基督诞生的分享者}

因此，为了解开这罪与死的锁链，全能的天主圣子，祂充满万有、包容万有，与圣父完全同等、同体，由祂而来、与祂同在，取了人的本性，万物的创造者和主屈尊成为有死的：祂为自己选择了一位祂所造的母亲，这位母亲未损贞洁的荣耀，提供了如此多的身体物质，以至于没有人类种子的玷污，这新人就拥有了纯洁和真理。因此，在由童贞圣母的子宫所生的基督内，人性与我们的并无不同，因为祂的诞生是奇妙的。因为祂是真天主，也是真人：在两个性体中都没有谎言。\Quote{圣言成了血肉}，这是藉着肉身的提升，而非天主性的亏损：祂如此调节了自己的大能和良善，以致通过取了我们的人性而提升了它，并没有因赋予而丧失自己的本性。在这次基督的诞辰中，按照达味的预言，\Quote{真理从地上生出，正义由天上俯视。}在这次诞辰中，依撒意亚的话也应验了：\Quote{愿大地产生并带来救恩，愿正义一齐萌生\BibleRef{依 45:8}}。因为人类肉身的土地，在第一个人身上受了诅咒，但在蒙福童贞女的这个后裔中，只产生了一个蒙祝福、免于其族类过犯的种子。每个人都在重生中分享这精神的起源；每个人在重生时，圣洗的水就像童贞女的子宫；因为那充满童贞女的同一圣神，也充满洗礼池，为使那神圣怀孕所推翻的罪恶，由此神秘的洗涤而被除去。\par

% structure: p0009 source=h2 target=subsection reason=relative-heading-rank; base=h2

\subsection{四、摩尼教徒因拒绝降生成人，陷入了可怕的邪恶}

在這一奧蹟中，親愛的諸位，摩尼教徒瘋狂的謬誤毫無分數，他們也無分於基督的再生，因為他們聲稱基督是由童貞瑪利亞肉體地誕生。如此，他們既不信祂真實的誕生，便也不接受祂真實的苦難；既不承認祂真實地被埋葬，就拒絕祂真正的復活。因為他們踏上了那可憎教義的危險道路，那裏全是黑暗與滑跌，他們便從謊言的懸崖墮入死亡的深淵，找不到任何可靠的立足之地；因為，除了他們其他一切魔鬼般的罪行外，就在基督敬拜的首要慶節上，正如他們最近的告白所顯示的，他們既在肉體上也在精神上縱情於污穢，喪失了自己的廉恥與信仰的純潔；以致他們在禮儀上與在教義上同樣褻瀆神明，同樣污穢不堪。\par

% structure: p0011 source=h2 target=subsection reason=relative-heading-rank; base=h2

\subsection{伍、其他異端含有部分真理，但摩尼教徒絲毫不含真理}

其他異端，親愛的諸位，雖然它們各具多樣性，都理應被定罪，然而它們各自在某些部分含有真理。\Person{亞略}主張聖子小於聖父、是一個受造物，並認為聖神如同其餘一切由同一聖父所造，如此他便陷入極大的褻瀆；但他沒有否認聖父的本體中永恆不變的天主性，儘管他無法在天主聖三的合一中看到它。\Person{馬其頓尼}在拒絕聖神的天主性時，缺乏真理之光，但他承認聖父與聖子具有同一能力與同一性體。\Person{撒伯里烏}因主張聖父、聖子、聖神中性體的合一不可分割而陷入無法解脫的錯誤，但他將本該歸於性體平等之事歸於性體的單一，並且因為他不能理解真實的天主聖三，他相信同一位格具有三重稱號。\Person{佛提諾}因心智盲目而誤入歧途，承認基督是出於我們性體的真實的人，但不相信祂在萬世之前是由天主所生的天主，因此失去了信仰的完整性，相信天主聖子取了一個真實的人性肉身，以至於聲稱其中沒有靈魂，因為天主性本身取代了靈魂。如此，如果公教信仰所詛咒的一切錯誤都被收回，那麼在一個個異端中可以找到某些可以與其該受責罰的背景分離的東西。但在摩尼教徒那可憎的教義中，絕對沒有任何可以被視為某種程度上可容忍的東西。\par

% structure: p0013 source=h2 target=subsection reason=relative-heading-rank; base=h2

\subsection{陸、基督徒必須持守同一信仰，不被引入歧途}

但是你們，親愛的諸位，我以不少於蒙福宗徒\Person{伯多祿}的熱切言辭勸告你們：\Quote{被揀選的種族，王家的司祭聖潔的邦國，天主自己的子民\BibleRef{伯前 2:9}，}建基於不可動搖的磐石、基督之上，並因主我們的救主真實地取了我們的肉性而與祂結合，你們要堅定於那信仰，就是你們在許多見證人面前所宣認的，並在其中藉著水與聖神重生，領受了救恩的傅油和永生的印記。但是\Quote{若有人傳給你們與你們所學過的不同，就當受詛咒\BibleRef{迦 1:9}。}拒絕將邪惡的傳說置於最清楚的真理之前，你們可能讀到或聽到任何違反公教與宗徒信經規律的事，當斷定其為完全致命與魔鬼的。不要被他們那虛假偽裝的守齋所迷惑，那些守齋趨向於毀滅人的靈魂，而非潔淨。他們確實披上虔誠與貞潔的外衣，但在這欺騙之下，他們隱藏自己行為的污穢，並從他們不敬虔的心靈幽深處射出箭矢，傷害單純的人；正如先知所說：\Quote{他們要在黑暗中射擊心裏正直的人。}堅固的堡壘是健全的信仰，真實的信仰，無須增減任何東西；因為除非它是唯一的，否則就不是信仰，正如宗徒所說：\Quote{只有一個主，一個信仰，一個洗禮，一個天主，眾人之父，祂超越萬有，貫徹萬有，且在我們眾人之內\BibleRef{弗 4:5}。}親愛的諸位，要以堅定不移的心靈緊抱這合一，並在其中\Quote{追求}一切\Quote{聖德\BibleRef{希 12:14}，}在其中實行主的誡命，因為\Quote{沒有信仰，是不可能中悅天主的，}並且沒有它，就沒有什麼是聖潔的，沒有什麼是純潔的，沒有什麼是活的：\Quote{因為義人因信德而生活，}誰若因魔鬼的欺騙而失去它，雖然活著，卻是死的，因為正如義德是因信德而獲得，同樣，永生也是藉著真實的信德而獲得，正如我們的主與救主所說：\Quote{永生就是：認識禰，唯一的真天主，和禰所派遣的耶穌基督。\BibleRef{若 17:3}}願祂使你們進步並堅持到底，祂與聖父及聖神，永生永王，萬世無疆。阿們。\par

\SourceRecord{Translated by Charles Lett Feltoe.；Nicene and Post-Nicene Fathers, Second Series；Vol. 12.；Edited by Philip Schaff and Henry Wace.；Buffalo, NY: Christian Literature Publishing Co.,；1895.；<http://www.newadvent.org/fathers/360324.htm>.}

% structure: p0001 source=h1 target=section reason=page-title

\section{讲道集 26}

\Emph{论圣诞庆节，第六篇}\par

% structure: p0003 source=h2 target=subsection reason=relative-heading-rank; base=h2

\subsection{一、圣诞清晨是默想诞辰最恰当的时刻}

在各日各时，亲爱的诸位，凡默想神圣事物者，常会想到我们的主与救主由童贞玛利亚所生，以激发心灵认识其创造主；无论心灵是沉浸于祈祷的叹息，或是赞美的欢呼，或是祭献的奉献，其属神的目光所最频繁、最信赖地注视的，莫过于天主圣子、同永圣父所生者，亦由人性诞生而生。但此应受天地朝拜的诞辰，没有比今日更向我们提示的了；今日，当晨曦仍然照耀自然界时，这奇妙奥秘的光辉便进入我们的感官。因为天使加俾额尔与惊讶的玛利亚的对话，以及她因圣神怀孕（这奇迹既被应许亦被相信），似乎不仅重现于记忆，而且重现于眼前。因为今日，世界的创造者由童贞玛利亚的胎中诞生，祂，那创造万有本性的，成了祂所造者的儿子。今日，天主的圣言披上了血肉，那从未被人眼所见者，也开始为我们手所触摸。今日，牧人们从天使的声音得知救主以我们血肉灵魂的实体诞生；今日，主羊群的领袖们预定了福音讯息的形式，使我们也能与天军一同说：\BibleQuote{天主受享光荣于高天，善意的人受享和平于地。}\par

% structure: p0005 source=h2 target=subsection reason=relative-heading-rank; base=h2

\subsection{二、基督徒本质上参与基督的诞辰}

因此，虽然那未受天主圣子尊威所轻视的婴孩，随着年岁增长而达至成熟，当祂苦难与复活的胜利完成后，所有为我们而行的谦卑行动都停止了，但今日的庆节为我们更新了由童贞玛利亚所生的耶稣的神圣童年；我们在朝拜救主诞辰时，发现我们正在庆祝我们自己生命的开端。因为基督的诞生是基督徒子民生命的泉源，头的诞辰即是身体的诞辰。虽然每个蒙召的人都有自己的次序，所有教会的儿女被时间间隔所分开，但整个信徒的身体在圣洗之泉中重生，与基督一同被钉于祂的苦难，与祂一同复活于祂的复活，与祂一同被置于父的右边于祂的升天，同样，他们也与祂一同生于这诞辰。因为任何信徒，无论在世界何处，若在基督内重生，便离开原初本性的旧路，经由重生进入新人；他不再被视为属地上父亲的族系，而是属于救主的后裔——救主成了人子，为使我们有能力成为天主的儿女。因为除非祂以这样的卑微降临到我们中间，没有人能凭自己的功德到达祂面前。愿属世的智慧不要在这点上蒙蔽蒙召者的心灵，愿属世思想的脆弱不要高抬自己反对天主恩宠的崇高，因为它很快会归回最卑微的尘土。在时代的终结，那从永远所预定的事得以完成；在现实面前，当预象和预表止息时，法律与先知成了真理。这样，亚巴郎被找到为万民之父，所应许的祝福藉他的后裔赐给了世界；不仅那些由血肉所生的以色列人，而且所有被接纳的义子的身体，都进入了为信德的儿女所预备的产业。不要被愚蠢争论的烦扰所动摇，不要让神圣之言的效果因人的计算而消散；我们与亚巴郎一同相信天主，\BibleQuote{没有因不信而动摇} \BibleRef{罗 4:20}，但\BibleQuote{最确切知道主所应许的，祂也能成就。}\par

% structure: p0007 source=h2 target=subsection reason=relative-heading-rank; base=h2

\subsection{三、与天主和平是祂赐予人的最佳恩赐}

因此，亲爱的诸位，救主不是由血肉种子而是由圣神所生，如此，第一次过犯的定罪没有触及祂。因此，所赐恩惠的伟大本身就要求我们以相称于其光辉的敬礼。因为，如蒙福的宗徒所教导的：\BibleQuote{我们领受了，不是世界的精神，而是那出于天主的精神，为使我们知道那些由天主赐给我们的事} \BibleRef{格前 2:12}；而这精神别无正确敬拜之法，除非我们将从祂所领受的献给祂。但在主恩惠的宝藏中，我们怎能找到比和平更适合此庆节之尊荣的呢？这和平在主的诞辰时首先由天使歌咏队所宣告。因为那和平产生天主的儿女，是爱德的乳母和合一的母亲；是有福者的安息和我们的永恒家乡；其本有工作与特殊职务乃是将那些从世界移开的人结合于天主。因此，宗徒劝勉我们趋向此善果，说：\BibleQuote{我们既因信德成义，就让我们与天主和平相处} \BibleRef{罗 5:1}。在这简短句子中，几乎总结了所有诫命；因为哪里有真和平，那里就不会缺少德行。但是，亲爱的诸位，与天主和平相处是什么意思？岂不是愿意祂所命令的，不愿意祂所禁止的吗？因为若人间的友谊寻求灵魂的平等和意愿的相似，而习惯的不同绝不能达到完全的和睦，那么，一个人若以天主不喜悦的事为乐，并渴望从他所知令天主不悦的事中得到快乐，他又怎能分享天主的和平呢？这不是天主儿女的精神；这样的智慧不为那义子的高贵家族所接纳。这蒙选的王族必须活出其重生的尊严，必须爱圣父所爱，在一切事上与其创造主无分歧，免得主再次说：\BibleQuote{我养育了儿女，提高了他们，他们却背叛了我；牛认识自己的主人，驴认识自己主人的槽，以色列却不知晓我，我的百姓也不认识我} \BibleRef{依 1:2}。\par

% structure: p0009 source=h2 target=subsection reason=relative-heading-rank; base=h2

\subsection{四、我们必须配得上作为天主的子女和朋友的召叫。}

亲爱的诸位，这份恩赐的奥义是伟大的，这份礼物超越一切礼物：天主竟称人为子，人竟称天主为父；因为这些称谓，我们领悟并认识到如此崇高的大爱。因为在自然的后裔和地上的家庭中，那些出身高贵的人因不良交往的过错而降低身份，不肖的子孙因祖辈的光辉而蒙羞；那么，那些因爱世界而不怕从基督的家族中被驱逐的人，将会落得怎样的下场？但如果父亲的光辉在后代身上重现能获得人的称赞，那么对于生于天主的人来说，重新获得造物主相似的光辉，并在自己身上彰显那生了他们的那一位，是多么光荣啊！正如主所说：\Quote{你们的光当在人前照耀，好使他们看见你们的善行，光荣你们在天之父\BibleRef{玛 5:16}？}我们确实知道，正如宗徒若望所说：\Quote{整个世界都卧在那恶者手下\BibleRef{若一 5:19}，}并且通过魔鬼及其天使的诡计，无数的企图要么以逆境恐吓努力向上的人，要么以顺境败坏他们；但是\Quote{在我们内的那一位比那反对我们的更大}，那些与天主和平、全心全意常对父说\Quote{愿你的旨意成就\BibleRef{玛 6:10}}的人，在任何战斗中都不会被击败，在任何攻击中都不会受伤。因为我们在告解中责备自己，并拒绝让心神顺从肉体的私欲，我们就激起了那罪恶创始者对我们的敌意；但是通过接受祂恩宠的侍奉，我们获得了任何事物都不能摧毁的与天主的和平，为使我们不仅因服从而向我们的君王投降，而且也因自由意志与祂结合。因为如果我们同心合意，如果我们愿意祂所愿意的，不赞同祂所不赞同的，祂将为我们完成所有的战争；那赐予意愿的那一位，也将赐予能力：使我们成为祂工作的同工，并带着信德的欢愉唱出那先知之歌：\Quote{主是我的光明，我的救恩：我还惧怕谁呢？主是我生命的保障：我还畏惧谁呢？}\par

% structure: p0011 source=h2 target=subsection reason=relative-heading-rank; base=h2

\subsection{五、基督的诞生是教会和平的诞生。}

那么，那些\Quote{不是由血统，也不是由肉情，也不是由人意，而是由天主而生\BibleRef{若 1:13}}的人，必须向圣父献上爱好和平的子女的同心同德；所有被收养的肢体都必须在新创造的首生者中相遇，祂来不是为实行自己的旨意，而是为实行那派遣祂者的旨意；因为圣父出于祂的恩宠，收纳为嗣子的不是那些不和谐的人，也不是那些不像祂的人，而是那些在感觉与情感上合一的人。那些按照同一模式被改造的人，必须有与模式相似的心神。主的诞辰是和平的诞辰：因为宗徒这样说：\Quote{祂是我们的和平，祂使双方合而为一；}既然我们或为犹太人或为外邦人，\Quote{藉着祂，我们在同一圣神内，得以进到父面前。}这正是祂在祂自愿预定的受难之日前，特别教导祂的门徒的话：\Quote{我把我的平安赐给你们，我把我的平安留给你们\BibleRef{若 14:27}；}并且为免祂的平安的特性因普遍的说法而不被注意，祂补充说：\Quote{不像世界所赐的那样\BibleRef{若 14:27}。}世界，祂说，有它的友谊，并且使许多分离的和好。也有在恶习上相等的心智，欲望的相似产生情感的等同。若有人偶然不喜低贱与不名誉的事，并将不合法的契约排除在他们的关系之外，这样的人们，即使他们是犹太人、异端者或异教徒，不属于天主的友谊，而属于这个世界的和平。但是属神的人和公教徒的和平，从天上降下并引领向上，拒绝与爱世俗的人相交，抵抗一切阻碍，从有害的快乐飞向真正的喜乐，正如主说：\Quote{因为你的财宝在哪里，你的心也必在哪里\BibleRef{玛 6:21}：}也就是说，如果你所爱的是下面的事，你将下降到最深处；如果你所爱的是上面的事，你将达到最高峰；愿和平的圣神引导我们并带我们到那里，我们的意愿和感觉是一致的，并在信德、望德和爱德上同心一意；因为\Quote{凡被天主圣神引导的，都是天主的子女\BibleRef{罗 8:14}}，祂与圣子及圣神一同为王，世世无穷。阿们。\par

\SourceRecord{Translated by Charles Lett Feltoe.；Nicene and Post-Nicene Fathers, Second Series；Vol. 12.；Edited by Philip Schaff and Henry Wace.；Buffalo, NY: Christian Literature Publishing Co.,；1895.；<http://www.newadvent.org/fathers/360326.htm>.}

% structure: p0001 source=h1 target=section reason=page-title

\section{讲道集 27}

\Emph{论圣诞庆节，第七篇。}\par

% structure: p0003 source=h2 target=subsection reason=relative-heading-rank; base=h2

\subsection{一、否认基督的天主性或人性同样危险}

真诚信友，对今日庆节怀有真诚与虔诚敬意之人，不会对主的降生成人抱有虚妄思想，也不会对祂的天主性怀有不敬之论。因为否认祂人性的真实性，或否认祂与圣父在光荣中的平等，同属危险的恶行。因此，当我们试图理解基督诞生的奥迹，即祂由童贞玛利亚诞生时，务要驱散一切世俗推理的阴云，并清除属世智慧之烟雾，使蒙光照的信德之眼得以清明：因我们所信靠的权威是神圣的，我们所遵循的教导是神圣的。不论我们在内心耳中聆听的是法律的证言、先知的神谕，还是福音的号角，那蒙圣神充满的蒙福若望以雷霆之声所宣告的，都是真实的：\BibleQuote{在起初已有圣言，圣言与天主同在，圣言就是天主。这圣言在起初就与天主同在。万物是藉着祂而造成的，凡受造的，没有一样不是藉着祂而造成的。}同样，同一宣讲者所说的话也是真实的：\BibleQuote{圣言成了血肉，寄居在我们中间；我们见了祂的光荣，正如父独生者的光荣。}因此，在两种性体内，同一天主子取了我们的本性，而不失自己的本性；按祂的人性更新了人，但在祂自己内却永恒不变。因为祂与圣父共有的天主性并未失去全能，\BibleQuote{奴隶的形态}也未凌驾于\OriginalTerm{天主的形态}{form of God}之上，因为那超越的永恒本体，为救赎人类而自谦自卑，将我们转入祂的光荣，但并未停止其为祂所是。因此，当天主独生子自称小于圣父，却又自称为与圣父平等时，祂在自己内显示出两种形态的真实性：不平等人性，平等则天主性。\par

% structure: p0005 source=h2 target=subsection reason=relative-heading-rank; base=h2

\subsection{二、降生成人改变了人类存在的一切可能性}

因此，天主子肉身的诞生并未减损或增添祂的尊威，因为祂不变的本体既不能减少也不能增加。因为\BibleQuote{圣言成了血肉}并非表示天主的本性变成了血肉，而是圣言将血肉纳入祂位格的统一中；在其中无疑接受了完整的人，这人在童贞玛利亚的胎中由圣神受孕，其童贞永久不失，天主子与祂如此不可分离地结合，以致那位无时间地从圣父本体而生者，自己在时间中由童贞玛利亚的胎中诞生。因为若非祂在我们的本性中卑微自下，而仍在自己内保持全能，我们便无法从永恒死亡的枷锁中得到释放。因此，我们的主耶稣基督，在诞生时是真人，却从未终止为真天主，在自己内开创了新的受造的开始，并在祂诞生的\Quote{方式}中，重新开始了人类属灵的生命，为消除我们按照肉体的出生所沾染的污秽，使那些被记载的人有可能无有罪孽种子地重生：\BibleQuote{他们不是由血气，也不是由肉欲，也不是由男欲，而是由天主\BibleRef{若 1:13}生的。}谁能领悟这奥迹？谁能表述这恩宠之举？有罪者回归无罪，旧人成为新人；陌生人获得义子名分，外人继承产业。不敬神者开始成为义人，贪婪者变为慷慨，不节制者变得贞洁，属地者成为属天者。这改变从何而来，除非来自至高者的右手？因为天主子来是要\BibleQuote{消灭魔鬼的作为\BibleRef{若一 3:8}}，并且祂如此将自己与我们结合，将我们与祂结合，以致天主下降到人的境地成为人被提升到天主那里。\par

% structure: p0007 source=h2 target=subsection reason=relative-heading-rank; base=h2

\subsection{三、魔鬼确切知道向每个人提供何种诱惑}

可是，真诚信友，面对天主向我们显示的无法言喻的仁慈，基督徒必须极其谨慎，以免再度陷入魔鬼的诡计，再次被已弃绝的错误纠缠。因为那古老的敌人不断\BibleQuote{伪装成光明的天使\BibleRef{格后 11:14}}，到处散布欺诈的罗网，竭尽全力败坏信徒的信德。他知道对谁施加贪婪的热衷，对谁以肚腹的诱惑进行攻击，在谁面前摆出放纵的快乐，在谁心中灌输嫉妒的毒液；他知道使谁被忧愁压倒，使谁被喜悦欺骗，使谁被恐惧惊吓，使谁被惊奇迷惑：没有一个人他不筛察其习性，不簸扬其忧虑，不窥探其情感；无论他见谁最专注于某种事务，便在那里寻找伤害他的机会。此外，他还有许多人被他捆绑得更紧，因他们适合他的计划，使他能利用他们的能力和口舌去欺骗别人。藉着他们，疾病的治愈、未来事件的预言、恶魔的安抚、幻象的驱除得到保证。此外还要加上那些虚假声称整个人类生命状况取决于星宿影响的人，以及那些把实际上属于天主或我们意志的事归之于不可改变的命运的人。然而，为了造成更大的损害，他们许诺说，如果那些不利的星宿受到恳求，命运可以改变。于是，他们邪恶的捏造自相摧毁：因为如果他们的预言不可靠，命运就无需惧怕；如果可靠，星宿便不应受敬拜。\par

% structure: p0009 source=h2 target=subsection reason=relative-heading-rank; base=h2

\subsection{四、某些人转向太阳并向它鞠躬的愚蠢做法应受谴责}

从这种教导体系也产生了一些愚昧之人的不敬行为，他们在黎明时分从高处崇拜初升的太阳：甚至有些基督徒认为这样做是恰当的，以至于在进入奉献给唯一永生的真天主的蒙福宗徒圣伯多禄的大殿之前，当他们登上通往高台的台阶时，他们转身向初升的太阳鞠躬，并弯下脖子向这灿烂的圆盘致敬。我们对发生这种情况感到悲痛和烦恼，这部分是由于无知的错误，部分是由于异教的精神：因为尽管他们中的一些人或许是在崇拜那美丽之光的造物主，而非作为受造物的光本身，但我们必须连这种仪式的表象也避免：因为如果有人放弃了神祇崇拜，在我们的崇拜中发现这种行为，难道他不会又回到他旧迷信的这一残余，视其为允许的，当他看到基督徒和异教徒都同样这样做时？\par

% structure: p0011 source=h2 target=subsection reason=relative-heading-rank; base=h2

\subsection{五、太阳和月亮为使用而创造，并非为崇拜}

因此，信友们必须放弃这种应受指责的做法，不可将对天主的独有尊荣与那些事奉受造物的人的仪式混为一谈。因为圣经说：\Quote{你要朝拜主，你的天主，惟独事奉祂。}\BibleRef{玛 4:10}。蒙福的约伯，按主所说，是\Quote{无可指责的人}，且\Quote{远离一切邪恶}\BibleRef{约 1:8}，他说：\Quote{我若见太阳发光，或明月运行，我的心里就偷偷欢喜，且用手亲吻；那岂不是莫大的罪恶，否认了至高的天主？}但太阳是什么，月亮又是什么，不过是可见受造界的元素和物质的光：一个更明亮，另一个稍弱。因为正如有白天和黑夜，造物主也设置了各种发光体，尽管在它们被造之前已有无太阳的白日和没有月亮的黑夜。但这些是为了服务于人的制造，让他这拥有理性的受造物能确知月份的区别、年岁的循环和季节的变化，因为通过不同时期的不等长度以及太阳升落变化的清晰指示，太阳结束一年，月亮开始月份。因为在第四天，正如我们所读，天主说：\Quote{让天上穹苍有光体，照耀大地，分开昼夜，作为标记和节令、日子和年岁，并让它们在穹苍中照耀大地。}\par

% structure: p0013 source=h2 target=subsection reason=relative-heading-rank; base=h2

\subsection{六、让我们醒来，善用我们所有的官能和机能}

醒来吧，人啊，认识你本性的尊贵。记念你是按天主的肖像受造的，这肖像虽在亚当内被败坏，却在基督内被重塑。善用可见的受造物，如你使用大地、海洋、天空、空气、泉源和河流：凡其中美丽奇妙之事，都归给造物主的赞美和光荣。不要被那使飞鸟、蛇、野兽、牲畜、苍蝇和蠕虫愉悦的光所奴役。将物质之光限定在你的肉体感官，并以你全部的心灵力量拥抱那\Quote{光照每一个人、来到这世界的真光，}\BibleRef{若 1:9}，先知也说到：\Quote{你们来归向祂，并受光照，你们的脸上就不会羞愧。}因为若我们\Quote{是天主的宫殿，天主的圣神住在}\BibleRef{格前 3:16}我们内，那么每位信友心中所拥有的，比他在天上所惊叹的更为伟大。因此，亲爱的诸位，我们并非命令或劝告你们轻视天主的工程，或认为良善的天主所造的美善之物与你们的信德有任何冲突，而是要你们合理且有节制地使用各类受造物和这世界的整个装备：因为正如宗徒所说：\Quote{所见的是暂时的；所不见的是永远的。}\BibleRef{格后 4:18}。因此，因为我们为今世而生，为来世而重生，不要让我们沉溺于暂时的美善，而要专注于永恒：为了能更近地注视我们的希望，让我们思考，在我们庆祝主诞生的奥迹时，天主的恩宠赐予了我们本性怎样的恩惠。让我们听宗徒的话：\Quote{因为你们已经死了，你们的生命与基督一同藏在天主内。当基督，你们的生命显现时，你们也要与祂一同出现在光荣中：}\BibleRef{哥 3:3}。祂与圣父及圣神永生永王，万世无疆。阿们。\par

\SourceRecord{Translated by Charles Lett Feltoe.；Nicene and Post-Nicene Fathers, Second Series；Vol. 12.；Edited by Philip Schaff and Henry Wace.；Buffalo, NY: Christian Literature Publishing Co.,；1895.；<http://www.newadvent.org/fathers/360327.htm>.}

% structure: p0001 source=h1 target=section reason=page-title

\section{讲道集 28}

\Emph{论圣诞节庆，第八篇。}\par

% structure: p0003 source=h2 target=subsection reason=relative-heading-rank; base=h2

\subsection{一、降生成人是喜乐的永不枯竭之源}

亲爱的诸位，虽然一切神圣的劝谕都催促我们\Quote{常在主内喜乐}\BibleRef{斐 4:4}，然而今天，当主诞生的奥迹清晰地照耀在我们身上时，我们无疑被激励要充满属灵的喜乐，以便我们能够归向那神圣仁慈的不可言喻的屈尊俯就，藉此，人的创造者屈尊成为人，而我们也在祂的本性中找到了我们所敬拜的祂。因为天主圣子，永恒而非受生的圣父的独生子，永恒地\Quote{具有天主的性体}，不变地、无时间地拥有与圣父毫无不同的属性，祂取了\Quote{奴仆的形像}，却没有丧失自己的尊威，为的是将我们提升到祂的状态，而不是将自己降低到我们的状态。因此，两种性体保持各自属性的同时，这种结合产生了合一的结果：凡属于天主性的，都与人性不可分离；凡属于人性的，都与天主性不可分割。\par

% structure: p0005 source=h2 target=subsection reason=relative-heading-rank; base=h2

\subsection{二、论童贞女的成孕}

亲爱的诸位，因此，在庆祝我们的主救主的诞辰时，让我们对蒙福童贞女的产子怀有纯洁的意念，相信在任何一个时刻，圣言的能力都不缺少她所成孕的肉身和灵魂，并且基督身体的圣殿并非预先获得形式与灵魂，好让它的居住者前来占据，而是藉着祂自己并在祂自己内，给予了新人一个开端，以致在唯一的天主子、天主而人者内，有天主性而无母，有人性而无父。因为她的童贞被圣神所孕育，同时毫无朽坏地生出了她族裔的后裔和创造者。因此，同一位主，正如福音所记载的，询问犹太人，根据圣经他们得知基督是谁的儿子，当他们回答说传统是祂要出自达味的后裔时，祂说：\Quote{那么，达味怎会受圣神感动称祂为主呢？他说：\InnerQuote{上主对我主说：你坐在我右边，直到我把你的仇敌放在你的脚下。}}犹太人无法解答所提出的问题，因为他们不明白，在同一基督内，达味的血统和天主性都曾被预言。\par

% structure: p0007 source=h2 target=subsection reason=relative-heading-rank; base=h2

\subsection{三、救赎人时，公义与仁慈必须兼顾}

但是，天主圣子的尊威——祂与圣父同等——在奴仆谦卑的形像中既不惧怕减少，也不需求增加；而祂为恢复人类所施行的仁慈，祂本可以单单凭藉天主性的能力来实现，以便将按天主肖像受造的受造物从残酷压迫者的轭下拯救出来。但因为魔鬼对最初之人的攻击并未表现得如此猛烈，以致在未得到他自由意志同意的情况下就将他拉到自己一边，所以必须这样摧毁人的自愿犯罪和敌意的欲望，以使公义的标准不致阻碍恩宠的赐予。因此，在整个全人类家族的普遍沦亡中，在天主计划隐秘处只有一个办法可以救助堕落者，那就是亚当的子孙中应当有一个生来就免于原罪、无罪无辜的人，藉祂的榜样和功德为其余的人赢得胜利。更进一步，因为自然生育不允许这样，并且不可能从我们有缺陷的种族中不藉种子而产生后裔，关于这事圣经说：\BibleQuote{谁能从污秽中生出洁净的呢？不是只有你吗？}\BibleRef{约 14:4}因此，达味的主成了达味的儿子，从所应许的枝条的果实中生出了一个无瑕疵者，两种性体结合成一位，以致藉着同一次成孕和诞生，我们的主耶稣基督出现了，祂内既有真天主性以行奇事，也有真人性以受苦难。\par

% structure: p0009 source=h2 target=subsection reason=relative-heading-rank; base=h2

\subsection{四、一切异端皆因不信基督的双重性体}

这么一来，亲爱的诸位，公教信仰可嗤笑那些狂吠反对它的异端者的错误，他们被世俗智慧的虚妄所欺骗，放弃了真理的福音，并且因为无法理解圣言的降生成人，便从光明之源为自己构筑了黑暗的缘由。因为在几乎调查了所有假信徒的意见——甚至那些胆敢否认圣神的人——之后，我们得出一个结论：几乎没有人误入歧途，除非他拒绝相信在基督内两个本性在一位格的宣认下的真实性。因为有些人只将人性归于主，另一些人只将天主性归于祂。有些人说，虽然祂内有真实的天主性，但祂的血肉是不真实的。另一些人承认祂取了真实的血肉，但声称祂并不具有天主圣父的本性；并通过将属于人的实体的东西归于祂的天主性，为自己造了一个较大的和一个较小的天主，然而在真实的天主性中是不可能存在等级的：因为凡小于天主的，就不是天主。另一些人认识到圣父与圣子之间没有区别，因为他们无法理解天主性的合一除非在位格的合一中，便主张圣父与圣子是同一的：以至于出生、被哺育、受苦、死亡、被埋葬、复活，都属于同一圣父，祂从头到尾承担了人和圣言两者的位格。有些人认为我们的主耶稣基督有一个不是出于我们实体的身体，而是取自更高更精微的元素；而另一些人则认为在基督的血肉中没有人的灵魂，而是圣言本身的天主性满足了灵魂的部分。但他们不智慧的主张进而成为这样一种形式：虽然他们承认主内有灵魂，却声称这灵魂没有理智，因为天主性本身对于在基督内作为人和作为天主的一切理性目的已经足够了。最后，同样一些人竟敢断言圣言的一部分转变为血肉，以至于在这一教条的多种多样中，不仅血肉和灵魂的本性，连圣言本身的本质也被解消了。\par

% structure: p0011 source=h2 target=subsection reason=relative-heading-rank; base=h2

\subsection{五、聂斯多略派与欧迪奇派在当前尤需避免}

还有许多其他令人惊愕的谬误，我们也不应使你们的耳朵厌倦，亲爱的诸位，去一一列举。但在所有这些各色各样的不虔敬——它们彼此紧密相连，因为一种亵渎形式与另一种之间存在关联——之后，我们恳请你们虔敬地留意避免这两种错误，尤其当前：一种以聂斯多略为始作俑者，曾一度试图得势，却未成功；另一种同样可咒骂，较近期以欧迪奇为其提倡导而兴起。前者竟敢主张蒙福的童贞玛利亚仅仅是基督人性的母亲，以便在她的受孕与生育中，不使人相信圣言与血肉曾发生过结合：因为天主子本身并未成为人子，而是出于祂单纯的屈尊，将自己与被造的人联结。这决不能被公教会的耳朵所容忍，这些耳朵已深深浸染在真理的福音中，以至于他们确知，除非祂本身就是那童贞女的儿子——而祂又是祂母亲的创造主——否则人类便无救恩的希望。另一方面，这亵渎性的更新近不虔敬的倡导者，承认在基督内两性的结合，却主张这一结合的结果是：两者中仅存其一，而另一的实体不再存在；这当然除非经由毁灭或分离，否则不可能终结。但这与纯正的信德如此对立，以至于若不失去基督信徒之名，便无法接受。因为如果圣言的降生成人是天主性与人性的结合，但正因其相遇，那原是二重的变成单一，那么从童贞女的子宫生出的就只是天主性，并且经历了接受营养和身体成长的虚假外表；略过人类状态的一切变化，被钉、死、埋葬的只是天主性：这样，按照那些这样思想的人，就没有理由希望复活，而基督也不是\BibleQuote{从死者中首生者}\BibleRef{哥 1:18}；因为祂不是一位应当被复活的那一位，如果祂不是一位能够被杀害的那一位。\par

% structure: p0013 source=h2 target=subsection reason=relative-heading-rank; base=h2

\subsection{六、天主性与人性自起初即在基督内同行}

亲爱的诸位，要远离你们的心中魔鬼启示的毒害谎言，并要知道圣子的永恒天主性与圣父同在时未曾成长，要明智地思考：对那同一本性，曾在亚当身上说过\BibleQuote{你是土，还要归于土}\BibleRef{创 3:19}，在基督身上却说\Quote{坐在我的右边}。根据那本性，基督与圣父平等，独生子从未低于圣父的崇高；祂与圣父同享的荣耀也不是暂时的拥有；因为祂就在圣父的右手边，关于这右手，在 \BibleBook{}{出谷纪} 中说\BibleQuote{上主，祢的右手在能力中受到了光荣}\BibleRef{出 16:6}；在 \BibleBook{}{依撒意亚} 中说\BibleQuote{上主，谁信了我们的报告？上主的手臂，向谁启示了呢？}\BibleRef{依 53:1}。因此，被接纳入天主子内的人，从祂在身体内的最初开端，就这样被纳入基督位格的合一中，以至于没有天主性，祂就不被孕怀；没有天主性，祂就不被带出；没有天主性，祂就不被哺育。在奇迹的行动中和在侮辱的承受中，是同一的位格；通过祂人性的软弱，被钉、死、埋葬；通过祂天主性的德能，第三天复活，升天，坐在圣父的右边，并且以人性的本性从圣父领受了那以天主性的本性祂自己同样所赐予的。\par

% structure: p0015 source=h2 target=subsection reason=relative-heading-rank; base=h2

\subsection{七、天主性的圆满通过头（基督）赐予身体（教会）}

亲爱的诸位，请以虔诚的心默想这些事，并常谨记宗徒的训诫，他劝告众人说：\Quote{你们要小心，免得有人用哲学和虚妄的欺骗，依照人的传统，而不是依照基督，把你们掳去；因为天主性的圆满完全以\OriginalTerm{肉身地}{bodily}的方式居住在祂内，你们也在祂内得以圆满\BibleRef{哥 2:8}。}祂不是说\Quote{属神地}，而是说\Quote{肉身地}，好使我们明白肉身的实体是真实的，那里有天主性的圆满在身体内的居所；当然，整个教会也被这圆满所充满，教会依附于元首，就是基督的身体；祂和圣父及圣神，永生永王，万世无疆。亚孟。\par

\SourceRecord{Translated by Charles Lett Feltoe.；Nicene and Post-Nicene Fathers, Second Series；Vol. 12.；Edited by Philip Schaff and Henry Wace.；Buffalo, NY: Christian Literature Publishing Co.,；1895.；<http://www.newadvent.org/fathers/360328.htm>.}

% structure: p0001 source=h1 target=section reason=page-title

\section{讲道词第三十一篇}

\Emph{主显节讲道（一）}\par

% structure: p0003 source=h2 target=subsection reason=relative-heading-rank; base=h2

\subsection{一、主显是圣诞的必要延续}

在不久前庆祝了无玷的童贞女为人类带来救主的日子之后，亲爱的诸位，可敬的主显节为我们延续了喜乐，免得我们在这些邻近而相似的奥迹中，使欢欣的热力和信德的热情冷却下来。因为关乎全人类的得救：天主与人之间的中保的婴孩，在他仍被拘留在小镇的时候，已经向全世界显示了出来。虽然他选择了以色列民族，并从该民族中拣选了一个家族，以摄取全人类的性体，但他不愿他诞生的早期日子隐藏在他母亲住处的狭小范围之内；却渴望很快被众人所认识，因为他屈尊为众人而生。因此，在东方地区，有一颗新光耀的星出现在三位贤士面前，这颗星比其他星辰更明亮、更美丽，能轻易吸引观看者的眼目和心灵，使他们立刻注意到这非同寻常的显现并非毫无意义。所以，那赐下记号的主，也赐给观看者对其的理解，并促使他们去探究他这样使其理解的事，而在探究之后，他把自己呈现给他们去寻获。\par

% structure: p0005 source=h2 target=subsection reason=relative-heading-rank; base=h2

\subsection{二、黑落德的阴谋徒劳无功。贤士的礼物具有自觉的象征意义}

这三位贤士跟随天上光芒的引导，以坚定的目光服从那引导他们的光辉的指示，被恩宠的光辉引向对真理的承认，因为他们以为君王的诞生是按人的意义宣告的，并且必须在一座王城中寻找。然而，那位取了奴仆的形体、来不是为审判而是为受审判的主，选择了白冷作为他的诞生地，耶路撒冷作为他的受难地。但黑落德听说一位犹太人的君王诞生了，怀疑是继承人，极为恐惧；为了消灭救恩之作者，他假意许诺归顺。他若效法贤士的信德，并把他图谋诡诈之事用于虔诚的目的，将是多么有福！愚昧嫉妒的盲目邪恶，你以为能靠你的疯狂推翻天主的计划吗？世界的统治者，赐予永恒国度的主，并不寻求暂时的。你为何试图改变那不可改变的被造秩序，并在别人的罪行中抢先？基督的死亡不属于你的时代。福音必须首先奠基，天主的国必须首先宣讲，首先给予病人治愈，首先行奇事。你为何希望自己承担属于他人工作之过犯？为何在无法实现你的邪恶阴谋时，却唯独使你自己背上愿望邪恶的罪责？你通过这种阴谋一无所获，也一事无成。那位自愿诞生的主，将自愿死去。因此，贤士们满足了他们的愿望，来到孩童主耶稣基督面前，同一颗星在他们前面引领。他们朝拜了在肉身中的圣言、在婴孩中的智慧、在软弱中的德能、在真实人性中的威严之主：并通过他们的礼物公开承认他们心中所信的，以显示他们信德和理解的奥迹。他们向天主献上乳香，向人献上没药，向君王献上黄金，自觉地尊崇在合一中的神性与人性：因为尽管每一性体有其本身属性，但二者的德能并无差异。\par

% structure: p0007 source=h2 target=subsection reason=relative-heading-rank; base=h2

\subsection{三、无辜者的屠杀与童贞女的受孕相一致，这再次教导我们纯洁的生活}

当贤士们返回本国，耶稣按天主的启示被带到埃及之后，黑落德的疯狂燃起徒劳的计谋。他下令杀死白冷所有的小孩子，由于不知道惧怕哪一个婴孩，他对所怀疑的年龄发布了一个普遍的判决。但是那邪恶的君王从世上除去的，基督却接纳他们进入天堂：对这些人，他尚未为他们倾流赎罪的血，就已赐予了殉道的尊严。因此，亲爱的诸位，请以忠信之心仰望永恒光明的慈悲光芒，并在朝拜为人类得救而施行的奥迹时，殷勤留意为你们所成就的事。热爱贞洁生活的纯洁，因为基督是童贞女之子。\Quote{应戒绝与灵魂交战的肉性私欲\BibleRef{伯前 2:11}，}正如我们所读到的，蒙福的宗徒在他的话语中劝勉我们说：\Quote{在恶意上应作小孩子\BibleRef{格前 14:20}，}因为荣耀的主使自己相似于必死之人的婴孩。追随天主子屈尊教导其门徒的谦卑。穿上忍耐的德能，藉此你们能赢得你们的灵魂；因为他是众人的救赎，也是众人的力量。\Quote{你们应追求天上的事，而非地上的事\BibleRef{哥 3:2}。}坚定地行走在真理和生命的道路上：不要被地上的事物阻碍，因为你们已预备了天上的事物，藉着我们的主耶稣基督，他与圣父及圣神，永生永王，直到永远。阿们。\par

\SourceRecord{Translated by Charles Lett Feltoe.；Nicene and Post-Nicene Fathers, Second Series；Vol. 12.；Edited by Philip Schaff and Henry Wace.；Buffalo, NY: Christian Literature Publishing Co.,；1895.；<http://www.newadvent.org/fathers/360331.htm>.}

% structure: p0001 source=h1 target=section reason=page-title

\section{第三十三篇讲道}

\Emph{论主显节，第三讲}\par

% structure: p0003 source=h2 target=subsection reason=relative-heading-rank; base=h2

\subsection{一、当我们还是罪人时，基督来拯救}

虽然我知道，亲爱的诸位，你们完全明白今日庆节的意义，福音的话也照例向你们阐明了，但为免我们有所遗漏，我仍将冒昧按主放在我口里的话来谈论这题目；这样，在我们的共同喜乐中，我们内心的热诚将因更明白庆节的缘由而更加真挚。天主的眷顾仁慈既决定在这末世拯救沉沦的世界，便在基督身上预定了万民的救恩；为使所有长期因邪恶错谬远离真天主崇拜的民族，甚至天主特选的子民以色列也几乎完全背离律法的规条，既都陷于罪恶之下，祂就怜悯众人。\par

因为正义处处衰微，全世界都交付于虚妄和邪恶，若非天主的大能延缓审判，全人类本应承受定罪。但震怒转为宽恕，为使所显示的恩宠更为显著，天主乐意在无人能夸耀自身功德的时候，运用赦罪的奥迹来消除人的罪恶。\par

% structure: p0006 source=h2 target=subsection reason=relative-heading-rank; base=h2

\subsection{二、来自东方的贤士是天主对亚巴郎应许的典型实现}

这不可言喻的仁慈的显示，亲爱的诸位，是在黑落德执掌犹太王权时发生的：当时，君王的正统继承已断绝，大司祭的权力已被推翻，一个外邦人获得了统治权；为使真君王的兴起获得预言的见证，这预言曾说：\Quote{权杖不离犹大，领袖不离他腰间，直到那应得者来到，万民都期望他。}关于这，曾经应许给最蒙福的始祖亚巴郎一个不可计数的子孙，不是由血肉种子，而是由丰饶的信德所生；因此它被比作众星之多，使他作为万民之父，可盼望一个属天的而非属地的后裔。因此，为产生这应许的后裔，以星宿为象征的继承人被一颗新星的出现所唤醒，好使诸天的侍奉服务于那以诸天为证之事。一颗比众星更亮的星唤醒了住在遥远东方的贤士，他们从这奇光的灿烂，并非不谙观测此类现象的人，领会了这标记的重要性：这无疑是由天主的神感在他们心中促成的，为使如此伟大景象的奥迹不向他们隐藏，且这在他们眼中不寻常的现象不晦暗于他们的心智。总之，他们审慎尽本分，准备了这样的礼物，以致在朝拜那一位的同时，也表明他们对祂三重职位的信仰：以黄金尊崇君王的位格，以没药尊崇人的位格，以乳香尊崇天主的位格。\par

% structure: p0008 source=h2 target=subsection reason=relative-heading-rank; base=h2

\subsection{三、特选的民族不再是犹太人，而是各邦国的信徒}

于是他们进了犹太王国的首府，在京城询问那他们得知生来为王的应被指给他们看。黑落德惊惶：他担心自身安全，为权力战栗，询问律法的司祭和经师，圣经关于基督诞生有何预言，他查明所预言的一切：真理启迪了贤士，不信蒙蔽了专家；属血肉的以色列不明白所读的，看不见所指示的；翻阅书卷，却不相信其言辞。犹太人，你的夸耀何在？你那出自亚巴郎血统的高贵出身何在？你的割损岂不成了未受割损\BibleRef{罗 2:25}？看，大的服事小的\BibleRef{创 25:23}，因你只在字面遵守那盟约的诵读，你成了外邦人的奴仆，他们进入了你祖业的份。愿外邦的丰盈进入先祖的家室，是的，愿它进入，并使应许之子在亚巴郎的后裔中获得祝福，而这祝福是血统的子孙所弃绝的。愿万民在三位贤士中崇拜宇宙的创造者；愿天主不仅只在犹太被人认识，而在全地，使祂的名在各地\Quote{在以色列中}为\Quote{大}。因为特选民族的后裔因不信而显为堕落，它的尊贵却因我们的信仰而成为众人共享。\par

% structure: p0010 source=h2 target=subsection reason=relative-heading-rank; base=h2

\subsection{四、无辜者的屠杀，以及基督随后逃往埃及，将真理带入埃及}

贤士朝拜了主，完成了他们所有的敬礼之后，按照梦中得到的警告，不按原路返回。因为他们既已相信基督，就不应再行走在旧生活的道路上，而应踏上新路，远离他们已弃绝的错误；这也是为了挫败黑落德的阴谋，他假借朝拜之名，正对婴孩耶稣策划恶毒计谋。因此，当他狡猾的希望破灭后，国王的怒火更盛。因为他推算贤士们指明的时间，便向白冷所有男婴大发暴怒，并对全城进行大屠杀，杀害了婴儿，使他们如此进入永恒的荣耀，以为若那里每一个婴孩都被杀，基督也必被杀。但祂——那位为救赎世界而推迟流血直到另一时刻的——由祂父母的帮助被带往埃及，如此寻得希伯来民族的古老摇篮，并藉更大上智安排的力量，行使真正若瑟的首牧职分；祂——那自天而降的\OriginalTerm{生命之粮}{Bread of Life}和\OriginalTerm{理性的食粮}{Food of reason}——除去了那比一切饥荒更甚的、埃及人心灵所苦的真理匮乏；若没有那次旅居，那\OriginalTerm{唯一牺牲}{one Victim}的象征便不完整；因为在那里，首先藉宰杀羔羊预表了\OriginalTerm{带来健康的十字架标记}{health-bringing sign of the Cross}和\OriginalTerm{主的逾越节}{Lord's Passover}。\par

% structure: p0012 source=h2 target=subsection reason=relative-heading-rank; base=h2

\subsection{V. 我们必须作为感恩的光明之子庆祝此庆节}

因此，亲爱的诸位，我们受这些神圣恩宠的奥迹教导，要以合理的喜乐庆祝我们初果的日子和万民蒙召的开始：\Quote{感谢}仁慈的天主\Quote{祂使我们堪当}，如宗徒所说，\BibleQuote{与圣徒在光明中分享基业，祂救了我们脱离黑暗的权势，把我们迁到祂爱子的国度里\BibleRef{哥 1:12}：}；因为正如依撒意亚先知预言：\BibleQuote{那坐在黑暗中的万民看见了大光，那些住在死影之地的人，光已照耀在他们身上。\BibleRef{依 9:2}}；关于他，他又对主说：\Quote{不认识你的民族，要呼求你；不认识你的百姓，要投奔你。}；这一天，\BibleQuote{亚巴郎看见了，喜乐了。\BibleRef{若 8:56}} 因为他明白，他信仰的子孙将因他的后裔——即基督——而得福，并预见藉信仰他将成为万民之父，\BibleQuote{将光荣归于天主，且坚信祂所应许的，祂也必能成就。\BibleRef{罗 4:21}}；这一天，达味在圣咏中歌颂说：\Quote{你所造的万民都要来朝拜你，上主，他们要光荣你的名。} 又说：\Quote{上主彰显了祂的救恩，将祂的正义在万民眼前公开显示。} 我们确知此事确实发生了，自那三位贤士在遥远之地被星星引导，认出并朝拜天地君王以来，[这星星对正确凝视的人不停每日出现。而且如果它能当基督隐藏于婴孩时使其被认识，那么当祂在尊威中为王时，它岂不更能启示祂？] 而且他们的朝拜确实敦促我们效法；好使我们尽可能事奉那邀请我们众人归向基督的恩慈天主。因为凡在教会内虔诚贞洁生活，并\BibleQuote{将心思放在上面的事，不放在地上的事\BibleRef{哥 3:2}}的人，在某种程度上像天上的光；当他本人保持圣洁生活的光明时，他就像一颗星为许多人指出通往主的道路。为此，亲爱的诸位，你们都应彼此帮助，好使你们在由正确信德和善工所达到的天主王国中，能像光明之子一样发光：藉着我们的主耶稣基督，祂与天主圣父及圣神，永生永王，于无穷世之世。阿们。\par

\SourceRecord{Translated by Charles Lett Feltoe.；Nicene and Post-Nicene Fathers, Second Series；Vol. 12.；Edited by Philip Schaff and Henry Wace.；Buffalo, NY: Christian Literature Publishing Co.,；1895.；<http://www.newadvent.org/fathers/360333.htm>.}

% structure: p0001 source=h1 target=section reason=page-title

\section{讲道 34}

\Emph{在主显节，第四篇}\par

% structure: p0003 source=h2 target=subsection reason=relative-heading-rank; base=h2

\subsection{I. 每年遵守主显节对基督徒有益}

亲爱的诸位，在那些见证天主仁慈之作为的日子，全心喜乐并隆重庆祝那些为我们的救恩而成就的事，是真虔诚的正确而合理的责任；因为时令循环的规律本身召唤我们如此虔诚地遵守，如今又将主显节之庆期带到了我们面前，这庆期是由主在圣子与圣父同永恒的、由童贞女所生之日不久后显现而祝圣的。而在此，天主眷顾为我们信德建立了一项伟大的保障，以便在每年重复对救主最早婴孩时期的敬礼时，祂内真实人性的产生能由最初的验证本身得以证明。因为正是这使不虔敬者成义，正是这使罪人成为圣人，即对我们主耶稣基督唯一的真天主性和真人性的信德：天主性，藉此祂在万世之前具有\Quote{天主的形体}而与圣父同等；人性，藉此祂在末世取了\Quote{奴仆的形体}而与人结合。因此，为了使这信德得以确证，以预先防备一切错误，天主的计划中有一项奇妙而慈爱的安排：一个居住在遥远东方、精通星象之术的民族，竟接受了那将要统治全以色列的婴孩诞生的征兆。因为一颗明亮新星的不同寻常的光辉出现在贤士们面前，当他们凝视其光芒时，心中充满了如此惊叹，以至于他们认为不能忽视如此明确地宣告给他们的事。而正如事件所显明的，天主的恩宠是这奇妙之事的原因：当整个白冷本身还不知道基督的诞生时，祂（恩宠）就将此事传递给了将要相信的万民，并藉着诸天的宣讲，宣告了人类言语尚不能解释的事。\par

% structure: p0005 source=h2 target=subsection reason=relative-heading-rank; base=h2

\subsection{II. 黑落德和贤士们原本都对那所指的国度抱有属世的观念；但后者学到了真理，前者却没有。}

但是，虽然使救主的诞生为万民所辨识是\OriginalTerm{神圣的谦卑}{Divine condescension}的职分，然而为了理解这奇妙的征兆，贤士们甚至可以从巴郎的古老预言中得到暗示，知道这事早已被预言，并因不断的重复而广为流传：\Quote{一颗星由雅各伯升起，一个人由以色列兴起，他将统治万民。}于是，这三人被天主藉着奇异星光的闪耀所激发，跟随那闪烁光芒的指引，以为他们会在耶路撒冷，即王城，找到那指定的婴孩。但发现自己的意见错误后，他们通过犹太人的经师和教师得知了圣经对基督诞生的预言；因此，受这双重见证的确认，他们以更加热切的信德寻求那同时被星光和确凿的预言话语所启示的那一位。当天主的谕旨通过大司祭的回答被宣告，而圣神的声音说：\BibleQuote{而你，白冷，犹大地啊，你在犹大的首领中绝不是最小的；因为将由你出来一位领袖，治理我的子民以色列\BibleRef{米 5:2}}，希伯来人中的首领相信他们所教导的，是多么容易、多么自然啊！但他们似乎与黑落德抱有物质的观念，将基督的国度视为与今世权势同等；因此他们盼望一个暂时的领袖，而他却害怕一个属世的对手。折磨你的恐惧，黑落德，是徒劳的；你试图将怒气发泄在你所怀疑的婴孩身上，是白费心机。你的国度不能容纳基督；世界之主并不满足于你统治的狭隘界限。你不想让祂在犹太为王的那一位，却无处不在为王；如果你自己服从祂的命令，你会统治得更幸福。你为什么不像你以诡诈的虚妄所承诺的那样真诚地去做呢？与贤士们同来，以恳求的朝拜敬拜真正的君王。但你，由于过于喜爱犹太人的盲目，不愿效法万民的信德，却将你顽梗的心引向残酷的诡计，尽管你注定既不能阻止你所害怕的那一位，也不能伤害你所杀害的人。\par

% structure: p0007 source=h2 target=subsection reason=relative-heading-rank; base=h2

\subsection{III. 贤士们的恒心带来了最重要的结果。}

于是，亲爱的诸位，贤士们听从星宿的指引进入白冷，正如福音作者告诉我们的：他们\Quote{极其高兴地欢喜起来，}\Quote{进了那家，看见婴孩和祂的母亲玛利亚，便俯伏朝拜了祂，打开自己的宝盒，给祂奉献了礼物，即黄金、乳香和没药。}\BibleRef{玛 2:10}这是多么奇妙的圆满认知的信德，不是由世俗智慧，而是由圣神的教训所教导的！这些离开本国、未曾见过耶稣、在祂的面容上未曾注意到任何足以促成如此系统化朝拜之事的人，何以在奉献礼物时遵循了这种方法呢？除非是：除了吸引他们肉眼的星宿外观之外，更灿烂的真理之光教导了他们的心灵，使他们在踏上艰辛旅程之前，必须明白那一位被预示了，祂应得黄金的君王尊荣、乳香的天主崇拜、没药的对有死性的承认。毫无疑问，就他们信德的光照而言，这样的信德和理解本身可能已足够，无需再用肉眼看那已用心灵完全注视过的事。但他们敏锐的勤勉，坚持到底直到找到婴孩，为将来的子民和我们的时代作出了贡献：如同主复活后宗徒多默触摸祂肉身伤口的痕迹，有益于我们众人一样，贤士们的来访为证实祂的婴孩时期也有益于我们。因此，贤士们看见并朝拜了犹大支派的婴孩，\Quote{按血统是达味的后裔}\BibleRef{罗 1:3}，\Quote{生于女人，生于法律之下}（\EditorialNote{迦拉达书第四章}），祂来\Quote{不是为废除，而是为成全}\BibleRef{玛 5:17}。他们看见并朝拜了那婴孩：身材幼小，无力帮助他人，不会说话，与普通人类孩童毫无区别。因为，正如那些证实祂不可见天主性尊威的见证是可信的，同样，也必须毫不怀疑\Quote{圣言成了血肉}，天主子永恒的本质取了人的真实本性：免得祂以后要行的不可言喻的奇事或祂要受的苦难，因其不一致而颠覆我们信德的奥迹：因为除了那些相信主耶稣既是真天主又是真人的人以外，无人能成义。\par

% structure: p0009 source=h2 target=subsection reason=relative-heading-rank; base=h2

\subsection{IV. \OriginalTerm{摩尼教异端}{Manichæan heresy}为否定真理而篡改\WorkTitle{圣经}}

亲爱的诸位，这无与伦比的信德、这万世宣扬的真理，受到\OriginalTerm{摩尼教}{Manichæans}人魔鬼般亵渎的反对：他们为杀害受欺骗者的灵魂，用邪恶和伪造的谎言编织了致命的邪恶教义网罗，在他们的疯狂见解的废墟上，人们陷入了如此深渊，以至于想象一个具有虚幻身体的基督，祂向人的眼睛和触摸呈现的没有一样是坚实、真实的，只显示了幻想之肉的空虚形状。因为他们认为，天主、天主子将自己置于一个女人体内，并将祂的尊威屈服于如此卑下，以致与我们的血肉本性结合，并以人类实体之真身体降生，这是不可信的——尽管这完全是祂权能的成果，而非祂受虐待，并且应相信这是祂光荣的屈尊，而非祂被玷污。因为若那可见的光不被包围它的任何污秽所玷污，太阳光线的光芒（无疑是有形受造物）也不被它渗入的任何肮脏或泥泞之地所污染，那么，有什么性质能够污染那永恒非受造之光的本质呢？因为祂与按自己肖像所造的受造物结合，为它提供了净化，自己却未受玷污，并治愈了它软弱的创伤，而未失去能力。由于这伟大而不可言喻的天主敬虔奥迹是由\WorkTitle{圣经}的一切见证所宣布的，我们所说的这些真理的反对者已经拒绝了通过梅瑟所赐的法律和先知们受天主默感的言论，并且篡改了福音和宗徒的书卷，删除或插入某些内容：在宗徒的名下和救主自己的话语下，伪造了许多虚假的卷册，用以加强他们虚谎的错误，并将致命毒药注入被欺骗者的心灵。因为他们看到一切都在反对和反驳他们，并且不仅\WorkTitle{新约}，而且\WorkTitle{旧约}都驳斥了他们亵渎和诡诈的愚蠢。然而，他们坚持疯狂的谎言，不断以欺骗扰乱天主的教会，劝说那些他们能诱捕的可怜人否认主耶稣基督真正取了人的本性；否认祂真正为世界的救恩被钉在十字架上；否认从祂被长矛刺透的肋旁流出了救赎的血和圣洗的水；否认祂被埋葬并于第三日复活了；否认在门徒眼前祂被举起升到诸天之上，坐在圣父的右边；并且，为了摧毁\WorkTitle{宗徒信经}的一切真理，使恶人无所畏惧、圣人无望可期，否认基督将审判生者死者；因此，那些被剥夺了这些伟大奥迹力量的人，可能学会敬拜太阳和月亮中的基督，并以圣神的名义敬拜\OriginalTerm{摩尼}{Manichæus}本人，即所有这些亵渎的发明者。\par

% structure: p0011 source=h2 target=subsection reason=relative-heading-rank; base=h2

\subsection{V. 避免与异端者一切来往，但要为他们向天主转求}

为坚固你们的心在信德与真理中，亲爱的诸位，愿今日的庆节帮助你们众人，并愿公教宣认因救主婴孩显现的见证而得到强化，同时我们谴责那些否认基督内有我们本性的肉身之人的亵渎言论。对此，蒙福的宗徒若望在毫不含糊的宣告中预先警告了我们，说：\BibleQuote{凡宣认耶稣基督是藉着肉身而来的神，便是出于天主；凡否认耶稣的神，便不是出于天主，这就是假基督。} 因此，不要让任何基督徒与这类人有任何关系，不要与他们结盟或来往。愿整个教会从以下事实中得益：因天主的仁慈，他们中有许多人已被识破，而他们自己的承认已揭示了他们的生活是如何亵渎。不要被他们关于食物与食物的区分、肮脏的衣物、苍白的面容所欺骗。禁食并非圣洁，如果它不是基于节制之原则，而是出于欺骗之意图。愿这成为他们危害无知之人、欺骗愚昧之人的终结；从此以后，任何人的跌倒都不可原谅；若某人今后被发现陷入可憎的谬误，他便不再被视为单纯，而是极其卑劣和乖僻。一项由教会认可且神圣建立的实践，我们非但不禁止，反而鼓励你们去做：你们甚至应为这样的人向主恳求；因为我们也含着泪和悲伤同情那些受骗灵魂的残骸，效法宗徒们的仁爱榜样，以便与软弱的人一同软弱，并\BibleQuote{与哭泣的人一同哭泣。} 因为我们希望，天主的仁慈能因许多眼泪和跌倒者的适当悔改而赢得；因为只要生命仍在身体内，就没有人的归正可被绝望，而我们都愿藉着主的帮助渴望所有人的革新，祂\BibleQuote{扶起被压制的人，释放被捆绑的人，开启瞎子；} 愿尊崇和光荣归于祂，直到永远。阿们。\par

\SourceRecord{Translated by Charles Lett Feltoe.；Nicene and Post-Nicene Fathers, Second Series；Vol. 12.；Edited by Philip Schaff and Henry Wace.；Buffalo, NY: Christian Literature Publishing Co.,；1895.；<http://www.newadvent.org/fathers/360334.htm>.}

% structure: p0001 source=h1 target=section reason=page-title

\section{讲道集 第三十六篇}

\Emph{论主显节，第六篇}\par

% structure: p0003 source=h2 target=subsection reason=relative-heading-rank; base=h2

\subsection{一、贤士的故事不仅是历史上过去的事实，而且对我们每日都有实用}

亲爱的诸位，基督救主初次显现于万民之日，我们应以恭敬的敬礼加以尊崇；今天，那三位贤士心中所怀的喜乐，也当存在于我们心中，那时他们因一颗新星的征兆和引领而惊醒，他们相信这是所应许的，便俯伏于天地君王之前。因为那一天并未如此过去，以致那时所启示的伟大工程也随之过去，而只有事情的传闻传给我们，供信德接纳、记忆庆祝；因为藉着天主屡次重复的恩赐，我们的时代日日享有了最初时代所拥有的果实。因此，虽然从福音读给我们的叙述，恰当地记录了那三天——那三位既未受先知预言教导，也未受法律见证教诲的人，从东方最远的地区前来承认天主——但我们在所有蒙召者的光照中，更清晰、更丰盛地看到同一件事现在正在延续，因为依撒意亚的预言得以实现，他说：\Quote{上主已在万民眼前显露祂的圣臂，地上的万民都看见了来自上主我们天主的救恩；}又说：\Quote{那些没有听过有关祂消息的人将要看见，那些没有听见的人将要领悟。}因此，当我们看见那些耽于世俗智慧、远离对耶稣基督信仰的人，从他们错误的深渊中被带出，并被呼召承认真光，这无疑是神圣恩宠的光辉在运作；凡照耀他们内心黑暗的新光，都来自同一颗星的光线：既令人惊奇，又带领被其光辉所照耀的心灵前去朝拜天主。但是，如果我们细心思考，想要看看那三种礼物如何也被所有以信德之脚前来归向基督的人所奉献，难道在真信徒的心中不也重复同样的奉献吗？因为承认基督为宇宙君王的人，从他心灵的宝库中献上黄金；相信天主的独生子将人的真实本性结合于自己的，献上没药；而宣认祂丝毫不亚于圣父的威严的，便以乳香的方式朝拜祂。\par

% structure: p0005 source=h2 target=subsection reason=relative-heading-rank; base=h2

\subsection{二、撒殚仍然继续黑落德的诡计，并好像扮演他来反对基督}

亲爱的诸位，仔细思考这些类比，我们也会发现黑落德的角色并不缺乏，魔鬼本身如今是不知疲倦的模仿者，正如那时他是秘密的挑唆者一样。因为万民的蒙召使他痛苦，他权力的每日败坏使他受折磨，他因处处被遗弃、真君王处处受到朝拜而悲伤。他预备计谋，策划阴谋，爆发凶杀，并且为了利用那些仍被他欺骗者的残余，他在犹太人中间因嫉妒而消耗自己，在异端者中间奸诈地埋伏，在异教徒中间燃烧出残酷。因为他看见永恒君王的权力是不可战胜的，祂的死已经消灭了死亡本身的权势；因此他武装了自己全部伤害的技巧，来攻击那些事奉真君王的人：用法律知识所产生的骄傲使一些人刚硬，用假信仰的谎言使另一些人堕落，并煽动其他人达到迫害的疯狂。然而这位\Quote{黑落德}的疯狂被那一位所战胜并归于虚无，祂甚至以殉道的荣耀加冕了婴儿，并赐予祂的信徒们如此不可战胜的爱，以致他们敢以宗徒的话说：\Quote{谁能使我们与基督的爱隔绝？是困苦吗？是窘迫吗？是迫害吗？是饥饿吗？是赤身露体吗？是危险吗？是刀剑吗？正如经上所载：为了你的缘故，我们终日被杀害，被视为待宰的羊。然而，借着那爱我们的主，我们在这一切事上都得胜有余}\BibleRef{罗 8:35}。\par

% structure: p0007 source=h2 target=subsection reason=relative-heading-rank; base=h2

\subsection{三、主动迫害的停止并不消除持续警醒的需要：撒殚只是改变了策略}

诸位所亲爱者，我们认为这样的勇气，并非仅在那时代，即世上的君王和今世的一切权势，以邪恶的暴行猛烈攻击天主子民，以为若能铲除基督徒之名于地上，便是他们的最大荣耀，却不知天主的教会因他们残酷的疯狂而增长，因为借着殉道者的折磨和死亡，那些被认为人数减少的，反而因榜样的力量而增加。总之，迫害者的攻击使我们的信德获得如此大的力量，以至于王室的尊荣也没有比世界的元首成为基督的肢体更大的装饰；他们的夸耀不在于生来身穿紫袍，而在于在圣洗中重生。然而，由于昔日风暴的压力已经平息，争斗停止后，一段时期的安宁似乎长久地向我们微笑，那些从和平统治本身产生的分歧需要谨慎提防。因为仇敌在公开迫害中证明无效之后，现在运用隐秘的技巧进行残忍的伤害，为要以享乐的绊脚石绊倒那些他不能用苦难打击的人。因此，看到君王的信德与他为敌，且不可分的天主圣三同一天主性在宫殿中如在圣堂中一样被虔诚敬拜，他因基督徒的血不被容许流出而悲伤，转而攻击他不能致死之人的生活模式。他将没收的恐惧变为贪婪的火焰，用贪欲腐蚀那些他用损失不能折其心志的人。因为那长期习惯已深深根植于他本性中的恶毒傲慢，并未放下仇恨，而是改变了性质，为要借着谄媚征服信友的心志。他用贪欲点燃那些他不能用折磨使之痛苦的人；他播下纷争，点燃情欲，鼓动舌簧，并为了使更谨慎的心不致退缩而远离他无法无天的诡计，他提供犯罪的机会；因为他一切计谋的唯一果实就是：那不以牛羊的祭献和焚香受敬拜的，当以各种邪恶行为向他致敬。\par

% structure: p0009 source=h2 target=subsection reason=relative-heading-rank; base=h2

\subsection{四、及时的悔改获得天主仁慈的眷顾}

因此，诸位所亲爱者，我们和平的状态有其危险，而那些不抵抗邪恶欲望的人，若以为他们信德特权的自由是安全的，那是徒然的。人的心由其行为的性质显明，心灵的风尚由其行动的本质暴露。因为有些人的情形，如宗徒所说：\Quote{他们自称认识天主，但以行为否认祂} \BibleRef{铎 1:16}。的确，当声音所听到的善不在良心中时，否认的罪过就真正招来了。人性软弱容易滑入过犯；因为没有罪是没有吸引力的，欺骗性的享乐很快就让人顺从。但我们应从肉身的欲望奔向属灵的救助；那认识自己天主的心灵应当远离仇敌的邪恶暗示。利用天主的容忍，不要固守你的罪，因为惩罚被推迟了。罪人不应对自己的免罚感到安全，因为如果他失去了悔改的时间，他将找不到仁慈的地方，正如先知所说：\Quote{在死亡中无人记念你；在阴府里谁会承认你？}又如那位体验到自己改过和复原困难的人，当投靠助人天主的仁慈，求祂——\Quote{那位扶起跌倒者，并举起所有被压碎者}——折断邪恶习惯的锁链。忏悔者的祈祷不会落空，因为仁慈的天主\Quote{必赐给敬畏祂者的心愿}，且要赐下所求的，正如祂赐下求告的根源。借着我们的主耶稣基督，祂和圣父及圣神，永生永王，万世无疆。阿们。\par

\SourceRecord{Translated by Charles Lett Feltoe.；Nicene and Post-Nicene Fathers, Second Series；Vol. 12.；Edited by Philip Schaff and Henry Wace.；Buffalo, NY: Christian Literature Publishing Co.,；1895.；<http://www.newadvent.org/fathers/360336.htm>.}

% structure: p0001 source=h1 target=section reason=page-title

\section{讲道集 39}

\Emph{论四旬期，一}\par

% structure: p0003 source=h2 target=subsection reason=relative-heading-rank; base=h2

\subsection{一、由希伯来人的榜样所见节制的益处}

昔日，当希伯来人民和以色列各支派因自己可耻的罪恶而被培肋舍特人的残酷暴政压迫时，为了能够战胜他们的敌人，正如神圣历史所记载，他们通过命令斋戒，恢复了身心的力量。因为他们明白，他们之所以招致这种艰苦而悲惨的奴役，是由于他们疏忽了天主的命令和邪恶的行为，并且除非他们首先抵制了罪恶，否则用武器争斗是徒劳的。因此，他们禁绝饮食，对自己施以严格自律的管教，并且为了征服敌人，首先战胜了自己内里对美味的诱惑。结果，他们斋戒时，那些凶猛的敌人和残忍的压迫者屈服于他们，而这些敌人当初在他们饱足时曾奴役他们。同样，我们这些身处诸多敌对和冲突之中的人，亲爱的诸位，若能善用同样的方法，便可通过一点谨慎而得到医治。因为我们的情况几乎与他们相同，正如他们受到肉体敌人的攻击，我们则主要受到属灵敌人的攻击。如果我们能藉着天主的恩宠改正自己的行为而战胜他们，我们肉体敌人的力量也将在我们面前削弱，并且通过我们自身的修正，我们将使那些对我们构成威胁的人变得软弱，这些人之所以可怕，不是因他们自己的功绩，而是因我们的缺失。\par

% structure: p0005 source=h2 target=subsection reason=relative-heading-rank; base=h2

\subsection{二、利用四旬期战胜敌人，并如此预备迎接复活期}

因此，亲爱的诸位，为了能够战胜我们所有的敌人，让我们通过遵守天主的命令寻求天主的助佑，因为我们知道，除非我们战胜自己，否则无法战胜我们的对手。因为我们与自己多有交锋：肉性渴望的与灵性相反，灵性渴望的与肉性相反。在这种分歧中，如果肉体的渴望更强，精神将可耻地失去其应有的尊严，而本该统治的却去服侍，这将是最灾难性的。但是，如果精神服从它的统治者，并喜悦于来自高天的恩赐，就会践踏世俗享乐的诱惑，不允许罪恶在必死的身体中为王，那么理性将维持有序的至高地位，属灵邪恶的计谋也不能摧毁它的堡垒：因为只有当肉体受精神的判断所统治，精神受天主的旨意所引导时，人才享有真正的平安和真正的自由。亲爱的诸位，虽然这种备战状态应当始终保持，以便我们警觉的敌人能被不懈的努力所战胜，但现在，当我们狡猾敌人的计谋比以往更加诡诈时，我们必须更加急切地寻求它，更加热心地培养它。因为他们知道，最神圣的四旬期日子现已临近，在遵守这些日子时，所有往昔的懈怠都受到惩罚，所有疏忽都得到警惕，他们便将所有的仇恨力量集中于这一点：即那些打算庆祝主的神圣逾越节的人，可能在某件事上被发现不洁，并在本应获得赎罪的地方产生绊脚石。\par

% structure: p0007 source=h2 target=subsection reason=relative-heading-rank; base=h2

\subsection{三、战斗是考验我们信德所必需的}

因此，亲爱的诸位，当我们临近四旬期的开始——这是更加谨慎事奉主的时候，因为我们仿佛进入了一场善功的竞赛——让我们预备好我们的灵魂与诱惑搏斗，并且明白，我们越热心于自己的救恩，对手的攻击就越坚定。但\Quote{但那在我们内的比那在世界上的更大\BibleRef{若一4:4}，} 因着祂，我们倚靠祂的力量而强大。因为主允许自己受试探者的试探，正是为了这个目的：使我们既因祂的榜样受教，又因祂的助佑得到坚固。因为祂战胜了那仇敌，正如你们所听过的，不是凭实力，而是藉引法律上的话，由此祂赋予人更大的尊荣，并对那仇敌施以更重的惩罚，因祂现在不是以天主，而是以人的身份征服了人类的仇敌。因此，祂战斗，是为使我们此后也能战斗；祂得胜，是为使我们也能同样得胜。因为，亲爱的诸位，没有诱惑的考验就没有有力的作为，没有考验就没有信德，没有仇敌就没有战斗，没有冲突就没有胜利。我们的生命正处于网罗中，处于战斗中；如果我们不愿受欺骗，就必须警醒；如果我们想要得胜，就必须战斗。因此，最有智慧的撒罗满说：\Quote{我儿，你去服事天主时，应预备你的灵魂受试探\BibleRef{德2:1}。} 因为他是一个满有天主智慧的人，知道追求宗教包含艰苦的斗争，也预见到战斗的危险，便事先警告那将要战斗的人；免得试探者乘其不备而来，发现他未曾准备，意外地伤害了他。\par

% structure: p0009 source=h2 target=subsection reason=relative-heading-rank; base=h2

\subsection{四、基督徒的甲胄兼具防御与攻击}

所以，亲爱的诸位，我们既受了神圣学识的教导，就应有意识地参与当前的竞赛与争斗，听宗徒所说：\BibleQuote{我们的战斗不是对抗血和肉，而是对抗率领者，对抗掌权者，对抗这黑暗世界的霸主，对抗在天界中的邪恶势力} \BibleRef{弗 6:12}，并且不要忘记，这些仇敌感觉到，我们为得救所作的一切努力都是针对他们的，正是由于我们寻求某种善，我们就在挑战仇敌。因为这是我们与他们之间一场古老的纷争，由魔鬼的恶意所培育，以致他们因我们成义而受苦，因为他们已从那些我们靠着天主的助佑正在前进的善事上堕落了。因此，如果我们被高举，他们就被打倒；如果我们被加强，他们就被削弱。我们的痊愈就是他们的打击，因为我们的伤口的治愈使他们受伤。\Quote{所以，要站稳，}亲爱的诸位，如宗徒所说，\BibleQuote{用真理作带，束起你们腰间的思虑，用和平福音的预备作鞋穿在脚上；在一切事上拿起信德的盾牌，使你们能熄灭恶者的一切火箭；并带上救恩的头盔，拿着圣神的剑，就是天主的圣言} \BibleRef{弗 6:14}。看，亲爱的诸位，我们用多么强大的武器，多么坚不可摧的防御被我们的统帅所武装，祂因多次胜利而闻名，是基督战斗的无敌大师。祂用贞洁的带子束了我们的腰，祂用和平的绳索穿了我们的脚：因为不束腰的士兵很快就被不洁的诱惑者击败，不穿鞋的容易被蛇咬。祂给了信德的盾牌保护我们全身；在我们头上放了救恩的头盔；祂用剑装备了我们的右手，即真理之言：使属灵的战士不仅安全无伤，而且有能力伤害攻击者。\par

% structure: p0011 source=h2 target=subsection reason=relative-heading-rank; base=h2

\subsection{五、四旬期所要求的禁食不仅是食物的克制，也是其他邪恶欲望的克制，尤其是愤怒。}

因此，亲爱的诸位，依靠这些武器，让我们积极主动且无所畏惧地进入摆在我们面前的竞赛：这样，在这四旬期的努力中，我们就不会满足于仅仅以此为目的，即认为只须禁戒食物是可取的。因为我们肉身的减损还不够，如果灵魂的力量没有增强的话。当外在的人稍被制服，就让内在的人稍得更新；当肉体的放纵被禁止时，让我们的思想因属灵的喜悦而振奋。让每个基督徒自我反省，严肃地探察内心；确保那里没有不和谐，没有不义的欲望被隐藏。让贞洁驱散不洁；让真理之光照亮欺谎的阴影；让骄傲的膨胀消退；让忿怒屈服于理性；让辱骂的利剑被粉碎，让舌头的责备被勒住；让复仇的念头消失，让伤害被遗忘。最后，让\BibleQuote{凡不是天父种植的植物，连根拔除} \BibleRef{玛 15:13}。因为只有在麦田里每棵异种被根除时，美德的种子才会在我们内好好滋养。因此，如果有人因对他人复仇的欲望所激发，以致将他投入监狱或捆绑起来，就让他赶紧宽恕，不仅宽恕无辜的人，也宽恕那看似应受惩罚的人，以便他能自信地使用主祷文中的词句说：\BibleQuote{宽免我们的债，如同我们也宽免欠我们债的人}。主特别强调这一祈求，好像整个祈祷的效力取决于这个条件，祂说：\BibleQuote{因为如果你们宽免别人的过犯，你们的天父也必宽免你们的；但如果你们不宽免别人，你们的父也不会宽免你们的过犯}。\par

% structure: p0013 source=h2 target=subsection reason=relative-heading-rank; base=h2

\subsection{六、善用四旬期将引领我们喜乐地参与复活节。}

因此，亲爱的诸位，要记住我们的软弱，因为我们容易陷入各种过犯，我们绝不能忽视这一特别的补救药和最有效的医治我们伤口的药方。让我们宽恕，以便得到宽恕；让我们赐予我们所恳求的赦免；当我们祈求被宽恕时，不要急切地想复仇。不要对穷人的呻吟充耳不闻，而要迅速以仁慈施舍给有需要的人，使我们配得在审判时获得怜悯。凡靠着天主的恩宠，全力以赴追求这完美的人，将忠诚地遵守这神圣的四旬期；脱离旧邪恶的酵母，在无酵的真诚和真理之饼中 \BibleRef{格前 5:8}，他将到达蒙福的逾越节，并因生命的更新而值得在基督我们的主内欢庆人类重生的奥迹，祂和圣父及圣神同享永生，永世无尽。阿们。\par

\SourceRecord{Translated by Charles Lett Feltoe.；Nicene and Post-Nicene Fathers, Second Series；Vol. 12.；Edited by Philip Schaff and Henry Wace.；Buffalo, NY: Christian Literature Publishing Co.,；1895.；<http://www.newadvent.org/fathers/360339.htm>.}

% structure: p0001 source=h1 target=section reason=page-title

\section{第四十讲道}

\Emph{论四旬期（二）}\par

% structure: p0003 source=h2 target=subsection reason=relative-heading-rank; base=h2

\subsection{一、进步与改善总是可能的}

亲爱的诸位，虽然随着复活节的临近，季节的复返已向我们指明了四旬斋戒，但我们的言语也必须加上劝言，在主助佑下，这些劝言对勤勉者不无裨益，对虔诚者也不是负担。因为既然这些日子的意义要求我们增加一切宗教功课，我相信，没有人不会因为被激励去行善工而感到欣喜。因为我们的本性，只要我们还属必朽，就是可变的，它在朝着最高德行的追求前进时，始终有后退的可能，也同样始终有进步的可能。而这正是成全之人真正的正义：他们绝不自以为已成全，以免在尚未完成的旅程中松懈了目标，因放弃进步的渴望而陷入失败的危险。\par

因此，亲爱的诸位，因为我们中没有一个人是那样成全和圣洁，以至于不能再更成全、更圣洁，所以让我们所有人，不分阶层，不论功绩，以虔诚的热忱，从我们已经达到之处，向着我们仍渴望达到之处奔跑，并为我们的常规虔诚功课添上一些必要的补充。因为在这些日子中不如往常更专注于宗教的人，在其他时候也显露出不够专注。\par

% structure: p0006 source=h2 target=subsection reason=relative-heading-rank; base=h2

\subsection{二、撒殚企图以新得补旧失}

因此，宗徒的宣告适时地在我们耳边响起，说：\BibleQuote{看，如今正是悦纳的时候；看，如今正是救恩的日子。}还有什么比这个时刻更蒙悦纳？还有什么日子比这些日子更适合于获得救恩？在这些日子里，向诸恶宣战，在一切德行上取得进步。你确实必须时常提防你救恩的敌人，免得任何漏洞暴露在诱惑者的圈套之下；但现在，当那个仇敌以更猛烈的仇恨向你发怒时，你必须使用更大的警惕和更敏锐的谨慎。因为如今在全世界上，他古老权势的力量已被从他手中夺走，无数被掳的器皿已从他的掌握中被解救出来。万民和万舌的民族正在挣脱他们残酷的掠夺者，如今找不到一个不反抗那暴君法律的人类种族；同时，在普天之下，千千万万的人正预备在基督内重生；当新受造物诞生临近时，属灵的邪恶正被附着的人驱逐出去。因此，被劫掠的仇敌的亵渎狂怒在焦躁不安，并寻求新的收获，因为它已失去了古老的权利。他不知疲倦，始终警觉，趁任何一只羊疏忽地从神圣的羊栈中走失时，他就攫取它，意图带领它们越过财宝的峭壁和奢华的斜坡，进入死亡的居所。于是他点燃他们的愤怒，喂养他们的仇恨，磨砺他们的欲望，嘲笑他们的节制，激起他们的贪饕。\par

% structure: p0008 source=h2 target=subsection reason=relative-heading-rank; base=h2

\subsection{三、基督的二性在受试探时显示}

因为连对我们的主耶稣基督，他也没有停止其诡诈的企图，那么还有谁他不敢试探呢？因为，正如福音记载所揭示的，当我们的救主——祂是真实的天主，为了显示祂也是真实的人，并消除一切邪恶和错误的见解——在四十昼夜的斋戒之后，经历了人性软弱的饥饿时，魔鬼因在祂身上发现了可能和必朽本性的迹象而欢喜，为要测试祂所畏惧的能力，就说：\BibleQuote{你若是天主子，就命这些石头变成饼。\BibleRef{玛 4:3}}毫无疑问，全能者能做到这事，而且造物主一声令下，受造物无论属何种类，都能容易转变成所命令的形态：正如在婚宴上，当祂愿意时，祂将水变成了酒；但在这里，更符合祂救恩计划的，是祂高傲仇敌的狡猾被主击败，不是凭借祂天主性的权能，而是凭借祂屈尊降卑的奥秘。最后，当魔鬼被打退，诱惑者在其一切计谋中受挫时，天使来到主前并侍奉祂，使祂既是真人又是真天主：祂的人性未因那些诡诈的问题而受玷污，祂的天主性因那些神圣的侍奉而得以彰显。因此，让魔鬼的子孙和门徒们蒙羞吧！他们满腹毒蛇的毒液，欺骗愚昧人，否认在基督内两种真实本性的同在，要么剥夺祂天主性中的人性，要么剥夺祂人性中的天主性；虽然这两种谬误都被一个双重的同步证据所摧毁：因为祂身体的饥饿显示了祂完美的人性，而行礼的天使则显示了祂完美的天主性。\par

% structure: p0010 source=h2 target=subsection reason=relative-heading-rank; base=h2

\subsection{四、斋戒不应仅限于禁食，而应导向善行}

因此，亲爱的诸位，既然我们由救赎者的训诫得知，\Quote{人生活不只靠饼，而也靠天主口中的每一句话}，并且基督徒无论戒律多少，理当更渴望以\Quote{天主的话}而非肉体食物来满足自己，就让我们怀着热忱的虔敬与恳切的信德，进入庆祝这隆重斋戒的礼仪，不是以无结果的禁食——那常因身体虚弱或贪吝的疾病而加于我们——而是以丰厚的慈善：使我们真能成为那真理本身所说的人：\BibleQuote{饥渴慕义的人是有福的，因为他们要得饱饫}\BibleRef{玛 5:6}。因此，但愿虔诚的善工成为我们的喜乐，并愿我们充满那滋养我们直到永生的食物。让我们因穷人的饱足而喜乐，因我们慷慨使他们满足。让我们以衣裳为乐，因我们用必需的衣物遮盖了赤身露体者。愿我们的仁爱为病患所感受，为虚弱者所感受，为流亡者的艰难所感受，为孤儿的贫乏所感受，为寡居者的忧伤所感受；在帮助这些人时，没有人不能施行一些善举。因为一个人的心大，他的收入不会小；慈悲与善行的尺度不取决于财富的多寡。善意之财从不真正缺少，即使在微薄的钱囊中。无疑，富人的花费更大，穷人的花费更小，但若行善者的心意相同，他们工作的果实并无差别。\par

% structure: p0012 source=h2 target=subsection reason=relative-heading-rank; base=h2

\subsection{五、还应促使个人改正与家庭和睦}

但是，亲爱的诸位，在这修德行善的机遇中，还有其他显著的荣冠，无需分发谷仓，无需施舍金钱，只要放荡被遏制，酗酒被弃绝，肉体的私欲被贞洁的法律所驯服：如果仇恨化为爱恋，敌意转为和平，温顺熄灭愤怒，宽恕原谅过错，总之，如果主人与仆人的行为如此井然有序，以致一方的管治更温和，另一方的纪律更完备。亲爱的诸位，正是经由这样的遵守，天主的仁慈将被获得，罪的控告被抹去，可敬的复活节得以虔诚地庆祝。而且，罗马世界的虔诚皇帝们长久以来以神圣的遵守护守着这一点；因为为尊敬主的苦难与复活，他们屈尊自己的崇高权力，放宽敕令的严厉，释放许多被囚者：这样，在普世因天主的仁慈得救的日子里，他们的仁慈，效法天上的美善，可被我们热切效仿。愿基督徒百姓因此效法他们的君王，并被这些君王的榜样在家中激励宽恕。因为私家的法律不应比公法更严厉。愿过错被宽恕，束缚被解开，冒犯被抹除，报复的企图落空，好使这神圣的庆节，借由神圣与人类的恩宠，找到所有喜乐、所有无辜的人：借着我们的主耶稣基督，他与圣父及圣神，是天主，永生永王，世世无尽。阿们。\par

\SourceRecord{Translated by Charles Lett Feltoe.；Nicene and Post-Nicene Fathers, Second Series；Vol. 12.；Edited by Philip Schaff and Henry Wace.；Buffalo, NY: Christian Literature Publishing Co.,；1895.；<http://www.newadvent.org/fathers/360340.htm>.}

% structure: p0001 source=h1 target=section reason=page-title

\section{讲道42}

\Emph{论四旬期，第四讲。}\par

% structure: p0003 source=h2 target=subsection reason=relative-heading-rank; base=h2

\subsection{一、\Emph{四旬斋戒是恢复纯洁的机会。}}

亲爱的诸位，当我提议向你们宣讲这至圣而重要的斋戒时，我最好以那位宗徒的话作为开始，因为基督在他内说话，并提醒你们我们所诵读的：\BibleQuote{看！如今正是悦纳的时候；看！如今正是救恩的日子。}\BibleRef{格后6:2}因为虽然没有一个时期不充满天主的恩赐，而且我们借着他的恩宠永远可以接近天主的仁慈，但现在，当我们被救赎的那一天再次来临，邀请我们履行一切虔诚的本分时，所有人的心灵都应被更大的热忱推动，走向灵性的进步，并以更大的信心受到鼓舞：好使我们能以洁净的身心，遵守主的苦难的超越奥迹。这些伟大的奥迹确实要求我们具有如此不懈的虔诚和不倦的敬意，以至于我们应在天主眼中始终保持不变，正如我们应当在复活节本身被找到一样。但是，因为很少有人有此恒心，而且因为考虑到肉体的软弱而放宽了更严格的遵守，并且因为一个人的利益涉及今生所有各种行为，甚至虔诚的心灵也会从尘世的灰尘中沾染一些污秽，所以天主的眷顾以极大的恩惠安排了四十天的纪律来医治我们并恢复我们心灵的纯洁，在此期间，其他时间的过犯可以通过虔诚的行为得到赎补，并通过纯洁的斋戒得以消除。\par

% structure: p0005 source=h2 target=subsection reason=relative-heading-rank; base=h2

\subsection{二、必须利用四旬期除去所有污秽，行善工不可吝啬}

因此，亲爱的诸位，当我们即将进入那些专用于斋戒的奥迹日子时，让我们注意遵守宗徒的训令，\BibleQuote{洁净自己，除去肉体和灵魂的一切污秽}\BibleRef{格后7:1}：好使我们通过控制两种本性之间的斗争，那在神引导下应管辖肉身的精神，能维护其统治的尊严：这样我们就不会得罪任何人，也不受责备者的斥责。因为如果守斋者的行为与完全纯洁的标准不符，我们就会正当地受到责难的攻击，并且由于我们的过失，不虔敬的舌头就会武装起来损害宗教。因为我们的斋戒主要不在于单纯禁食，也不仅仅是从身体欲望中撤去美味，除非心灵被召回远离恶行，舌头被约束不诽谤。这是一个温良和忍耐、和平与安宁的时期：此时一切恶习的污秽都应被根除，而美德的延续应由我们获得。现在，让虔诚的心灵大胆地习惯宽恕过错、忽略侮辱、忘记冤屈。现在，让忠信的精神以正义的武器训练自己，左右兼顾，好使良心在坚定不移的正直中不受干扰，无论在荣誉或羞辱、坏名声或好名声中，不被赞美所骄傲，也不被辱骂所疲倦。宗教的自制不应忧郁，而应真诚；从不缺少神圣喜乐慰藉的人，不应听到抱怨的咕哝。在施行慈悲工作时，不应害怕世俗财物的减少。基督徒的贫穷总是富足的，因为所有的比他缺少的更多。穷人在这个世界上的劳作也不用惧怕，因为他被赐予在万有之主内拥有一切。因此，行善的人绝不应害怕行善的能力会枯竭；因为在福音中，那寡妇的虔诚因她的两个小钱而受到了赞扬，而一杯凉水的自愿施舍也得到赏报。因为我们的爱德程度是由我们情感的真挚决定的，而怜悯他人的人永远不会缺少怜悯。撒勒法的圣寡妇发现了这一点，她在饥荒时期把她仅有的一天的食物献给了有福的厄里亚，将先知的饥饿置于自己的需要之前，毫不吝啬地放弃了少许面粉和一点油。但她没有失去她以全心所给的，而在她虔诚慷慨倒空的器皿中，新的丰裕之源产生，以致她的物产丰满并不因她用于神圣目的而减少，因为她从未害怕陷入匮乏。\par

% structure: p0007 source=h2 target=subsection reason=relative-heading-rank; base=h2

\subsection{三、如同救主一样，魔鬼也试图用我们的虔诚来陷害我们}

可是，亲爱的诸位，不要怀疑那与一切德行对抗的魔鬼，妒忌你们这些我们相信是你们自己受到激励而来的美好意愿，并为此目的，他正在磨砺他的恶毒的武力，为了使你们本身的虔诚成为自己的陷阱，并力图以自夸击败那些他未能以不信任击败的人。因为骄傲的恶习与善行相邻，傲慢总埋伏在德行近旁：因为按着所记载的，\BibleQuote{凡夸耀的，当在主内夸耀}\BibleRef{格前 10:17}，过着可赞美生活的人很难不被人的赞美所捕获。那个极恶的仇敌有什么意图他不敢攻击？有谁的禁食他不试图破坏？因为正如福音诵读中所显示的，他甚至对世界救主本身也没有停止他的诡计。因为他极其害怕那四十昼夜的禁食，他极其狡猾地想发现这种克制的能力是赐予祂的还是祂自己固有的：因为如果基督在肉体上同样受制于相同的条件，他就无需害怕他所有叛变阴谋的失败。于是他首先狡猾地试探祂是否就是万物的创造者，能够随意改变物质事物的性质；其次，在人类肉体的形态下是否隐藏着天主性，对祂来说，以空气为战车、带领祂尘世的肢体穿越空间是容易的。但当主宁愿以祂真实人性的正直来抵抗他，而不显示祂天主性的能力时，他便转向他第三个设计的狡猾，想要以帝国之欲诱惑祂——在那位身上天主权能的迹象已经失败——并应许世上的王国来引诱祂崇拜自己。但魔鬼的狡猾被天主的智慧化为愚拙，以致骄傲的仇敌被他从前所束缚的束缚了，并不害怕攻击那位理应为世界而被杀者。\par

% structure: p0009 source=h2 target=subsection reason=relative-heading-rank; base=h2

\subsection{四、悖逆者甚至将禁食转为罪}

那么，我们不仅要在口腹的诱惑中，也要在克己的目的上，警惕这仇敌的诡计。因为他知道如何藉着食物带给人类死亡，也知道如何藉着我们自己的禁食来伤害我们；他利用\OriginalTerm{摩尼教人}{Manichæans}作为工具，如同从前他驱使人取用禁止之物，如今他反过来唆使人避开允许之物。这确实是有益的操练，使人习惯于清淡饮食并克制对美味的食欲；但那些人的教条化是有祸的，他们的禁食本身竟转为罪。因为他们谴责受造物的本性，以损害造物主，并声称他们因吃那些东西而被玷污，而他们主张那些东西的创造者是魔鬼，不是天主：尽管绝对没有存在的事物是邪恶的，也没有任何事物在本质上被包含于恶中。因为良善的造物主造了万物都是好的，宇宙的创造者是惟一的，\Quote{祂创造了天、地、海和其中的万有}。其中凡赐予人作为食物和饮品的，按其种类都是圣洁而洁净的。但如果以无节制的贪婪取用，那是过度使饮食者蒙羞，并非食物或饮品的本性使他们污秽。\BibleQuote{因为凡洁净的人，一切都是洁净的；但为污秽和不信的人，没有一样是洁净的，连他们的理智和良心也都污秽了。}\BibleRef{铎 1:15}\par

% structure: p0011 source=h2 target=subsection reason=relative-heading-rank; base=h2

\subsection{五、你的禁食要合理适时}

但是你们，亲爱的诸位，公教母亲的圣洁子女，在真理的学校里受圣神的教导，要以合理的节制来规范你们的自由，知道连合法的事也应节制，并在更严格的时期区分不同的食物，意在放弃某些种类的使用，并非谴责其本性。所以，不要受那些单凭自己规条而被败坏之人的错误感染，\BibleQuote{事奉受造物而不事奉造物主}\BibleRef{罗 9:26}，并以愚蠢的禁食来事奉天上的光体：因为他们选择在每周的第一天和第二天禁食，以尊敬太阳和月亮，在他们这一件悖逆的事上，已证明他们对天主两次不忠、两次亵渎，不仅将禁食设立为对星辰的敬拜，而且蔑视主的复活。因为他们拒绝人救恩的奥迹，拒绝相信我们的主基督在我们本性的真实肉体中真正诞生、真正受苦、真正被埋葬、真正复活了。结果，他们以禁食的阴郁谴责我们喜乐的日子。并且为了掩饰他们的不信，他们竟敢参加我们的聚会，在奥迹的共融中，他们有时为了确保掩饰，以不配的嘴唇接受基督的身体，却完全拒绝饮我们救赎的血。我们将此告知你们，圣洁的弟兄们，使你们能藉这些记号辨识这类人，并使那些不敬的伪装被发现的人，藉司铎的权威从圣者的团体中被驱逐。因为被祝福的宗徒保禄以他的远见警告天主的教会说：\BibleQuote{但是弟兄们，我恳求你们，留意那些制造分裂和绊脚石、反对你们所学之道理的人，并要远离他们。因为这样的人并不事奉主基督，而是事奉自己的肚腹，并以甜言蜜语和动听的言辞欺骗无辜人的心。}\BibleRef{罗 16:17}\par

% structure: p0013 source=h2 target=subsection reason=relative-heading-rank; base=h2

\subsection{六、以生活的改善使你的禁食成为现实}

因此，亲爱的诸位，你们既已从我们的这些劝告中充分受到教诲，这些劝告我们曾多次在你们耳边反复宣讲，以抗议那可憎的谬误；现在请以虔诚的热忱进入神圣的四旬期，并准备藉你们自己的仁慈工作以赢得天主的仁慈。熄灭你们的愤怒，消除仇恨，珍惜合一，并在真正谦卑的职分上彼此争先。公道地管理你们的奴隶和属下，不要让他们中的任何人因监禁或锁链而受苦。放弃报复，宽恕过犯：以温柔代替严厉，以温和代替愤怒，以和平代替纷争。愿所有的人发现我们自制、和平、仁慈：好使我们的斋戒蒙天主悦纳。因为一句话，当我们戒绝一切邪恶时，我们就向祂献上真正节制和真正虔诚的祭献；全能的天主在一切事上帮助我们，祂与圣子及圣神共享同一神性、同一尊威，至于无穷之世。阿们。\par

\SourceRecord{Translated by Charles Lett Feltoe.；Nicene and Post-Nicene Fathers, Second Series；Vol. 12.；Edited by Philip Schaff and Henry Wace.；Buffalo, NY: Christian Literature Publishing Co.,；1895.；<http://www.newadvent.org/fathers/360342.htm>.}

% structure: p0001 source=h1 target=section reason=page-title

\section{\Emph{讲道集 46}}

\Emph{论四旬期，第八篇}\par

% structure: p0003 source=h2 target=subsection reason=relative-heading-rank; base=h2

\subsection{一、守四旬期不仅须避免身体的不洁，也须避免思想和信仰的错误}

亲爱的诸位，我们确知你们的虔诚是如此热切，以致在作为主逾越节前驱的斋戒中，你们中间许多人已先于我们的劝勉。然而，因为正确的节制工夫不仅需要肉身的克苦，也需要心灵的净化，我们期望你们的操守如此完备：当你们削减属于肉身私欲的享乐时，也应摒除来自心中妄念的错误。因为那心不被任何谬误污染的人，以真实而合理的净化，为自己准备迎接逾越节——我们全部信仰奥迹汇聚于此。诚如宗徒所言：\BibleQuote{凡不出于信德的，都是罪，\BibleRef{罗 14:23}}，那被谎言之父以幻象欺骗、又不以基督真实血肉为食粮之人的斋戒，将是无益而虚空的。因此，我们既须全心服从天主的诫命和健全的道理，也须竭尽远见，戒除邪恶的想象。因为，当这可敬的节期再度来临，整个教会普遍受劝勉去理解其救恩的奥迹时，我们那狡猾奸诈的仇敌此时更阴险地施以欺骗。那才真是基督复活的真切宣认者和敬拜者，他不混淆基督的苦难，也不在基督肉身的诞生上受骗。因为有些人以基督十字架的福音为耻，竟无耻地否认祂为救赎世界所受的惩罚，并否认主内真实肉身的本性，不明白天主圣言那不动情、不变的神性竟为人的救恩如此屈尊，以致由祂的大能不丧失自己的属性，又由祂的仁慈承担了我们的属性。因此，在基督内，有双重形像，却只有一个位格；同时是圣子和人子的天主之子是唯一的主，祂因慈爱的计划、而非因必然的法则接受了奴仆的形态，因为祂由自己的大能成为卑微，由自己的大能成为可受苦的，由自己的大能成为可死的；为摧毁罪恶与死亡的专制统治，祂内的软弱本性能够承受惩罚，而刚强本性却不失其任何光荣。\par

% structure: p0005 source=h2 target=subsection reason=relative-heading-rank; base=h2

\subsection{二、基督的一切行动彰显双重本性的临在}

因此，亲爱的诸位，当你们在阅读或聆听福音时，发现我们的主耶稣基督的某些事受到伤害，某些事因奇迹而光照，如此在同一人身上，时而人性显现，时而神性闪耀，你们切不可将这些归于幻象，仿佛在基督内只有人性或只有神性；而要忠信地相信二者，谦卑地钦崇二者；这样，在圣言与血肉的结合中就不会有分离，身体性的证据不至于显得虚幻，因为耶稣内已显明了神性的标记。关于祂两种本性的见证都是真实而丰富的，并因天主上智的深奥，这一切都达成同一目的：至圣者、不可分的圣言与可受苦的血肉联合，神性可被视为在一切事上参与血肉，血肉也参与神性。因此，那愿逃避谎言、作真理门徒的基督徒心灵，应自信地运用福音故事，仿佛仍与宗徒们一起，以属灵的领悟和肉身的视觉器官，辨识主面前的可见行为。将以下归于人：祂由女人作为婴孩出生；归于天主：祂母亲的童贞在怀孕和生育中均未受损。认出\Quote{奴仆的形态}包裹在襁褓中，躺在马槽里，但承认那是主的形态，由天使传报，\Quote{被元素宣告，}被贤士朝拜。理解为祂的人性：祂没有回避婚宴；承认那神性的：祂将水变为酒。让你们的自身情感向你们解释祂为何为死去的朋友流泪；当那同一朋友在坟墓中朽坏四天后，仅凭祂声音的命令被带活并复活时，认明祂的神能。用唾沫和泥土和泥是肉身的工作；但以此涂抹并照亮瞎子的眼睛，则是那能力的确凿标记，这能力为彰显其光荣而保留了它，并未允许在祂生命之初显现。用睡眠休息解除身体的疲劳，真正属于人性；但以斥责的命令平息狂风巨浪，真正属于神性。为饥饿者预备食物，表示人的仁慈和爱人之心；但用五个饼和两条鱼使五千人（除妇女儿童外）吃饱，谁敢否认那是神性之工？这神性借着真实肉身的机能合作，不仅在人身上显示自己，也显示人在自己内；因为人本性中古老的原初创伤，除非天主圣言从童贞玛利亚的胎中为自己取了血肉，使血肉与圣言在同一人内共存，否则无法得医治。\par

% structure: p0007 source=h2 target=subsection reason=relative-heading-rank; base=h2

\subsection{三、持守信经的宣信}

亲爱的，你们要坚定不移地持守这信德，即对主的\OriginalTerm{降生成人}{Incarnation}的信仰，藉此整个教会成为基督的身体；并且要远离异端者的一切谎言。要记住，只有当你们的心灵不受错误观念的任何玷污时，你们的慈悲善工才有益于你们，你们的严格克己才能结出果实。要弃绝此世智慧的言论，因为天主憎恨它们，没有人能藉此达至\OriginalTerm{真理}{Truth}的知识。要牢记你们在信经中所宣认的。相信天主子\OriginalTerm{与圣父同永恒}{co-eternal with the Father}，万物藉祂而造成，没有祂什么也不能造成，在末期祂也按照肉身诞生。相信祂在身体内被钉十字架、死亡、复活，并被高举于天上万有之上，坐在圣父的右边，还要以祂升天时同一样的肉身，降来审判生者死者。因为这就是宗徒向所有信众所宣讲的，说：\BibleQuote{如果你们与基督一同复活了，就追求天上的事，那里有基督坐在天主的右边；你们要思念天上的事，不要思念地上的事。因为你们已经死了，你们的生命与基督一同藏在天主内。当基督、我们的生命显现时，你们也要与祂一同显现在光荣中。\BibleRef{哥 3:1}}\par

% structure: p0009 source=h2 target=subsection reason=relative-heading-rank; base=h2

\subsection{四、利用四旬期在基督徒职责的各个方面进行全面的改善}

所以，亲爱的，既然依靠如此伟大的恩许，你们不仅在望德上，也要在行为上成为属天的人。虽然我们的心思必须时刻专注于灵魂和身体的圣洁，但现今在这四十天的斋戒期间，要更加积极地从事虔诚的善工，不仅在于施舍，这在证明悔改上非常有效，也在于宽恕过犯，并对那些被控有罪的人心怀慈悲，以免天主在我们之间所定的条件在我们祈祷时不利于我们。因为我们按照主的教导说：\BibleQuote{宽免我们的罪债，如同我们也宽免得罪我们的人\BibleRef{玛 6:12}}，我们应当全心实行我们所说的。因为唯有如此，我们接下来所求的才会实现，即我们不陷于\OriginalTerm{诱惑}{temptation}，并脱离一切\OriginalTerm{邪恶}{evils}：藉着我们的主耶稣基督，祂与圣父及圣神，永生永王，于无穷世。阿们。\par

\SourceRecord{Translated by Charles Lett Feltoe.；Nicene and Post-Nicene Fathers, Second Series；Vol. 12.；Edited by Philip Schaff and Henry Wace.；Buffalo, NY: Christian Literature Publishing Co.,；1895.；<http://www.newadvent.org/fathers/360346.htm>.}

% structure: p0001 source=h1 target=section reason=page-title

\section{讲道集 第49篇}

\Emph{论四旬期，第十一篇}\par

% structure: p0003 source=h2 target=subsection reason=relative-heading-rank; base=h2

\subsection{一、四旬期斋戒人人有责}

在所有时日与季节中，确实，可爱的诸位，天主仁慈的某些标记总是显现，没有哪一段时期缺乏神圣的奥迹，为的是使我们既然从四面八方看到救恩的证明，便能更热切地接受天主慈悲永不间断的召唤。然而，在恩宠的各种工程与恩赐中，为恢复人灵所赐予的一切，如今更清晰、更丰富地呈现在我们面前，因为此时要庆祝的不是信德的孤立部分，而是其全部。因为随着复活节的临近，人们守那最隆重、最具约束力的斋戒，其遵守毫无例外地加给所有信友；因为没有一个人圣洁到不需要变得更圣洁，也没有一个人虔诚到不可能更虔诚。因为，在这生命的不确定中，谁能被认为既不受诱惑，也无过失呢？有谁不希望添加自己的德行，或除去自己的恶习呢？既然逆境伤害我们，顺境败坏我们，而且完全得不到我们想要的东西与完全拥有它同样危险。在财富的丰裕中有陷阱，在贫穷的窘迫中也有陷阱。前者使我们骄傲自大，后者激起我们的怨言。健康考验我们，疾病考验我们，因为前者滋生粗心，后者滋生忧愁。平安中有陷阱，恐惧中也有陷阱；无论那沉溺于世俗思想的灵魂是被快乐占据还是被忧虑占据，都不重要；因为既在空虚的享乐中憔悴，又在折磨人的焦虑中劳苦，同样是不健康的。\par

% structure: p0005 source=h2 target=subsection reason=relative-heading-rank; base=h2

\subsection{二、大路拥挤，救恩的窄路几乎空无一人}

因此，那真理的保证完全应验了，我们由此得知：\BibleQuote{通往生命的门是窄的，路是狭小的}\BibleRef{玛 7:14}；而通往死亡的路是宽的，挤满了大队的人，走在安全路上的人却寥寥无几。为何左边的路比右边的更拥挤，岂不是因为众人倾向于世俗的快乐和肉身的财富吗？虽然他们所渴望的短暂而不确定，但人们为了贪图快乐更愿意忍受辛劳，而不是为了热爱德行。因此，渴求可见之物的人不计其数，而偏爱永恒胜过暂时的人却几乎找不到。所以，既然蒙福的宗徒保禄说：\BibleQuote{那看得见的，是暂时的；那看不见的，才是永远的}\BibleRef{格后 4:18}，德行的道路在某种程度上是隐藏而隐蔽的，因为\BibleQuote{我们得救，还是在于望德}\BibleRef{罗 8:24}，真信德最爱那无需任何肉身的介入就能获得的一切。因此，要使我们任性的心远离一切罪恶，并在周围无数享乐的诱惑下，不让心灵的力量向任何攻击让步，这是一项伟大的工作和辛劳。谁\BibleQuote{触摸沥青，而不被它玷污}\BibleRef{德 13:1}？谁不被肉身削弱？谁不被尘土弄脏？最后，谁有那样的纯洁，以至于不被那些没有它们就无法生活的事物所污染？因为神圣的教导借着宗徒的口命令：\BibleQuote{有妻子的，要像没有一样；哭泣的，要像不哭泣的；欢乐的，要像不欢乐的；购买的，要像一无所得的；享用这世界的，要像不享用的，因为这世界的局面正在逝去。}因此，那在贞洁的清醒中度其朝圣时光、不在所经之事上停留的心灵是有福的，它如同地上的旅客而非拥有者，既不缺乏人的情感，又安息于天主的恩许。\par

% structure: p0007 source=h2 target=subsection reason=relative-heading-rank; base=h2

\subsection{三、撒殚在此季节被激起新的努力}

而且，可爱的诸位，没有哪个时期比现在更需要并赋予这种刚毅，因为借着遵守特别的严格，养成了一个必须坚持的习惯。因为你们很清楚，这是魔鬼在整个世界狂怒的时期，基督徒的军队必须与他战斗，任何变得冷淡、懒惰或沉溺于世俗忧虑的人，现在都必须以属灵的武装装备起来，借着天上的号角燃起他们的热情投入战斗，因为那\BibleQuote{因嫉妒使死亡进入世界}\BibleRef{智 2:24}的人，现在被最强烈的嫉妒所吞噬，被最大的烦恼所折磨。因为他看见整个人类的新群体被引入天主义子的名分，新生的后裔借着教会的贞洁丰饶而倍增。他看见自己所有的暴虐权力被剥夺，从他曾经占据的心中驱逐出去，而无论男女，成千上万的老、幼、壮年人都从他手中被夺走，没有人因自己的罪或原罪而被排除，因为成义不是凭功绩偿付，而是纯粹白白赐予的。他还看见那些失足、被他诡诈的陷阱所欺骗的人，在忏悔的泪水中被洗净，并借着宗徒的钥匙打开仁慈之门，被接纳和好之恩。此外，他感到主的苦难之日临近，他被那十字架的力量所粉碎，这十字架在基督——完全无罪的那一位——身上，是世界的赎价，而非罪的刑罚。\par

% structure: p0009 source=h2 target=subsection reason=relative-heading-rank; base=h2

\subsection{四、按照天主诫命的标准自我省察是四旬期的适当操练}

因此，为使烦扰的仇敌的恶意无法借其愤怒得逞，必须唤起更热切的虔诚去遵行天主的诫命，好使我们能以心灵与身体的准备，进入那所有天主慈悲的奥秘汇聚的时期，呼求天主的引导与帮助，好使我们能借祂成就一切，没有祂我们一无所能。因为命令加于我们，是为了让我们寻求那颁布命令者的帮助。谁也不得以自己的软弱为借口，因为那赐予意愿的，也赐予力量，正如蒙福的宗徒雅各伯所说：\BibleQuote{你们中若有人缺乏智慧，就该向那慷慨施恩予众人、不责斥人的天主祈求，就必赐给他。\BibleRef{雅 1:5}}哪一位信友不知道他应当培养哪些德行、对抗哪些恶习呢？谁是如此偏袒或不谙审断自己的良心，以致不知道什么应当除去、什么应当发展呢？确实，没有人如此缺乏理智，以致不了解他生活方式的特性，或不知道他心中的隐秘。那么，他不要在事事上自满，也不要按肉身的欢愉判断自己，而要将每一种习惯放在天主诫命的天平上，在那里，有些事被命令去做，另一些被禁止，他可以按此标准衡量他生活的行为，以真正的均衡审断自己。因为天主有计划的慈悲在祂的诫命中设立了最明亮的镜子，使人可以看见自己心灵的面貌，并认识到它与天主肖像的相符或相异；其特殊目的是，至少在我们被救赎和恢复的日子里，我们暂时抛开肉身的忧虑和不安的劳碌，从世间事务转向天上的事物。\par

% structure: p0011 source=h2 target=subsection reason=relative-heading-rank; base=h2

\subsection{五、要获得我们自身罪过的赦免，就必须宽恕他人}

但正如经上所载：\BibleQuote{我们众人在许多事上都失足\BibleRef{雅 3:2}}，所以首先应激发仁慈的情感，忘记他人对我们的冒犯；免得我们因任何报复的渴望而违背那最神圣的盟约——我们借以在天主经中约束自己的盟约。当我们说\BibleQuote{宽免我们的罪债，如同我们也宽免那些亏负我们的人}时，让我们不要在心硬中不肯宽恕，因为我们要么被报复的欲望所占据，要么被温和的宽仁所充满。对于常处于诱惑危险中的人来说，更渴望的是自己的过犯无需受罚，而不是别人的过犯受罚。还有什么比在教会里，也在所有人的家中，都存在罪过的赦免，更适合基督徒信仰的呢？让恐吓被放下，让束缚被解开；因为那不肯解开它们的，将用它们把自己捆绑得更惨烈。因为一个人对另一个人所决定的，他依据自己的条款对自己宣告。而\BibleQuote{怜悯人的人是有福的，因为他们要蒙受怜悯\BibleRef{玛 5:7}}。祂在审判中是公正而仁慈的，容许一些人处在别人的权下，为的是在公平的管理下，纪律的有益性和宽仁的仁慈得以保存，并且没有人敢拒绝宽恕他人的缺点，而这宽恕正是他自己希望得到的。\par

% structure: p0013 source=h2 target=subsection reason=relative-heading-rank; base=h2

\subsection{六、与仇敌和好以及施舍也是四旬期的责任}

此外，正如主所说：\BibleQuote{缔造和平的人是有福的，因为他们要称为天主的子女\BibleRef{玛 5:9}}，所以让所有纷争和敌意都被放下，不要让任何人以为，忽略恢复兄弟间的和平还能在逾越节庆中有份。因为在高天之上与圣父一起，那不与众弟兄共融于爱德中的，将不被算作祂的儿子。此外，在施舍和照顾穷人的事上，让我们基督徒的斋期丰盛而富有；让每个人将他所禁绝的美味施舍给软弱和贫困的人。务要注意，使众人同声赞美天主；让那施舍一部分财产的人明白，他是天主慈悲的仆人；因为天主将穷人的事交在慷慨的人手中，为使那些借洗礼的水或忏悔的泪水得以洗净的罪，也能借施舍而被涂抹；因为经上记载：\BibleQuote{水能灭火，施舍能除罪\BibleRef{德 3:30}}。经由我们的主耶稣基督，等等。\par

\SourceRecord{Translated by Charles Lett Feltoe.；Nicene and Post-Nicene Fathers, Second Series；Vol. 12.；Edited by Philip Schaff and Henry Wace.；Buffalo, NY: Christian Literature Publishing Co.,；1895.；<http://www.newadvent.org/fathers/360349.htm>.}

% structure: p0001 source=h1 target=section reason=page-title

\section{讲道 第51篇}

\Emph{在四旬期第二主日前星期六——关于显圣容，圣玛窦福音 \BibleRef{玛 17:1-13} 的讲道}\par

% structure: p0003 source=h2 target=subsection reason=relative-heading-rank; base=h2

\subsection{一、伯多禄的宣认引领至显圣容}

亲爱的诸位，那藉我们肉身的耳朵到达我们心灵内在听觉的福音选段，召叫我们去理解一个伟大的奥迹；我们赖天主恩宠的帮助，若能将注意力转向之前所述之事，便更能达至这一理解。\par

人类的救主耶稣基督，在建立那召叫恶人归义、死人复活的信德时，常以劝诫性的教导和奇迹般的行为教诲祂的门徒，为使他们相信祂——基督——既是天主的\OriginalTerm{独生子}{Only-begotten of God}，又是人子。因为二者缺一，对救恩毫无助益；若相信主耶稣基督只是没有人性的天主，或只是没有\OriginalTerm{天主性}{Godhead}的人，两者同样危险，因为二者必须同被宣认，正如在祂的天主性内存在着真实的人性，同样在祂的\OriginalTerm{人性}{Manhood}内也存在着真实的天主性。因此，为坚固他们对这一信德最有益的认知，主曾在众人对祂的各种意见中，问祂的门徒们自己相信或认为祂是谁；当时宗徒伯多禄因至高圣父的启示，超越有形之物，以心灵之眼超拔人性，看见祂是永生天主之子，并宣认了天主性的光荣，因为他所看的并非血肉之躯的本质；基督对如此崇高的信德深为喜悦，以致伯多禄领受了圆满的祝福，并被赋予那不可侵犯的磐石的神圣坚固，教会要建立在这磐石上，征服地狱之门和死亡的律法；以至于在解或缚任何人的恳求时，在天上唯那藉伯多禄的判决所定夺的，方得生效。\par

% structure: p0006 source=h2 target=subsection reason=relative-heading-rank; base=h2

\subsection{二、同上（续）}

然而，亲爱的诸位，这崇高备受赞扬的理解，也必须在基督较低层次的本性之奥迹上受到教导，免得宗徒的信德，虽被高举至宣认基督内的\OriginalTerm{天主性}{Deity}的光荣，却认为接受我们的软弱与那不能受苦的天主不相称、不合宜，并相信人性在祂内已如此光荣化，以至于不能承受刑罚或在死亡中解体。因此，当主说祂必须上耶路撒冷，从长老、司祭长和经师们受许多苦，并在第三日复活时，蒙受天上光照、心中因宣认而火热的真福伯多禄，以他自以为忠诚坦率的蔑视，拒绝了那些嘲弄的侮辱和最残酷死亡的耻辱，却被耶稣慈爱的斥责所制止，并被激励去分享祂的苦难。因为紧接着的救主劝言注入并教导了：凡愿意跟随祂的人，应否认自己，并在永恒之事的希望中视短暂之物的损失为轻；因为唯那不惧为基督丧失自己灵魂的人，才能拯救自己的灵魂。为使宗徒们全心怀有这种愉快的恒常勇气，对背起十字架的严酷毫不颤抖，不以基督的惩罚为耻，也不认为祂所受的于自己有辱（因为受苦的苦楚应在不蔑视祂光荣权能的情况下彰显），耶稣带着伯多禄、雅各伯和他的兄弟若望，与他们一同上了一座极高的山，向他们显示了祂光荣的明亮，因为虽然他们已在祂身上认识了天主的威严，却不知道祂身体的力量——祂的天主性就涵蕴其中。因此，正确而意味深长地，祂曾应许旁边的一些门徒在未尝死味之前，必要看见\Quote{人子在自己的王国里来临}，即在君临的辉煌中，这辉煌特别属于祂所取人性的本性，祂愿意向这三个人显明。因为天主性本身那不可言喻、不可接近的景象，是为心地纯洁之人保留至永生的，他们尚披戴可朽血肉之躯时，决不可能注视并看见。因此，主在蒙选的见证人前显示祂的光荣，并将祂与他人共有的形像装扮得如此辉煌，以致祂的面容如同太阳的光辉，祂的衣服洁白如雪。\par

% structure: p0008 source=h2 target=subsection reason=relative-heading-rank; base=h2

\subsection{三、显圣容的目的与意义}

在这显圣容中，首要目的是要从门徒心中除去十字架的绊脚石，并防止他们的信德因祂自愿受难的屈辱而动摇，因为向他们揭示祂隐藏尊威的卓越。但同样不乏远见的是，也为圣教会的望德奠下基础，使基督的整个身体都能意识到它将要接受的改变的特性，并使肢体能应许自己分沾那已在他们元首身上闪耀的光荣。关于这一点，主在论及祂来临的威严时曾亲自说：\BibleQuote{那时，义人将在他们父的国里发光如同太阳\BibleRef{玛 13:43}。}而真福宗徒保禄也证明同一件事，说：\BibleQuote{因为我断定现时的苦楚，与将来在我们身上要显示的光荣，是不能相较的\BibleRef{罗 8:18}。}又说：\BibleQuote{因为你们已经死了，你们的生命与基督一同藏在天主内；当基督，我们的生命显现时，那时，你们也要与祂一同出现在光荣之中\BibleRef{哥 3:3}。}但为坚固宗徒并助成他们一切知识，那奇迹还传递了更进一步的教导。\par

% structure: p0010 source=h2 target=subsection reason=relative-heading-rank; base=h2

\subsection{四、梅瑟和厄里亚显现的意义}

对于\Person{梅瑟}和\Person{厄里亚}，即法律和先知，显现并同主说话；在这五个人面前，最真实地应验了所说的：\Quote{\BibleQuote{凭两三个见证人，每句话都得证实} \BibleRef{申 19:15}}。还有什么比这话更稳固、更坚定呢？旧约和新约的号角在宣告这话时联合，古代见证的文献与福音的教导结合。因为两盟约的篇章相互印证，那位在奥秘帷幕下由先前的预象所应许的，如今因现时荣耀的光辉而清晰地显明。因为正如\Person{若望}所说：\Quote{\BibleQuote{法律藉梅瑟而赐予，恩宠和真理却由耶稣基督而来} \BibleRef{若 1:17}}。在他身上，预象的应许和法律规条的目的都得以实现：因为他既以自己的临在教导预言的真理，又藉恩宠使诫命成为可能。\par

% structure: p0012 source=h2 target=subsection reason=relative-heading-rank; base=h2

\subsection{\Person{伯多禄}的建议违反天主的次序}

因此，宗徒\Person{伯多禄}因这些奥秘的启示而激动，轻视世俗、蔑视尘世，被一种对永恒事物的狂热渴望所攫住，对整个显现充满狂喜，渴望在因主荣耀的显现而蒙福的地方与\Person{耶稣}同住。因此他说：\Quote{主，我们在这里真好！你若愿意，就让我们搭三个帐棚，一个为你，一个为\Person{梅瑟}，一个为\Person{厄里亚}。}对于这个提议，主没有回答，表示他所想要的虽非邪恶，却违反天主的次序：因为世界除非藉\Person{基督}的死，不能得救；并且藉主的榜样，信友们被召叫相信，尽管对幸福的应许不应有任何怀疑，但我们应该明白，在此生试炼中，我们应祈求坚忍之力而非荣耀，因为统治的喜乐不能先于受苦的时刻。\par

% structure: p0014 source=h2 target=subsection reason=relative-heading-rank; base=h2

\subsection{圣父从云中发声的意义}

当他还在说话时，忽然有一片光耀的云彩遮蔽了他们，并有声音从云中说：\BibleQuote{这是我的爱子，我所喜悦的，你们要听从他！}圣父确是在圣子内，而在主的光辉（他已调和以使门徒能看见）中，圣父的性体并没有与独生子分离：但为了强调双重位格，正如圣子身体的荣光使圣子显现于他们眼前，同样圣父的声音从云中发出使圣父被他们听见。当这声音发出后，\BibleQuote{门徒就俯伏在地，非常害怕。}他们因圣父和圣子的威严而颤抖：因为现在他们更深地领悟了两位不可分的天主性；在恐惧中，他们没有将一位与另一位分离，因为他们的信德没有怀疑。因此，这是一个广阔而丰富的见证，包含比耳朵所听到的更圆满的意义。因为当圣父说：\BibleQuote{这是我的爱子，我所喜悦的，等等}，难道不是清楚地说：\BibleQuote{这是我的儿子}，他的本质是从我而来、与我同在的吗？因为生者不先于受生者，受生者也不后于生者。\BibleQuote{这是我的儿子}，他在天主性、能力、永恒上都不与我分离。\BibleQuote{这是我的儿子}，不是被收养的，而是真生的，不是从别处受造，而是由我所生；也不是从另一种本性而像我，而是由我的本性所生，与我同等。\BibleQuote{这是我的儿子}，\BibleQuote{万物是藉着他造的；凡受造的，没有一样不是藉着他造的}，因为凡我所做的事，他也照样做；我所作的一切，他都与我不可分离且毫无差异地一同做：因为子在我内，我也在子内，我们的合一从不分开；虽然我是生者，他是被我所生者，但你们若对他有任何不能归于我的想法，便是错误。\BibleQuote{这是我的儿子}，他不是以强夺的方式，也不是以贪婪的方式寻求那与我的平等，而是存于我荣耀的形态中，为了实施我们共同救赎人类的计划，他将不变的天主性降低至奴仆的形象。\par

% structure: p0016 source=h2 target=subsection reason=relative-heading-rank; base=h2

\subsection{我们所应听从的是谁}

\BibleQuote{你们要听从他}因此，毫不犹疑地听从他在我内完全喜悦的那一位，藉他的宣讲我得以显现，藉他的卑微我得以受光荣；因为他是\BibleQuote{真理和生命}，他是我的\BibleQuote{能力和智慧}。\BibleQuote{你们要听从他}，法律奥秘所预言的，先知的口所歌咏的那一位。\BibleQuote{你们要听从他}，他用他的血救赎世界，他束缚魔鬼，夺走他的财物，他摧毁罪的枷锁和过犯的契约。听从他，他开启天堂之路，藉十字架的刑罚为你预备登天的阶梯？你为什么因被救赎而战栗？你为什么害怕你的伤口被治愈？愿基督所愿的和我所愿的成就。抛弃一切肉性的恐惧，以信德的恒心武装自己；因为你不应当在救主的苦难中害怕那藉他恩赐你连在自己结束时也无需害怕的事。\par

% structure: p0018 source=h2 target=subsection reason=relative-heading-rank; base=h2

\subsection{圣父的话对整个教会具有普遍意义}

这些事，亲爱的诸位，并不仅仅是为了那些亲耳听见的人们的益处而说的，而是藉着这三位宗徒，整个教会都学到了他们亲眼所见、亲耳所闻的一切。愿所有人的信德，都按照至圣福音的宣讲而坚定，愿无人以基督的十字架为耻，因为世界正是藉着它被救赎的。愿无人因畏惧为义而受苦，或对许诺的应验产生怀疑，因为我们要经过劳苦才能抵达安息，经过死亡才能获得生命；我们的一切卑微软弱都由祂所承担，只要我们持守对祂的承认与爱慕，我们就必如祂一样得胜，并领受祂所应许的。因为，无论是为了履行祂的命令，还是为了忍受逆境，圣父预先宣告的声音应常在我们的耳中回响：\BibleQuote{这是我的爱子，我所喜悦的，你们要听从他！}祂和圣父及圣神，永生永王。阿们。\par

\SourceRecord{Translated by Charles Lett Feltoe.；Nicene and Post-Nicene Fathers, Second Series；Vol. 12.；Edited by Philip Schaff and Henry Wace.；Buffalo, NY: Christian Literature Publishing Co.,；1895.；<http://www.newadvent.org/fathers/360351.htm>.}

% structure: p0001 source=h1 target=section reason=page-title

\section{讲道 54}

\Emph{论受难，第三篇；在复活节前的星期日讲道。}\par

% structure: p0003 source=h2 target=subsection reason=relative-heading-rank; base=h2

\subsection{一、基督的两重本性}

在天主慈悲的一切工程中，亲爱的诸位，从起初赐予人类救恩的，没有比基督被钉在十字架上为世界更奇妙、更崇高的了。因为此前的所有世代的奥迹都指向这一奥迹，天主旨意在祭献、先知预兆和法律规条中所规定的一切变化，都预示这是确定的，并应许了它的实现：以致如今预像和象征已终结，我们在相信那已经完成的事上得利，而从前我们在期待它而得利。因此，亲爱的诸位，在一切关于我主耶稣基督受难的事上，天主教的信德坚持并要求我们承认在救主内两性相遇，且在其属性存留的同时，两性的结合得以完成，以致从那时起，按照人类需求所要求的，在荣福童贞的胎中，\Quote{圣言成了血肉}，我们不可将祂视为没有那属人的天主，也不可视为没有那属天的世人。每一性确实藉着其不同的行动表达其实在的存在，但各自并不脱离与另一性的联系。任何一方都无所欠缺；在尊威中谦卑圆满，在谦卑中尊威圆满；合一并不带来混淆，区分也不破坏合一。一方可受苦难，另一方不可侵犯；然而屈辱属于同一主体，荣耀亦然。祂同时存在于软弱和大能中；同时能死，又是死亡的征服者。因此，天主取了完整的全人，并以祂的慈悲和大能将两性如此结合，以致每一性都存在于另一性中，且都没有越出其自身属性而进入另一性。\par

% structure: p0005 source=h2 target=subsection reason=relative-heading-rank; base=h2

\subsection{二、两性共同行动，人的苦难并非被迫，而是符合天主旨意}

但因为那在永恒之前为我们恢复而制定的奥迹的设计，若不藉着人的软弱和天主的德能，就无法完成，所以两种\Quote{形态}都共同实施各自所特有的事，即圣言实行圣言的事，血肉实行血肉的事。一是以奇迹闪耀，另一则屈服于伤害。一不离开与父荣耀的平等，另一不离开我们本性的特质。然而，即使是祂对苦难的忍受，也并未如此使祂参与我们的谦卑，以至于与天主性的德能分离。恶者的疯狂加诸于主的一切嘲弄和侮辱、一切迫害和痛苦，并非因必然所忍受，而是因自由意志所承担：\BibleQuote{因为人子来，是为寻找并拯救那遗失的\BibleRef{路 19:10}：}祂如此利用迫害者的恶行来救赎众人，以致在祂死亡与复活的奥迹中，连那些谋杀者，若他们相信，也能得救。\par

% structure: p0007 source=h2 target=subsection reason=relative-heading-rank; base=h2

\subsection{三、\Person{犹达斯}的恶名从未被超越}

因此，\Person{犹达斯}，你被证明比众人更罪大恶极、更不幸；因为当悔改本应召你回到主那里时，绝望却把你拖向绞索。你应该等待你罪行的完成，并在基督的血为所有罪人流出来之前，推迟你那可怕的绞刑。在主的许多奇迹和恩赐中，本可唤醒你的良心，至少那些神圣的奥迹——你在逾越节晚餐中所接受的，当时甚至已因神性知识的记号而被察觉出卖——本可挽救你于轻率的坠落。你为何不信赖祂的良善？祂并未将你从祂圣体圣血的共融中驱逐，当你与人群和持武器的一队人前来捉拿祂时，祂也没有拒绝给你和平之吻。啊，无可挽回的人！啊，\Quote{一去不返的灵！}你随从你内心的狂怒，魔鬼站在你右边，将你为所有圣者性命所准备的恶行转向你自己的毁灭，以致因为你的罪行超过了所有惩罚的尺度，你的邪恶使你成为自己的审判官，你的惩罚允许你成为自己的刽子手。\par

% structure: p0009 source=h2 target=subsection reason=relative-heading-rank; base=h2

\subsection{四、基督自愿以祂的荣耀换取我们的软弱}

因此，\BibleQuote{天主在基督内使世人与自己和好\BibleRef{格后 5:19}，}造物主亲自穿着将被恢复至造物主肖像的受造物；在施行了神性奇迹之后——这些奇迹的施行，曾由预言之神预告：\BibleQuote{那时瞎子的眼要开了，聋子的耳要听得见；那时瘸子必跳跃如鹿，哑子的舌头要说话\BibleRef{依 35:5}；}耶稣知道已到了为祂光荣的受难实现的时候，便说：\BibleQuote{我的心灵忧伤至死\BibleRef{玛 26:38}；}又说：\BibleQuote{父啊，若是可能，就让这杯离开我\BibleRef{玛 26:38}。}这些话表达了一定程度的恐惧，显示祂愿意通过分担我们的软弱来医治我们软弱的感受，并通过经历痛苦来制止我们对忍受痛苦的恐惧。因此，主在我们的本性中与我们一同颤抖，以祂自己能力的圆满完全覆盖我们的软弱和脆弱。因为祂来到这世界，是一位从天而降的富足而仁慈的商人，通过奇妙的交换，与我们订立了救恩的契约：接受我们的，赐予祂的；以尊荣换侮辱，以救恩换痛苦，以生命换死亡；祂——原本有超过一万两千天使大军可以为消灭迫害者而服务——更愿接纳我们的恐惧，而不是运用祂自己的能力。\par

% structure: p0011 source=h2 target=subsection reason=relative-heading-rank; base=h2

\subsection{五、\Person{圣伯多禄}首先受益于他师傅的贬抑}

如此多的羞辱赐予所有信友，至福的宗徒伯多禄是最先证明的：在猛烈的威胁与残酷的惊吓使他动摇之后，他迅速转变，恢复力量，从伟大的典范中寻得疗愈，以致忽然动摇的肢体回归元首的坚固。因为奴仆不能\Quote{大于主，门徒也不能大于师傅}，若那战胜死亡者未先畏惧，他便无法战胜人性的战栗。因此，主\Quote{回头看伯多禄}，在司祭的毁谤、见证的谎言、鞭打祂并唾弃祂之人的凌辱中，以那双已预见其惊慌的眼睛注视那惊慌的门徒；真理的目光进入他心中，他必须被矫正，仿佛主的声音在那里回响，说道：\Quote{伯多禄，你要往哪里去？为何退缩于己？转向我，信赖我，跟随我：这是我受难的时刻，你受苦的时刻尚未来到。为何畏惧你也将战胜的事？不要让那与我同享的软弱使你困惑。我曾为你畏惧；你当对我满怀信心。}\par

% structure: p0013 source=h2 target=subsection reason=relative-heading-rank; base=h2

\subsection{VI. 犹太人的疯狂计谋反倒成了自取灭亡}

\BibleQuote{到了早晨，众司祭长和民间的长老聚集商议，要处死耶稣\BibleRef{玛 27:1}。}这个早晨，犹太人，对你们而言不是日出，而是日落；惯常的白昼并未光顾你们的眼目，但最黑暗之夜笼罩着你们邪恶的心。这个早晨为你们推倒了圣殿及其祭台，废除了法律和先知，毁灭了王国和司祭职，使你们的一切庆节变为永久的哀悼。因为你们议定了一个疯狂而流血的计谋，你们这些\Quote{肥牛犊}，你们这些\Quote{许多公牛}，你们这些\Quote{咆哮}的野兽，你们这些狂犬，要将生命的主宰和荣耀的主交付死亡；并且，仿佛你们暴怒的极恶能因借助那统治你们省份之人的判决而减轻，你们把捆绑的耶稣带到比拉多的审判台前，使那惊恐的法官被你们持续的呼喊压倒，竟选择释放一个凶手，而要求将世界的救主钉在十字架上。基督被判刑后——这更多是由于比拉多的怯懦而非权力，他用洗净的手却以污秽的唇说出那宣告祂无罪的同一张嘴，将耶稣送上了十字架——百姓放纵自己，顺从司祭的眼色，向主堆积许多侮辱，疯狂的暴民将愤怒倾泻在祂身上，而祂温柔地自愿忍受一切。可是，亲爱的弟兄们，由于整段故事过长，今日无法全部讲述，让我们将余下的部分推迟到星期三，那时将重读主的受难史。因为主将应允你们的祈祷，使我们靠着祂自己的恩赐完成我们的许诺：藉着我们的主耶稣基督，祂永生永王。阿们。\par

\SourceRecord{Translated by Charles Lett Feltoe.；Nicene and Post-Nicene Fathers, Second Series；Vol. 12.；Edited by Philip Schaff and Henry Wace.；Buffalo, NY: Christian Literature Publishing Co.,；1895.；<http://www.newadvent.org/fathers/360354.htm>.}

% structure: p0001 source=h1 target=section reason=page-title

\section{第五十五篇讲道}

\Emph{第四讲：论主的受难，圣周三宣讲。}\par

% structure: p0003 source=h2 target=subsection reason=relative-heading-rank; base=h2

\subsection{一、两个强盗的悔改与亵渎之间的差异是人类的一个预象。}

亲爱的诸位，我们理应满足你们的期望，这必须赖主对你们祈祷的丰厚回应来实现：那使你们热切祈求的，也必使我们胜任所托。\par

在最近论述主的受难时，我们达到了福音叙述中\Person{比拉多}屈服于犹太人恶喊，要求将耶稣钉十字架的那一点。因此，当天主性隐藏在脆弱肉身中所允许的一切都完成时，天主之子耶稣基督被钉在他也曾背负的十字架上，两个强盗同样被钉，一个在祂右边，另一个在左边；以便即使在十字架的事件中也显示出祂在审判中必须对所有的人所作的分别：因为信道的强盗的信仰是那些将得救之人的预象，而亵渎者的邪恶预表了那些将被定罪的人。因此，基督的受难包含我们救恩的奥秘，而犹太人的不义为惩罚祂所准备的工具，救赎主的大能却为我们变成了登上光荣的阶梯：主耶稣为所有人的得救如此承受了这受难，以至于当祂被钉在木头上悬挂时，祂为杀害祂的人恳求圣父的仁慈，并说：\BibleQuote{圣父，赦免他们，因为他们不知道他们所做的是什么}\BibleRef{路 23:34}。\par

% structure: p0006 source=h2 target=subsection reason=relative-heading-rank; base=h2

\subsection{二、在嘲弄中，司祭长们显出了对圣经的完全无知。}

但是，救主为之祈求宽恕的司祭长们却用嘲弄的利刺加重了十字架的酷刑；对于祂，他们无法再用手发泄更大的愤怒，便用舌头的武器攻击祂，说：\BibleQuote{祂救了别人，却不能救自己。祂是以色列的王，现在从十字架上下来，我们就信祂}\BibleRef{玛 27:42}。从什么错误的泉源，从什么仇恨的池塘，你们这些犹太人，啜饮这样恶毒的亵渎？什么师傅教导了你们，什么教训使你们确信，你们应当相信祂是以色列的王和天主子，而祂要么不允许自己被钉十字架，要么摆脱了捆绑的钉子？法律的奥秘、逾越节的神圣仪式、先知的口从未告诉你们这一点：而你们确实真实且屡次读到那适用于你们可憎的恶行和主的自愿受苦的经文。因为祂亲自通过\Person{依撒意亚}说：\BibleQuote{我把我的背转给鞭打，我的面颊转给掌掴，我没有转脸躲避唾污的羞辱。}\BibleRef{依 50:6} 祂亲自通过\Person{达味}说：\BibleQuote{他们给我苦胆作食物，我渴时，他们给我酸醋喝}，又说：\BibleQuote{许多恶犬围绕我，恶徒的集会包围我。他们刺透了我的手和脚，我数算了我的所有骨头。他们却注视着我，他们分了我的外衣，为我的长袍拈阄。}而且，恐怕你们自己的恶行的过程似乎已被预言，而受难者身上没有预表任何权能，你们确实没有读到主从十字架上下来，但你们读到了：\BibleQuote{主在木头上为王了。}\par

% structure: p0008 source=h2 target=subsection reason=relative-heading-rank; base=h2

\subsection{三、十字架的胜利立时且有效。}

因此，基督的十字架象征着预像的真正祭台，在这祭台上，藉着带来救恩的牺牲，人性的奉献应被举行。在那里，无玷羔羊的血涂抹了远古过犯的后果；在那里，魔鬼仇恨的整个暴政被粉碎，谦卑胜过骄傲的抬举而得胜光荣：因为信仰的效果如此迅速，以至于与基督同钉的强盗中，那相信基督是天主子的人，成义后进入了乐园。谁能揭示如此巨大恩惠的奥秘？谁能陈述如此奇妙改变的力量？在顷刻之间，长期作恶的罪债被消除；他紧贴着十字架，在他挣扎的灵魂的残酷折磨中，他转向基督；对他，自己的邪恶带来了惩罚，现在基督的恩宠赐予了冠冕。\par

% structure: p0010 source=h2 target=subsection reason=relative-heading-rank; base=h2

\subsection{四、当悲剧的最后一幕结束时，犹太人该如何感受？}

然后，已经尝了那醋，即那葡萄园的产物，这葡萄园不顾神圣的栽种者而堕落，变成了外邦葡萄的酸涩，主说：\BibleQuote{完成了。}；也就是说，圣经被满足了：我再也没有狂暴民众的愤怒要忍受了：我已经承受了我预言要受的一切。软弱的奥秘完成了，让权能的证明产生吧。于是祂低下头，交付了灵魂，将那应在第三天复活的身体赐予了安眠的安息。当生命的作者经历这奥秘的阶段，且在天主尊严如此大的屈尊时，整个世界的根基被震动，当一切受造物以其动荡谴责他们的邪恶罪行，且世界的元素本身对罪犯宣告了明确的判决时，你们这些犹太人，当宇宙的判决反对你们，你们的邪恶无法被撤销（因为罪行已经完成），你们有什么想法，有什么内心反省？什么困惑遮盖了你们？什么折磨抓住了你们的心？\par

% structure: p0012 source=h2 target=subsection reason=relative-heading-rank; base=h2

\subsection{五、在预备复活节共融时，贞洁与爱德是最需要的两件事。}

因此，亲爱的诸位，天主的仁慈如此伟大，竟愿意使甚至来自这样一个民族的许多人因信德成义，并将我们这些一度在旧日愚昧的深沉黑暗中丧亡的人，收纳于圣祖的行列和选民之中，所以让我们不要懈怠，也不要懒惰，而要攀登希望的顶峰；并要明智而忠信地反省：我们是从怎样的俘虏之境、从何等悲惨的奴役之中，以何等的赎价被赎回，被何等强有力的手臂领出，好让我们在我们的身体上光荣天主：使我们甚至藉着我们生活方式的正义，显示出祂居住在我们内。因为没有比仁慈的温柔和纯洁的贞洁更高贵或更卓越的德行了，所以让我们特别以这些武器装备自己，好使我们能够从地上被举起，如同藉着积极爱德和光明洁德的双翼，获得天上的位置。凡藉着天主的恩宠充满这渴望，并因自己的进步而不自夸，却夸耀于主的人，便适当地尊崇逾越奥迹。毁灭的天使不会越过他的门槛，因为门框上涂有羔羊的血和十字架的标记。他不畏惧埃及的灾祸，并留下那些被同一水淹没的敌人，而他自己正是藉此水而得救。所以，亲爱的诸位，让我们以净化了的心灵和身体，拥抱我们救恩的奇妙奥迹，并从一切\Quote{旧恶的酵母，让我们庆祝}主的逾越节，以合宜的敬礼：这样，在圣神的引导下，我们可能\Quote{隔离}不受任何试探\Quote{与基督的爱，}基督藉祂的血使万有和平，已回归圣父荣耀的崇高，却仍不弃舍服事祂的人的卑微。愿光荣与尊崇归于祂，直到永远。阿们。\par

\SourceRecord{Translated by Charles Lett Feltoe.；Nicene and Post-Nicene Fathers, Second Series；Vol. 12.；Edited by Philip Schaff and Henry Wace.；Buffalo, NY: Christian Literature Publishing Co.,；1895.；<http://www.newadvent.org/fathers/360355.htm>.}

% structure: p0001 source=h1 target=section reason=page-title

\section{讲道集 58}

\Emph{（论受难，第七篇）}\par

% structure: p0003 source=h2 target=subsection reason=relative-heading-rank; base=h2

\subsection{一、基督在逾越节受难的理由}

亲爱的诸位，我确实知道，逾越节庆典所分沾的奥秘是如此崇高，不仅超越我卑微悟性的浅薄领会，甚至超越伟大心智的能力。但我不可因天主工作的伟大而怀疑或羞于履行我的职责；因为即使无法解释，我们也不可对人得救的奥秘保持沉默。但靠你们的祈祷相助，我们相信天主的恩宠将被赐予，用祂启示的甘露浇灌我们心灵的不毛之地；好使牧者之口宣扬那些对圣羊群之耳有益的事。因为当赏赐一切美善的主说：\Quote{你张开口，我要使它充满。}我们也敢用先知的话回答说：\Quote{主，求祢开启我的嘴唇，我的口要传扬祢的赞美。}因此，亲爱的诸位，再次着手讲述主受难的福音叙事时，我们明白这是天主计划的一部分：那亵渎的犹太首领和不圣洁的司祭，他们曾多次寻找机会向基督发泄愤怒，却不得不在逾越节时才获得行使暴怒的能力。因为那些长久以来在奥秘预象中所应许的事，必须全然清楚应验；例如，那真正的羔羊必须取代作为其预象的羔羊，那唯一的祭献必须终结众多不同祭献。因为所有那些通过梅瑟关于羔羊祭献的神圣规定，都预表了基督，并真实宣告了基督的被杀。因此，为使预象让位于实体，使预象在实体面前消失，古老的礼规被新的圣事取代，牺牲变成了祭品，血被血擦除，法律规定的庆节因改变而得以圆满。\par

% structure: p0005 source=h2 target=subsection reason=relative-heading-rank; base=h2

\subsection{二、犹太首领不但废除了自己的法律，而且在毁灭基督时未能领会新的安排}

因此，当大司祭们召集经师和民间长老到公议会，所有司祭的心思都专注于对耶稣行恶的意图时，法律教师们便使自己置身法律之外，并因自愿的失职而废除了祖传的礼仪。因为当逾越节庆开始时，那些本应装饰圣殿、清洁器皿、准备祭品、以更神圣的热心从事法律所规定的洁净之礼的人，却被叛徒般的仇恨所狂怒，全神贯注于一项工作，以一致的残酷图谋同一罪行，尽管他们注定因惩罚无辜和定正义的罪而一无所获，除了未能领会新奥秘和违反旧礼仪。因此，首领们为防止圣日发生骚乱\BibleRef{玛 26:5}，所表现的热心并非为了庆节，而是为了滔天罪行；他们的焦虑不是服务于宗教事业，而是自陷于罪。因为这些谨慎的大司祭和焦虑的司祭害怕在重要庆节发生暴动，不是怕百姓做错事，而是怕基督逃脱。\par

% structure: p0007 source=h2 target=subsection reason=relative-heading-rank; base=h2

\subsection{三、耶稣建立圣体圣事，对叛徒犹达斯仁慈到底}

但耶稣坚定于祂的旨意，毫不畏惧地执行父的旨意，完成了新约，建立了新的逾越节。因为当门徒与祂一同斜倚在神秘晚餐时，当盖法厅中正讨论如何处死基督时，祂制定了祂体血的圣事，教导他们该向天主奉献何种祭品，而且连那出卖者也没有被排除在这奥秘之外，为表明他并非因个人委屈而被激怒，而是早已自愿决定背叛。因为他自己是他败坏的根源和背信的原因，跟随魔鬼的引导，拒绝让基督做导师。因此，当主说：\Quote{我实在告诉你们，你们中有一人要出卖我。}祂表明祂清楚知道出卖者的良心，不是用严厉或公开的斥责来使叛徒困惑，而是用温和沉默的警告来对待他，好使那从未因被拒绝而迷失的人，更易于因悔改而矫正。唉，不幸的犹达斯，你为何不利用如此大的容忍？看，主宽恕了你邪恶的企图；基督不向任何人揭露你，只向你自己揭露。你的名字和面貌未被发现，只有你心中的秘密被真理和仁慈的话语触及。宗徒品位的尊荣没有被拒绝，圣事的分沾也没有被拒绝。回心转意吧；抛开你的疯狂，变得明智。慈悲邀请你，救恩敲门，生命召唤你回到生命。看，你那清白无罪的同伴们对你的罪行暗示感到战栗，所有人为自己颤抖，直到背叛的作者被揭露。因为他们悲伤并非由于良心的控告，而是由于人善变的无常；害怕每个人所知道的自己的罪行不如真理本身所预见的真实。但你在圣人们的恐惧中滥用主的耐心，相信你的厚颜隐藏了你。你在罪上加添无耻，不被如此明确的考验所吓倒。当其他人停止吃那主已设定审判的食物时，你却没有从盘中缩回你的手，因为你的心没有被从罪行中扭转。\par

% structure: p0009 source=h2 target=subsection reason=relative-heading-rank; base=h2

\subsection{四、受难中的各种事件进一步阐释，以及基督苦难的真实性得以肯定}

因此，亲爱的诸位，正如福音作者\Person{若望}所叙述的，当主将那蘸过的饼递给负卖祂的人，更清楚地指明他时，魔鬼完全占据了犹达斯，并且现在，藉着他真实的邪恶行为，占据了那早已被其恶计束缚的人。因为他只是身体与同席的人躺在一起，心灵却在酝酿司祭的仇恨、证人的虚伪和愚昧暴民的狂怒。最后，主看见犹达斯执意犯下何等严重的罪行，说：\BibleQuote{你所要做的，快去做吧！ \BibleRef{若 13:27}} 这不是命令的声音，而是许可的声音；不是恐惧的声音，而是准备好的声音。祂，那掌管万有的，显示祂不阻拦背叛者，并如此施行圣父为救赎世界的旨意，既不促进也不害怕迫害者所预备的罪行。因此，当犹达斯在魔鬼的怂恿下离开基督，并自宗徒团体的合一中割裂时，主不受任何恐惧的干扰，只一心关切祂来救赎之人的得救，将不受迫害者攻击的所有时间用于奥妙的交谈和神圣的教导，正如\BibleBook{若望}{福音}所记载的：举目向天，为整个教会恳求圣父，使圣父所有并要赐给圣子的人都合而为一，并保持不分，以归于救赎者的光荣；最后又加上那祈祷，祂在其中说：\BibleQuote{父啊！如果可能，就让这杯离开我吧！ \BibleRef{玛 26:39}} 在此不可认为主耶稣想要逃避受难与死亡——祂已将这圣事的奥迹交托给门徒们保管——因为祂亲自禁止满怀虔诚信德与爱德而激动抽剑的\Person{伯多禄}，说：\BibleQuote{父所赐给我的杯，我岂能不喝呢？ \BibleRef{若 18:11}}；并且也因主根据若望福音所说的那句话是确定的：\BibleQuote{天主竟这样爱了世界，甚至赐下了自己的独生子，使凡信祂的人不至丧亡，反而获得永生 \BibleRef{若 3:16}}；同样，宗徒\Person{保禄}也说：\BibleQuote{基督爱了我们，并为我们舍弃了自己，作为献给天主的馨香牺牲。 \BibleRef{弗 5:2}} 因为藉着基督的十字架拯救众人，是圣父与圣子共同的意愿与共同的计划；那在永恒之前已慈悲地决定并不可更改地预定了的事，绝不能受到任何干扰。因此，在取了真实而完整的人性时，祂接受了身体真实的感受和心灵真实的情感。并不能因为祂的一切充满圣事、充满奇迹，就推论祂流出虚假的眼泪，或假装饥饿而进食，或假装睡眠。祂是在我们的谦卑中被轻看，因我们的忧伤而忧伤，带着我们的痛苦被钉在十字架上。因为祂的怜悯承受了我们必朽的苦难，目的在医治它们；祂的大能迎战它们，目的在征服它们。\Person{依撒意亚}最清楚地预言了这一点，说：\BibleQuote{祂承担了我们的罪，为我们而受痛苦；我们以为祂是在痛苦中，在鞭伤中，在折磨中。但祂是为了我们的过犯而被刺透，为了我们的罪恶而被击伤；因祂的创伤我们得了痊愈。 \BibleRef{依 53:4-5}}\par

% structure: p0011 source=h2 target=subsection reason=relative-heading-rank; base=h2

\subsection{五、基督的顺服是教会不朽的教训}

因此，亲爱的诸位，当天主子说：\BibleQuote{父啊！如果可能，就让这杯离开我吧！ \BibleRef{玛 26:39}} 祂运用了我们本性的呼喊，为人类的脆弱与战栗辩护：好使我们的忍耐得以坚固，恐惧得以驱除，在我们必须承受的事中。最后，既然祂已在一定程度上缓和了我们软弱的恐惧，便不再求这个——我们保留它们并无益处——祂转入另一种心境，说：\BibleQuote{然而，不要照我的意愿，而要照你的意愿。 \BibleRef{玛 26:39}} 并且又说：\BibleQuote{如果这杯不能离开我，除非我喝它，就愿你的旨意成就。 \BibleRef{玛 26:42}} 头的这些话是整个身体的救恩：这些话教导了所有的信友，激发了所有宣信者的热忱，加冕了所有的殉道者。因为谁能战胜世界的仇恨、诱惑的狂风、迫害者的恐怖呢？除非基督以全体之名并为全体向圣父说了：\BibleQuote{愿你的旨意成就？ \BibleRef{玛 26:42}} 那么，所有教会之子——那些以如此高价购买、如此白白地成义的人——都应学习这些话；当某种猛烈诱惑的冲击临到他们时，让他们使用这强有力祈祷的帮助，好能战胜恐惧与战栗，学习耐心忍受。从这一点起，亲爱的诸位，我们的讲道必须转入对主受难细节的默想，并且为了避免以冗长令你们厌烦，我们将分次进行共同的任务，将其余部分推迟到周三。天主的恩宠必会赐予你们，如果你们祈求祂赐给我履行职务的能力：以上所求，是靠我们的主耶稣基督，等等。\par

\SourceRecord{Translated by Charles Lett Feltoe.；Nicene and Post-Nicene Fathers, Second Series；Vol. 12.；Edited by Philip Schaff and Henry Wace.；Buffalo, NY: Christian Literature Publishing Co.,；1895.；<http://www.newadvent.org/fathers/360358.htm>.}

% structure: p0001 source=h1 target=section reason=page-title

\section{讲道五十九}

\Emph{（论受难，第八篇：圣周三）}\par

% structure: p0003 source=h2 target=subsection reason=relative-heading-rank; base=h2

\subsection{一、基督的被捕成就了祂自己的永恒旨意}

亲爱的诸位，在我们上次讲道中讨论了主被捕之前发生的事件，如今靠天主的恩宠，照我们所承诺的，该讨论受难本身的细节了。当主藉祂神圣祈祷的言语清楚表明，神性和人性真实而圆满地存在于祂内，显示不愿受苦出自一面，而决心受苦则出自另一面——即排除一切软弱的恐惧并坚固其崇高的大能——之后，祂便回归祂的永恒旨意，以无罪的\Quote{奴仆的形状}面对那藉犹太人双手猛烈攻击祂的魔鬼：使那位独有无罪人性的那一位，担当众人的案情。因此，黑暗之罪攻击了真光，尽管他们带着火把和灯笼，却仍逃不脱自己不信的黑夜，因为他们不认识光明之源。他们逮捕祂，祂已准备好被捉拿；他们带走祂，祂甘心被带领；因为，若祂愿意抵抗，他们邪恶的手绝不能伤害祂，但如此将阻碍世界的救赎，而祂本要为众人的得救而死，却一个也救不了。\par

% structure: p0005 source=h2 target=subsection reason=relative-heading-rank; base=h2

\subsection{二、比拉多在容许自己被犹太人误导一事上的罪过何其重大}

因此，祂容许百姓的狂怒——由司祭煽动——所敢做的一切加诸自身，被带到亚纳斯那里，他是盖法的岳父；然后亚纳斯将祂解送到盖法那里。在诬告者的疯狂指控和伪证者的谎言之后，两位大司祭的代表将祂移交到比拉多处听审；他们既忽视天主的法律，又高喊\Quote{除了凯撒，我们没有君王}，仿佛他们效忠于罗马法律，将全部判决权留给了总督，实则是寻求施行残酷的帮手，而非案件的仲裁。因为他们将耶稣捆绑得紧紧的，被许多拳打耳光打得伤痕累累，被吐唾沫，已为他们的喊声所定罪；以致在他们自己判决的众多迹象中，比拉多不敢宣告这位众人皆欲其死的人无罪。实际上，审讯本身就表明他在被告身上找不到任何罪过，而且在他的判决中他没有坚持自己的主见：因为他作为法官宣判了一位他宣告为无罪的人有罪，并求那不义之民承担那正直者的血——这血，他凭自己的确信感到，并从他妻子的梦中得知，自己不可沾染。洗手不能洗净那污秽的灵魂，洒水的手指也不能赎那邪恶心思所赞同的罪行。比拉多的过错确实小于犹太人的罪行，因为是犹太人用凯撒的名字恐吓他，用仇恨的言语责骂他，迫使他犯下恶行。但他也未能免于被指控，因为他迎合了那些鼓噪者，放弃了自己的判断，顺从了他人的控告。\par

% structure: p0007 source=h2 target=subsection reason=relative-heading-rank; base=h2

\subsection{三、然而犹太人的罪过远为重大}

因此，亲爱的诸位，比拉多屈从于那不可和解的百姓的疯狂，容许耶稣受许多嘲弄的侮辱，受过度凌辱的折磨，并将祂——为鞭子撕裂、荆棘冠冕、穿着嘲弄的紫红袍——显示在迫害者的眼前，无疑是想平息仇敌的心，使他们饱尝残酷仇恨之后，不再进一步迫害这位遭受如此多样磨难的人。但他们的愤怒仍在熊熊燃烧，他们向他喊叫，要求释放巴拉巴，这样耶稣就承担了十字架的刑罚；于是，当众人同意地低语说\Quote{祂的血归在我们和我们子孙身上}时，那些恶者便照他们一贯要求的那样，为自己得到了定罪的结局。\Quote{他们的牙齿是武器和箭矢，他们的舌头是利剑}，正如先知所见证的。因为他们徒然避免亲手钉死荣耀的主，却用致命的舌箭和话语的毒箭攻击祂。在你们身上，虚假的犹太人和不圣洁的百姓首领，这罪恶的全部重量都落在你们身上；尽管这恶行的巨大也涉及总督和士兵，但你们是首要和主要的罪犯。在基督的定罪中，无论比拉多的判决或执行他命令的军营所行的任何不义，都只能使你们更令人憎恨，因为你们狂怒的冲动甚至不让那些对你们不义行为不悦的人免于罪责。\par

% structure: p0009 source=h2 target=subsection reason=relative-heading-rank; base=h2

\subsection{四、基督背负自己的十字架是给教会永恒的教训}

于是主被交给他们野蛮的意愿，并为了嘲弄祂的王权，被命令背负祂自己的死亡工具，为使依撒意亚先知所预见的得以应验，说：\Quote{看，有一个婴孩诞生了，有一个儿子赐给了我们，祂的权柄落在祂的肩膀上。}因此，当主背着那要成为祂权柄令牌的十字架木时，在恶人眼中固然是极大的嘲弄，但在信者心中却展示了一个伟大的奥迹：因为祂，这荣耀的克胜魔鬼者，这击溃敌对权势的强者，正以高贵的姿态背着祂胜利的战利品，用祂不可征服的忍耐的肩膀，将救恩可敬的记号带入所有领域。仿佛甚至那时祂已藉祂工作的这一单纯象征来坚固所有跟随祂的人，并说：\Quote{谁不背起自己的十字架跟随我，就不配是我的。} \BibleRef{玛 10:38}\par

% structure: p0011 source=h2 target=subsection reason=relative-heading-rank; base=h2

\subsection{五、十字架从主转移到基勒乃人西满身上，象征着外邦人参与祂的苦难}

但当群众同耶稣前往刑场时，遇见一个基勒乃人西满，将十字架的木头放在他身上，代替主背着，藉此预先预示外邦人的信德，因为基督的十字架对他们不是羞辱，而是光荣。所以，这不是偶然的，而是象征性和奥妙的：当犹太人向基督发怒时，一个外邦人被找来分担祂的苦难，正如宗徒所说：\BibleQuote{若我们与祂一同受苦，也必与祂一同为王。}\BibleRef{弟后 2:12} 因此，不是希伯来人，也不是以色列人，而是一个外邦人，代替了救主承受祂至圣的屈辱。因为藉此转移，无玷羔羊的赎罪祭和一切奥迹的应验，从割礼之人传到了未受割礼之人，从按肉体的子孙传到了按精神的子孙：正如宗徒所说：\BibleQuote{我们逾越节的羔羊基督已经被杀献祭了。}\BibleRef{格前 5:7} 祂将自己作为崭新而真实的和好祭品献给圣父，被钉十字架，不是在圣殿内——那里的敬拜已告结束，也不是在被判为有罪而必遭毁灭的城内，而是在城外，\BibleQuote{在营外}\BibleRef{希 13:12} 目的是，当旧的象征性祭品终止时，新的祭品被放在新的祭台上，基督的十字架不是圣殿的祭台，而是世界的祭台。\par

% structure: p0013 source=h2 target=subsection reason=relative-heading-rank; base=h2

\subsection{六、我们不仅要看十字架，还要看它的意义}

因此，极亲爱的各位，基督被高举在十字架上，你们心灵的眼睛不要只停留在那些邪恶罪人看见的景象上，关于他们，梅瑟曾说过：\BibleQuote{你的性命必悬在你眼前，你昼夜恐惧，你的性命毫无保障。}\BibleRef{申 28:66} 因为他们在被钉的主身上，只想到自己的恶行，没有那种使真实信德得以成义的恐惧，而是那种折磨邪恶良心的恐惧。但我们这些被真理之神光照的心智，要以纯洁而自由的胸怀，培育那照亮天地之十字架的光荣，并用内在的目光看见主论及祂即将来临的苦难时所说的话：\Quote{人子受光荣的时候来到了。} 祂接着又说：\Quote{现在我的心灵烦乱。我该说什么呢？父啊，救我脱离这个时刻，但我正是为这个时刻而来的。父啊，光荣祢的子。} 当父的声音从天上传来，说：\Quote{我已经光荣了，还要再光荣。} 耶稣回答旁人说：\Quote{这声音不是为我，而是为你们。现在是这世界的审判，现在这世界的元首要被赶出去。至于我，当我从地上被举起来时，我要吸引众人归向我。}\par

% structure: p0015 source=h2 target=subsection reason=relative-heading-rank; base=h2

\subsection{七、十字架的能力普遍吸引人}

啊，十字架的神奇能力！啊，不可言喻的苦难光荣，其中包含了主的审判台、世界的审判和被钉者的权能！因为祢吸引万有归向祢，主！当祢整日向悖逆不信的百姓伸出双手时，反抗祢的子民\BibleRef{依 65:2}，最后整个世界都承认祢的威严。祢吸引万有归向祢，主！当时所有元素联合宣告对犹太人罪行的判决，天上的光变暗，白昼变成黑夜，大地也异常震动，整个受造界拒绝为那些恶人服务。祢吸引万有归向祢，主！因为圣殿的帐幔裂开，至圣所不再为那些不配的大司祭存在；这样，预像变成了真理，预言变成了启示，法律变成了福音。祢吸引万有归向祢，主！因此，从前在犹太人唯一的圣殿中以模糊记号举行的事，如今要由万民在全世界公开而明确地举行。因为现在有更高贵的肋未人，有更大尊荣的长老，有更神圣傅油的司铎：因为祢的十字架是一切福泽的泉源，一切恩宠的源头，藉着它，信徒获得软弱中的力量，羞辱中的光荣，死亡中的生命。现在，各式各样的肉性祭品已经停止，祢的圣体圣血唯一的奉献满足了所有不同的祭品：因为祢是真正的\BibleQuote{除免世罪的天主羔羊，}\BibleRef{若 1:29} 在祢身上完成了所有奥迹，正如只有一个祭品代替许多祭品，同样只有一个王国代替许多民族。\par

% structure: p0017 source=h2 target=subsection reason=relative-heading-rank; base=h2

\subsection{八、我们不可为自己而活，而要为那为我们死去的基督而活}

因此，诸位最亲爱的，让我们明认那万民的蒙福导师，宗徒保禄所明认的，他说：\BibleQuote{这话是确实的，值得完全接纳：基督耶稣来到世界，为拯救罪人}\BibleRef{弟前 1:15}。因为天主对我们的仁慈更加奇妙之处，在于基督不是为义人或圣洁的人而死，而是为不义和邪恶的人而死；虽然天主性的本性不能承受死亡的毒刺，但祂在诞生时却从我们取了祂能为我们奉献的。因为祂早就以祂死亡的力量威胁我们的死亡，藉欧瑟亚先知的口说：\BibleQuote{死亡啊，我将成为你的死亡；阴府啊，我将成为你的毁灭}\BibleRef{欧 13:14}。因为祂藉着死亡，经历了阴府的规律，却藉着复活打破了它们，如此摧毁了死亡的连续性，使之成为暂时的而非永恒的。\BibleQuote{因为正如在亚当内众人都死了，照样在基督内众人都要复活}\BibleRef{格前 15:22}。因此，诸位最亲爱的，愿圣保禄所说的话实现：\BibleQuote{他们活着，今后不再为自己生活，而是为那为他们死而复活的那位生活}\BibleRef{格后 5:15}。并且因为旧事已过，一切都成了新的，愿没有人停留在他旧的属肉体的生活中，而是让我们都藉着每日的进步和虔敬的增长而更新。因为一个人无论怎样成义，只要他仍在此生中，就总能更加被悦纳和更好。那不前进的，就是在后退；那无所获的，就在失去某些东西。那么，让我们以信德的步伐，藉着怜悯的工作，在义德的爱中奔跑，好使我们灵性地遵守我们救赎的日子，\BibleQuote{不可用旧酵母，即恶毒和邪恶的酵母，而要用纯洁和真诚的无酵饼}\BibleRef{格前 5:8}，从而堪当分享基督的复活，祂与圣父及圣神，永生永王，万世无疆。阿们。\par

\SourceRecord{Translated by Charles Lett Feltoe.；Nicene and Post-Nicene Fathers, Second Series；Vol. 12.；Edited by Philip Schaff and Henry Wace.；Buffalo, NY: Christian Literature Publishing Co.,；1895.；<http://www.newadvent.org/fathers/360359.htm>.}

% structure: p0001 source=h1 target=section reason=page-title

\section{讲道 62}

\Emph{（论受难，第十一篇）}\par

% structure: p0003 source=h2 target=subsection reason=relative-heading-rank; base=h2

\subsection{一、受难的奥迹超越人的理解}

我们所渴望的、普世亦当向往的主受难庆节已经来临，它不容我们在灵性喜乐的喧嚣中缄默不语：因为虽然常就同一题旨作合宜而相称的讲论是困难的，然而司铎不得回避以讲道向民众的耳朵施予有关天主慈悲这伟大奥迹的训诲，因为题目本身既不可言说，反使讲者易于出口，且所说的话绝不能全然落空，纵使永远说不够。那么，让脆弱的人性臣服于天主的荣耀，永远承认自己不堪担当祂慈悲工程的阐述。让我们在思想中劳苦，在洞见上失败，在言辞上结结巴巴：甚至我们关于主威严的正确思想也当不足，这是好的。因为，记起先知所说：\Quote{你们要寻求主，就必得刚强：不断寻求祂的面，}谁也不得以为自己已寻获所求的一切，免得他停止努力而无法靠近。而在所有令人惊叹的沉思疲惫的天主工程中，有什么比救主的受难更令我们的心灵愉悦而又困惑呢？尽管我们可能默想祂与圣父同一同等的全能，但天主内的谦卑比大能更令人惊异，而把握天主威严的完全空虚比感悟祂内\Quote{奴仆的形象}的无限提升更为困难。但借着忆起：尽管造物主与受造物、不可侵犯的天主与可受苦的肉体绝对不同，然而两种本性的属性在基督的同一位格内相遇，以致在祂软弱与大能的作为中，同一位格既领受卑下亦享荣耀，这极大地有助于我们的理解。\par

% structure: p0005 source=h2 target=subsection reason=relative-heading-rank; base=h2

\subsection{二、信经以圣伯多禄的宣认为教会基本教义}

在所爱的诸位，在信经之初所领受的这项信仰规则中，基于宗徒教导的权威，我们承认我们的主耶稣基督——我们称祂为全能天主圣父的独生子——也是由圣神从童贞玛利亚所生。当我们宣认祂被钉十字架、死亡和第三日复活时，我们并不否认祂的威严。因为凡属于天主的和凡属于人的，都由祂的人性和神性同时完成，以致因着可受苦与不可受苦的结合，祂的大能不受祂软弱的波及，祂的软弱也不被祂的大能所克服。真福宗徒伯多禄确当受称赞，因为他宣认了这结合；当时主询问门徒们对祂的认识，他迅速抢先于众人说：\BibleQuote{祢是基督，永生天主之子。}而他确实看见这一点，不是凭借血肉的启示——那可能会阻碍他内在的视觉——而是藉由圣父之神在他信者的心中工作，使他在准备治理整个教会时，先学会他将要教导的；并且为使他将要宣讲的信德得以坚固，他领受了这保证：\BibleQuote{你是伯多禄，在这磐石上，我要建立我的教会，阴间的门决不能战胜她。}因此，基督信仰的坚强——它建立于不可攻破的磐石之上，不惧怕死亡之门——承认唯一的主耶稣基督是真天主亦真人，同时相信祂是童贞玛利亚之子，又是祂母亲的造主；祂虽是时间的创造者，却也在时代的末叶出生；祂是全能的主，却又是必死的人类的一员；祂不知罪，却以有罪肉身的相似形状为罪人祭献。\par

% structure: p0007 source=h2 target=subsection reason=relative-heading-rank; base=h2

\subsection{三、魔鬼的诡计反噬自身}

为使人类脱离致命过犯的束缚，祂向狂怒的魔鬼隐藏了祂威严的大能，以我们脆弱谦卑的本性对抗他。因为如果那残忍而骄傲的仇敌能够知道天主慈悲的计划，他必会设法安抚犹太人的心使之温和，而不是以不义的怨恨煽动他们，免得他在攻击一位对他毫无亏欠者的自由时，丧失他对所有俘虏的奴役。于是他被自己的恶意挫败：他向天主之子施加了刑罚，这刑罚却转为所有人类之子的医治。他流了义血，这血成为世界的赎价与赎罪的饮品。主承担了祂按自己旨意的目的所选择的一切。祂允许狂人将罪恶的手加于祂身上：这些手在为自身毁灭效力时，却服务于救赎者的工程。然而祂甚至对杀害自己的人怀有如此深厚的慈爱，以致在十字架上向圣父祈祷，不求为自己报复，而求饶恕他们，说：\BibleQuote{父啊，宽恕他们，因为他们不知道他们做的是什么}\BibleRef{路 23:34}。而那祈祷的力量如此之大，以致许多曾说过\BibleQuote{他的血归在我们和我们的子孙身上}\BibleRef{玛 27:25}的人，由宗徒伯多禄的宣讲转为痛悔，一日之内约有三千犹太人受洗：他们都成了\BibleQuote{一心一意}\BibleRef{宗 4:32}的人，现在已准备好为曾要求将他钉死的那一位而死。\par

% structure: p0009 source=h2 target=subsection reason=relative-heading-rank; base=h2

\subsection{四、为何犹达斯不能通过基督获得宽恕}

这宽恕是叛徒犹达斯所不能达到的：因为他，这丧亡之子，魔鬼站在他右边，在基督完成普世救赎的奥迹之前，就自陷于绝望。因为主为罪人而死，或许他本可找到救恩，如果他未曾匆忙去上吊。但那颗邪恶的心，时而沉溺于偷窃的欺诈，时而忙于背信弃义的计谋，从未对救主仁慈的明证怀有任何领会。这两只邪恶的耳朵曾听到主说话，当祂说：\BibleQuote{我来不是召义人，而是罪人 \BibleRef{玛 9:13}}，和\BibleQuote{人子来，是为寻找及拯救迷失了的人 \BibleRef{路 19:10}}，但这些话并未让他的心领会基督的仁慈，祂不仅治愈了身体的软弱，也医治了病弱灵魂的创伤，对瘫子说：\BibleQuote{孩子，放心罢！你的罪赦了 \BibleRef{玛 9:3}}；也对被带到祂面前的犯奸淫的妇人说：\BibleQuote{我也不定你的罪；去罢，从今以后不要再犯罪了}，为要在祂的一切工程中显示祂是作为世界的救主，而非审判者而来。但这邪恶的叛徒拒绝理解这一点，并对自己采取了措施，不是出于悔改的自责，而是出于丧亡的疯狂，因此那位将生命之作者卖给杀害祂的人的人，甚至在死亡中增加了那定他罪的罪的分量。\par

% structure: p0011 source=h2 target=subsection reason=relative-heading-rank; base=h2

\subsection{五、基督被钉十字架的残酷被其奇妙的德能所消融}

因此，那些假见证人、残酷的民众首领、邪恶的司祭通过一个懦弱的总督和无知的兵丁对主耶稣基督所做的一切，自古以来既是可憎恶的，也是可欣喜的。因为虽然主的十字架是犹太人残酷意图的一部分，但它因他们所钉死的那一位而具有奇妙的德能。民众的愤怒针对一个人，而基督的仁慈却为全人类。他们所施加的残酷，祂自愿承受，为使祂永恒旨意的工程能通过他们不受阻碍的罪行得以实现。因此，福音中所详尽叙述的整个事件顺序，信友必须如此接受：通过默然相信主受难时所发生的一切，我们应当明白，不仅罪的赦免由基督完成，而且正义的标准也得到了满足。但为能更详尽地讨论这一点，愿我们靠主的帮助将这部分的主题保留到本星期的第四天。我们盼望，天主的恩宠将因你们的恳求而赐予，帮助我们实现我们的承诺：藉着我们的主耶稣基督，等等。阿们。\par

\SourceRecord{Translated by Charles Lett Feltoe.；Nicene and Post-Nicene Fathers, Second Series；Vol. 12.；Edited by Philip Schaff and Henry Wace.；Buffalo, NY: Christian Literature Publishing Co.,；1895.；<http://www.newadvent.org/fathers/360362.htm>.}

% structure: p0001 source=h1 target=section reason=page-title

\section{讲道 63}

\Emph{(论主的受难，第十二篇：于星期三宣讲)}\par

% structure: p0003 source=h2 target=subsection reason=relative-heading-rank; base=h2

\subsection{一、天主选择以在软弱中成就的全能来拯救人}

诸位亲爱的，主的受难的荣耀（我们曾允诺今天再次谈论）首要地因其谦卑的奥迹而令人惊叹，这奥迹既赎回了我们，又教训了我们，使那位付出赎价者也将祂的义德分施给我们。因为天主圣子的全能——祂以同一性体与圣父相等——本可仅凭祂的意愿就将人类从魔鬼的辖制下解救出来，但若不符合神圣的作为：以曾被战胜的事物来战胜仇敌邪恶的抵抗，并以那使整个族类陷入奴役的本性来恢复我们本性的自由，那就更好了。然而，当福音作者说：\Quote{\BibleQuote{圣言成了血肉，居住在我们中间} \BibleRef{若 1:14}}，以及宗徒说：\Quote{\BibleQuote{天主在基督内使世界与自己和好} \BibleRef{格后 5:19}}时，这表明至高圣父的独生子进入了与人性谦卑的如此结合，以至于当祂取了我们的肉身和灵魂的实体时，祂仍是同一的天主圣子，高举了我们的特性，而非祂自己的特性：因为需要被加强的是软弱，而非能力，以致在受造者与造物主结合时，被取者中不应缺少任何神圣性，取者中也不应缺少任何人性。\par

% structure: p0005 source=h2 target=subsection reason=relative-heading-rank; base=h2

\subsection{二、天主的计划始终被部分理解，如今普遍适用}

这仁慈与公义的天主计划，虽然在过去的世代中多少被笼罩在黑暗里，却并未完全隐藏，以至于从一开始直至主来临前最堪受赞美的圣人们无法理解它：因为藉着基督而来的救恩，已由先知的话语和事件的象征所预许，不仅预言它的人获得了这救恩，所有相信他们预言的人也都获得了。因为同一个信德使历代圣人们成义，并且因耶稣基督——天主与人之间的中保——我们所承认成就的一切，或我们祖先虔敬接受为将要成就的一切，都属于信者同一的望德。犹太人与外邦人之间没有区别，因为正如宗徒所说：\Quote{\BibleQuote{割礼不算什么，未受割礼也不算什么，只在于遵守天主的诫命} \BibleRef{格前 7:19}}，如果以圆满的信德遵守这些诫命，它们就使基督徒成为亚巴郎真正的子孙，即完美的，因为同一位宗徒说：\Quote{\BibleQuote{因为你们凡在基督耶稣内受洗的，就是穿上了基督。不再有犹太人或希腊人，不再有奴隶或自由人，不再有男人或女人，因为你们众人在基督内合而为一。如果你们属于基督，那么你们就是亚巴郎的后裔，按照恩许作承继人} \BibleRef{迦 3:27}}。\par

% structure: p0007 source=h2 target=subsection reason=relative-heading-rank; base=h2

\subsection{三、神圣的元首与肢体的结合不可分离}

因此，诸位亲爱的，毫无疑问，人的本性已被天主圣子接纳为如此结合，以至于不仅在祂——那首生于一切受造物的人——内，而且在祂所有的圣人内，都有同一的基督，并且正如元首不能与肢体分离，肢体也不能与元首分离。因为虽然不在今生，而是在永恒中，天主将是\Quote{\BibleQuote{万物中的万有} \BibleRef{格前 15:28}}，但甚至现在，祂也是祂圣殿——即教会——不可分离的居住者，正如祂亲自应许说：\Quote{\BibleQuote{看！我天天与你们在一起，直到今世的终结} \BibleRef{玛 28:20}}。与此一致，宗徒说：\Quote{祂是身体——教会——的元首，祂是元始，是死者中的首生者，为使祂在一切事上居首位，因为天主乐意使一切圆满（天主性）居住在祂内，并藉着祂使万有与自己和好。} \BibleRef{哥 1:18-20}\par

% structure: p0009 source=h2 target=subsection reason=relative-heading-rank; base=h2

\subsection{四、基督的受难提供了救赎的奥迹和当效法的榜样}

这些以及许多其他引证向我们的心提示了什么？无非是我们在一切事上应更新为祂的肖像，祂虽保持着\Quote{\BibleQuote{天主的形体} \BibleRef{斐 2:6}}，却屈尊\Quote{取了罪过的肉身的形体}。因为祂承担了我们由罪恶而来的一切软弱，却不染罪恶，所以祂感到饥、渴、睡、疲劳、忧伤、哭泣，并忍受最剧烈的痛苦直至死亡，因为除非那唯独全人类本性无罪的祂容许被恶人之手杀害，无人能从死亡的罗网中解脱。因此，我们的救主天主子为所有信祂的人预备了奥迹与榜样，使他们藉着重生领悟奥迹，藉着效法追随榜样。因为蒙福的宗徒伯多禄教导此事说：\Quote{基督为我们受了苦，给你们留下榜样，使你们追随祂的足迹。祂没有犯过罪，口中从未出过诡诈。祂被辱骂时，不还口；受苦时，不恐吓，却将自己交托给不义的审判者。祂亲身在木头上承担了我们的罪，使我们死于罪恶，而活于正义。} \BibleRef{伯前 2:21-24}\par

% structure: p0011 source=h2 target=subsection reason=relative-heading-rank; base=h2

\subsection{五、基督没有废除，而是成全并提升了法律}

亲爱的诸位，既然没有一位信徒被拒绝领受恩宠的恩赐，同样也没有人在基督徒的纪律方面不是欠债者；因为，虽然奥秘的法律的严厉已被除去，但自愿遵守它的益处却增加了，正如圣史若望所说：\Quote{因为法律是藉梅瑟传授的，但恩宠和真理却是由耶稣基督而来的\BibleRef{若 1:17}}。因为一切依照法律先行的事——无论是肉身的割损、大量的祭献，还是遵守安息日——都是为基督作证，并预兆基督的恩宠。祂就是\Quote{法律的终向\BibleRef{罗 10:4}}，不是废除，而是成全其意义。尽管祂同时是旧约和新约的作者，但祂改变了与许诺相关的象征礼仪，因为祂实现了许诺，并以被预报者的来临终结了预报。然而，在道德诫命方面，旧约的规条并未被废除，反而有许多被福音教训所拓展；因此，带来救恩的事物比那些应许救主的事物更完美、更清晰。\par

% structure: p0013 source=h2 target=subsection reason=relative-heading-rank; base=h2

\subsection{第六章 基督苦难的当下效果每日由基督徒体验，尤其在圣洗圣事中}

因此，天主圣子为世界与天主和好所行所教的一切，我们不仅知道是过去的历史，更在它现在功效的力量中体会。祂就是那一位，由童贞母亲因圣神降生，并用同一圣神使祂无玷的教会丰产，好藉圣洗圣事的重生，使无数子女生于天主，关于他们经上说：\Quote{他们不是由血气，也不是由肉欲，也不是由男欲，而是由天主生的\BibleRef{若 1:13}}。祂就是那一位，因全世界的义子名分，亚巴郎的苗裔得了祝福\BibleRef{创 22:18}，族长藉信德而非血肉之子的出生，成了万民之父——即恩许之子。祂就是那一位，不排除任何民族，从天下万国中组成一个圣羊群，每日实现祂所许诺的，说：\Quote{我还有别的羊，还不属于这一栈，我也该把他们引来，他们要听我的声音，这样，将只有一个羊群，一个牧人\BibleRef{若 10:16}}。因为尽管祂首先对真福伯多禄说：\Quote{喂养我的羊}，然而唯一的主指导所有牧人的职责，并以如此丰美水草丰盛的牧场喂养那些来到这磐石的羊，以致无数的羊群因祂爱情的丰盈而得滋养，并毫不迟疑为牧人的缘故而牺牲，正如善牧自己甘心为羊舍命。祂的苦难不仅由殉道者光荣的勇气所分担，也在重生的行动中由所有新生者的信德所分担。因为弃绝魔鬼和信仰天主、从旧状态过渡到新生命、脱去属地的形像、穿上属天的样式——这一切是一种死亡和复活，使那被基督接纳并接纳基督的人，在来到洗礼池之后与之前不再是同一人，因为重生者的身体成了被钉者的身体。\par

% structure: p0015 source=h2 target=subsection reason=relative-heading-rank; base=h2

\subsection{第七章 基督徒的善行只是基督善行的一部分}

亲爱的诸位，这改变是至高者的手笔，祂\Quote{在一切人身上行一切事}，以致藉着在每个信徒身上所见的良好生活方式，我们认出祂是一切义行的创始者，并感谢天主的仁慈，祂以无数恩宠的恩赐如此装饰整个教会身体，以致通过同一光的多条光线，同样的光辉遍照各处，而任何基督徒所做的好事，都不能不是基督荣耀的一部分。这就是那成义并光照每个人的真光。这就是那从黑暗权势中拯救我们并转移我们进入天主圣子国度的那一位。这就是那藉着新生命提升心神愿望并熄灭肉体情欲的那一位。这就是那藉着以\Quote{纯洁和真诚的无酵饼}，藉着弃绝\Quote{邪恶的旧酵母\BibleRef{格前 5:8}}，并以主自己来陶醉和喂养新受造物，从而恰当地庆祝主的逾越节的那一位。因为领受基督的体血所带来的，无非是我们转变为我们所领受的，并在灵魂和身体上处处携带祂——我们曾在祂内、与祂同死、同葬、同复活，正如宗徒所说：\Quote{因为你们已经死了，你们的生命与基督一同藏在天主内。当基督，你们的生命，显现时，你们也要与祂一同出现在光荣中\BibleRef{哥 3:3}}。祂与圣父等。\par

\SourceRecord{Translated by Charles Lett Feltoe.；Nicene and Post-Nicene Fathers, Second Series；Vol. 12.；Edited by Philip Schaff and Henry Wace.；Buffalo, NY: Christian Literature Publishing Co.,；1895.；<http://www.newadvent.org/fathers/360363.htm>.}

% structure: p0001 source=h1 target=section reason=page-title

\section{讲道集 67}

\Emph{（《论受难》第十六篇，主日宣讲。）}\par

% structure: p0003 source=h2 target=subsection reason=relative-heading-rank; base=h2

\subsection{一、默想基督受难的预言，是虔诚喜乐的伟大泉源}

亲爱的诸位，信友的心思本当时常惊叹天主的化工，并特别将理性能力用于那些能增进信德的反思。因为，只要虔诚的心注意普遍享有的恩惠或祂的恩宠的特殊恩赐，就能远离许多虚荣，从肉身的牵挂退入精神的幽居。但在主受难的季节，这更应当热切而彻底地实行，好使那时在神圣读经中所读的，确能以悟性的耳朵领受，并且那些言语伟大的题目，因着其下的奥迹真相，更显伟大。因为使我们举心向上的首要原因，乃是先知们的歌声早已咏唱过福音真理所叙述的事，不是作为将要发生，而是作为已经发生；而人的耳朵尚未得知将要成就的事，圣神却已宣告为实现了。因为达味王——按肉躯说，基督是他的后裔——在主被钉十字架之日一千一百多年前就已结束了他的一生，并未遭受他所述及加于自身的任何刑罚。但是，因为借他的口发言的那一位，要从他的血统取得受难的肉躯，所以十字架的故事正确地在祂肉体祖先的身上预先应验了。因为达味确实在基督内受苦，因为耶稣确实在从达味所取的肉躯内被钉十字架。\par

% structure: p0005 source=h2 target=subsection reason=relative-heading-rank; base=h2

\subsection{二、天主的预知并不使犹太人的恶行得以推诿}

既然犹太人的不虔敬对威严之主所做的一切，早在如此久远以前就被预言了，而先知的语言与其说是关乎未来，不如说是关乎过去，那么，这除了向我们启示天主永恒旨意的不可改变的计划以外，还有什么呢？在祂那里，将要决定的事已决定，将要发生的事已成就。因为，既然我们行为的特质和一切愿望的实现都为天主所预知，那么祂自己的化工，祂岂不更知晓吗？祂认为正确地记录那些无可阻碍必定成就的事，如同已经成就一样。因此，当宗徒们充满圣神，遭受基督仇敌的恐吓和虐待时，也同心合意地向天主说：\Quote{确实，在这城中，黑落德和般雀·比拉多，连同外邦人和以色列众民，聚集起来反对祢的圣仆耶稣，祢所傅油的，要行祢的手和祢的旨意所预定要成就的事。}那么，难道基督迫害者的恶行出自天主的计划？那无可比拟的罪行竟是由天主的手预备和推动的吗？显然，我们绝不可这样想至高的公义：关于犹太人的恶意所预知的，与关于基督受难所预定的，是大为不同，甚至完全相反的。他们杀害祂的欲望，与祂受死的意愿，并非出于同一根源；他们的滔天罪行与救赎主的忍受，也不是同一位圣神的产物。主并未激起，而是允许那些狂徒的凶恶之手；祂虽预知必须成就的事，却没有强制成就，即便祂正是为成就此事而取了肉躯。\par

% structure: p0007 source=h2 target=subsection reason=relative-heading-rank; base=h2

\subsection{三、基督绝非杀害祂者的罪过的根源}

事实上，被钉者的情形与钉死祂者的情形是如此不同：基督所承担的，不可逆转；而他们所行的，却可被抹去。因为那来拯救罪人的，连对杀害祂的人也拒绝施恩，反将恶人的邪恶转化为信徒的善行，为使天主的恩宠越发奇妙，由于仁慈地施行恩宠，不是根据人的功德，而是根据天主智慧和知识的丰富宝藏，因为连那流了救主血的人，也被接纳到圣洗的水中。因为，正如圣经——它记载了宗徒的行实——所载，当有福的宗徒伯多禄的讲道刺透犹太人的心，他们承认了自己罪行的不义，说：\Quote{诸位弟兄，我们该做什么？}那同一位宗徒说：\Quote{你们悔改，各人因耶稣基督之名领洗，以赦免你们的罪，并领受圣神的恩赐。因为这恩许是为你们和你们的儿女，并为一切远方的人，凡我们的主天主所召叫的，}不久之后，圣经接着说：\Quote{于是，领受他话的人就领了洗；那一天，约增加了三千人。}\BibleRef{宗 2:37}因此，主耶稣基督在愿意忍受他们狂暴的愤怒时，完全不是他们罪行的根源；祂并未强迫他们渴望这事，而是允许他们能够如此，并利用盲目民众的疯狂，正如祂利用背叛者犹达斯的奸诈一样——祂曾以善言善行恩赐他，将他从所设想的那可怕罪行中召回：收纳他为门徒，擢升他为宗徒，以征兆警告他，并接纳他参与神圣奥迹的启示，好使那未缺少任何仁慈纠正之恩的人，毫无犯罪的借口。\par

% structure: p0009 source=h2 target=subsection reason=relative-heading-rank; base=h2

\subsection{四、犹达斯罪行的巨大}

但哦，你这最不虔敬的人，\Quote{你是客纳罕的苗裔，而非犹大的苗裔}，不再\Quote{是蒙拣选的器皿}，而是\Quote{丧亡与死亡之子}，你曾想魔鬼的煽动于你更为有利，以致被贪婪的火炬点燃，急切想得到三十块银钱，却没有看见你将失去的财富。因为即使你不认为主的许诺是可信的，又有什么理由偏爱这么小的一笔钱，胜过你已经领受的呢？你曾习惯命令邪魔，医治病人，与其他宗徒同受尊荣，而且为满足你对获利的渴望，你可以从你所保管的钱囊中偷窃。但你那贪恋禁物的心，更被那不大允许的事强烈刺激：使你喜悦的，与其说是价银的数目，不如说是罪恶的极大。因此，你邪恶的交易之所以可憎，不仅因为你把主看得如此廉价，更因为你出卖了那曾是救赎者，甚至是你自己的救赎者，却对自己毫无怜悯。你的惩罚被交在你自己的手中，是公义的，因为没有人比你更残酷地致力于你自己的灭亡。\par

% structure: p0011 source=h2 target=subsection reason=relative-heading-rank; base=h2

\subsection{五、基督的受难是为了我们的救赎，藉着奥迹与榜样。}

因此，按照祂旨意的目的，在指定的时刻，耶稣基督被钉十字架、死去并埋葬，这并非祂自身状况所需的定命，而是将我们从囚禁中救赎的方法。因为\Quote{圣言成了血肉}，为的是从童贞女的胎中祂能取我们受苦的本性，并使那不能加于天主圣子的事，能加于人子。因为虽然在祂诞生之初，天主性的标记已照耀在祂身上，祂身体成长的整个历程充满奇事，然而祂真实地承担了我们的软弱，没有罪的玷污，却未免除自己任何人性的脆弱，为能将祂所是的赐予我们，并在祂内医治我们所有的一切。因为祂，全能的医师，为我们这悲惨处境预备了双重的药方，其中一部分是奥迹，另一部分是榜样，藉着奥迹，神能得以赐予；藉着榜样，人性的软弱得以驱除。因为正如天主是我们成义的根源，人也有义务向祂奉献虔诚。\par

% structure: p0013 source=h2 target=subsection reason=relative-heading-rank; base=h2

\subsection{六、我们只能藉着跟随基督的足迹达到祂的完美。}

因此，亲爱的诸位，藉着这不可言喻的健康恢复，我们没有留下任何骄傲或怠惰的余地：因为我们没有一样不是领受的，并且我们被明确警告，不可疏忽对待天主恩宠的恩赐。因为那如此及时来援助我们的，公正地用诫命督促我们；那引领我们走向荣耀的，慈悲地激励我们顺服。因此，主祂自己正当地成为我们的道路，因为除非通过基督，没有人能到基督那里去。但通过祂并朝向祂行走的，是那踏上祂忍耐与贬抑之道路的人，在那条路上，你确知不乏劳苦的热浪、悲伤的云彩、恐惧的风暴。恶人的罗网、不信者的迫害、强权者的威胁、骄傲者的侮辱，都在那里；而万军之主、荣耀之王，以我们软弱的形态和有罪肉体的模样，经历了这一切，为的是在此生危险之中，我们渴望的不再是逃避和躲开它们，而是忍受并战胜它们。\par

% structure: p0015 source=h2 target=subsection reason=relative-heading-rank; base=h2

\subsection{七、基督在十字架上呼喊\Quote{被遗弃}，是为了教导我们，人性若无神性是不够的。}

因此，主耶稣基督，我们的元首，在祂自身内代表祂身体的所有肢体，并在十字架的惩罚中替祂所救赎的人说话，发出了祂曾在圣咏中发出的呼喊：\BibleQuote{天主，我的天主，求祢眷顾我：祢为何舍弃了我？} 那呼喊，亲爱的诸位，是一种教导，而非怨言。因为在基督内，天主与人只有一个位格，祂不能被那与祂不可分离者所舍弃；祂是替我们这些战栗软弱的人发问：为什么那害怕受苦的肉身没有被垂听？因为在受难开始时，为了医治和纠正我们软弱的恐惧，祂曾说：\BibleQuote{父啊！若是可能，就让这杯离开我；但不要照我，而要照祢所愿的；} 又说：\BibleQuote{父啊！若是这杯不能离开我，除非我喝它，愿祢的旨意成就。} \BibleRef{玛 26:39} 因此，既然祂已征服了肉身的战栗，并已接受了父的旨意，将一切对死亡的恐惧践踏在脚下，那时正在施行祂计划的工程，为何在祂对如此胜利的凯旋时刻，祂却寻求被舍弃——即不被垂听——的原因和理由呢？除非是为了表明：祂为人的恐惧而辩护时所怀的情感，与祂按照父的永恒法令为世界的和好而作出的决然选择，是大不相同的。因此，这\Quote{不被垂听}的呼喊，正是对伟大奥迹的阐明，因为如果我们的软弱在祂内得到了所求，救赎者的能力就不会赐予人类任何东西。亲爱的诸位，今天这些话足矣，免得我们因讲论太长而使你们疲倦：其余的事，我们推迟到星期三再讲。如果你们祈求，主必垂听你们，使我们能藉着那生活且统治万世万代者的恩惠，履行我们的诺言。阿们。\par

\SourceRecord{Translated by Charles Lett Feltoe.；Nicene and Post-Nicene Fathers, Second Series；Vol. 12.；Edited by Philip Schaff and Henry Wace.；Buffalo, NY: Christian Literature Publishing Co.,；1895.；<http://www.newadvent.org/fathers/360367.htm>.}

% structure: p0001 source=h1 target=section reason=page-title

\section{讲道 68}

\Emph{（论受难，第十七篇：于星期三宣讲。）}\par

% structure: p0003 source=h2 target=subsection reason=relative-heading-rank; base=h2

\subsection{I. 基督的天主性在受难中从未离弃祂}

最亲爱的诸位，我们现已应允要讲述的上次讲道，其论证已到我们谈论被钉上十字架的主向圣父发出的呼喊之处：我们告诫单纯而不加思索的听者，不要将\Quote{我的天主，等等}这话理解为耶稣被钉在十字架上时，圣父天主性的全能已离祂而去；因为天主性和人性在祂内如此完全结合，以致这结合不能因刑罚或死亡而被破坏。因为每一实体保有其特性，天主既未远离祂身体的苦难，也未因肉体而成为可受苦的，因为受苦者内的天主性并未实际受苦。因此，按照降生成人的圣言之本性，那在万有中间被造的，正是那藉着祂万有得以被造的。祂被恶人之手逮捕，正是那不受任何限制的。祂被钉以铁钉，正是那不受任何创伤影响的。最后，祂经历死亡，正是那从未停止永恒的，因此两件事实都藉着无疑的标记得以确立，即基督内屈尊的真实和威严的真实；因为天主的能力与人的软弱结合，目的是，天主在使我们的成为祂的同时，也使祂的成为我们的。因此，圣子并未与圣父分离，圣父也未与圣子分离；不变的天主性和不可分离的天主圣三不容许任何分裂。因为虽然降生成人的任务专属天主独生子，但圣父与圣子的分离，并不比肉体与圣言的分离更多。\par

% structure: p0005 source=h2 target=subsection reason=relative-heading-rank; base=h2

\subsection{II. 基督的死亡是出于祂的旨意，但在拯救他人时，祂却不能救自己}

因此，耶稣大声呼喊说：\Quote{祢为什么舍弃了我？}为要向众人宣告，祂不应该被救拔、被保护，而应被交在残忍之人手中，即成为世界的救主和众人的赎主，不是出于悲惨，而是出于慈悲；不是由于援助的缺乏，而是由于决心赴死。但我们必须感到，祂的生命代祷的能力何其大，祂凭自己的内在能力死而复活。因为蒙福的宗徒说：圣父\Quote{没有怜惜自己的圣子，反而为我们众人把祂交出来} \BibleRef{罗 8:32}；他又说：\Quote{基督爱了教会，并为她舍弃了自己，为圣化她。}因此，主被交付于受难，既是圣父的旨意，也是祂自己的旨意，所以圣父不仅\Quote{舍弃}了祂，而且在某种意义上祂也舍弃了自己，不是仓促逃避，而是自愿退隐。因为被钉者的能力向那些恶人自我抑制，为了利用一个隐秘的计划，祂拒绝利用自己公开的能力。因为如果祂抵抗迫害者，那么祂来藉受难毁灭死亡和死亡之主的意图如何能拯救罪人呢？这便是犹太人的信念：耶稣被天主舍弃了，所以他们才能对祂行这样不圣洁的残忍行为；因为他们不明白祂奇妙忍耐的奥秘，便以亵渎的嘲笑说：\Quote{祂救了别人，却不能救自己。如果祂是以色列的君王，现在就从十字架上下来，我们就相信祂} \BibleRef{玛 27:42}。愚蠢的经师和邪恶的司祭啊，救主的能力并不按你们盲目的意愿而显示，也不因服从亵渎者的恶舌而延迟人类的救赎；因为如果你们愿意承认天主圣子的天主性，你们就会观察祂无数的作为，而这些作为本应坚定你们那虚假承诺的信德。但是，如果正如你们自己所承认的那样，祂确实救了别人，那么许多在众目睽睽之下行的大奇迹，为什么丝毫没有软化你们心灵的刚硬呢？除非是因为你们一直如此抗拒圣神，以致将天主的一切恩惠都转为你们的毁灭？因为即使基督从十字架上下来，你们仍将留在你们的罪恶中。\par

% structure: p0007 source=h2 target=subsection reason=relative-heading-rank; base=h2

\subsection{III. 当时正从旧约过渡到新约时期}

因此，空虚狂欢的侮辱被鄙视了，主在恢复迷失和堕落者时的慈悲并未因辱骂或斥责而偏离其目的之路。因为一只无与伦比的祭品正为世界的救恩献给天主，基督真羔羊的宰杀，通过如此多世代所预言的，正在将应许之子民转移至信仰的自由中。新约也在被批准，在基督的血中，永恒国度的继承者正在被登记；大司祭进入至圣所，无玷的司祭穿过祂肉体的帐幔，以在天主前代祷。最后，如此明显的过渡正从法律到福音、从会堂到教会、从许多祭献到唯一祭品进行着，以致当主交付灵魂时，那悬挂在前方、隔绝圣殿内部及其神圣至圣所的奥秘帐幔，被突然的力量从上到下撕裂了，因为真理正在取代预像，并且在所宣布的祂面前，前驱已不再需要。此外，所有元素都陷入可怕的混乱，大自然本身从基督的钉死者那里收回了支持。尽管管理钉十字架的百夫长因所见而恐惧，说\Quote{这人真是天主子}，但犹太人更比所有坟墓和岩石刚硬的邪恶心灵，并未被报告有任何痛悔刺透：因此看来罗马士兵当时比以色列的司祭更易承认天主子。\par

% structure: p0009 source=h2 target=subsection reason=relative-heading-rank; base=h2

\subsection{IV. 让我们在这一神圣季节中藉着斋戒和善工获益}

因为，犹太人被剥夺了这些奥迹所赋予的一切圣化，将他们的光明变为黑暗，将他们的\Quote{庆节变为悲哀}，所以，亲爱的诸位，让我们俯伏我们的身体和灵魂，朝拜天主的恩宠，这恩宠已倾注在万民身上，日复一日地恳求仁慈的圣父和富有的救赎者赐予我们帮助，使我们能逃脱今生的一切危险。因为狡猾的诱惑者无处不在，没有什么能逃脱他的网罗。我们必须以坚定的信德永远抵抗他，有天主眷顾的帮助，这帮助在一切危险中向我们伸出，这样，虽然他从不停止攻击，但他绝不能成功突破。亲爱的诸位，愿所有人都虔诚地遵守并善用斋戒，不要让任何过度损害这种克己的益处，我们已证明这种克己对灵魂和身体都是有益的。因为属于清醒和节制的事情，在此季节必须更加勤谨地遵守，以便从短暂的熱忱中养成持久的习惯；无论是在慈悲的工作中，还是在严格的克己中，信友们不可让任何时刻闲置，因为随着年岁增长、时间流逝，我们应当增加善行的储存，而不是浪费机会。对于虔诚的意愿和热心的灵魂，天主的仁慈必会赐予，祂使我们能够获得祂使我们渴望的一切，祂与我们的主耶稣基督、祂的圣子，以及圣神，永生永王。阿们。\par

\SourceRecord{Translated by Charles Lett Feltoe.；Nicene and Post-Nicene Fathers, Second Series；Vol. 12.；Edited by Philip Schaff and Henry Wace.；Buffalo, NY: Christian Literature Publishing Co.,；1895.；<http://www.newadvent.org/fathers/360368.htm>.}

% structure: p0001 source=h1 target=section reason=page-title

\section{讲道集 第七十一篇}

\Emph{（论主的复活，之一；于复活节前夕的圣周六宣讲。）}\par

% structure: p0003 source=h2 target=subsection reason=relative-heading-rank; base=h2

\subsection{一、我们都必须分享基督复活的生命}

在上次讲道中，亲爱的诸位，我认为恰当地向你们解释了我们在基督十字架上的参与，由此信徒的生活本身就包含了逾越节的奥迹，这样在庆节中所敬礼的，便由我们的实践来庆祝。你们自己已证明这多么有益，并通过你们的虔敬学会了，更长的斋戒、更频繁的祈祷和更慷慨的施舍对灵魂和身体有多大益处。因为几乎没有人没有从这操练中获益，也没有人在良心的深处没有储存一些他可以合理欢喜的东西。但这些益处必须以持久的谨慎来保持，免得我们的努力落入懈怠，魔鬼的恶意偷走天主恩宠所赐予的。因此，既然我们通过四十天的遵守希望产生这样的效果，即我们在主受难的时候应该感受到十字架的某些东西，我们必须努力也被发现分享基督的复活，并\Quote{出死入生}\BibleRef{若一 3:14}，当我们还在这身体内时。因为当一个人通过某种过程从一件事物变成另一件事物时，不再是原来的他对他是终结，而成为原来不是的是开始。但问题是，人是对什么死去或活着：因为有一种死亡，是生命的原因，有一种生命，是死亡的原因。除了在这个短暂的世界上，两者都被追求，以至于我们暂时的行为的性质决定了永恒报应的区别。因此，我们必须向魔鬼死去，向天主活着；我们必须向不义灭亡，以便向正义复活。让旧的沉没，使新的兴起；并且，正如真理所说，\Quote{没有人能事奉两个主人}\BibleRef{玛 6:24}，不要让那使站立者跌倒的成为主，而要让那使跌倒者兴起得胜的成为主。\par

% structure: p0005 source=h2 target=subsection reason=relative-heading-rank; base=h2

\subsection{二、天主没有将祂的灵魂撇在阴府，也没有让祂的肉身见腐朽}

因此，既然宗徒说：\Quote{第一个人出于地，属于土；第二个人出于天，属于天。那属于土的怎样，凡属于土的也怎样；那属于天的怎样，凡属于天的也怎样。我们既带有那属于土的肖像，也该让我们带有那属于天者的肖像}，我们应当大大欢喜这改变，藉此我们通过祂不可言喻的仁慈，从卑下的尘世被提升到天上的尊荣，祂降卑到我们的处境，为提拔我们到祂的处境，不仅取了有罪本性的实质，也取了它的状况，让不可受苦难的天主性受到有死之人所遭遇的一切痛苦的冲击。因此，为使门徒们不安的心不被长久的悲伤折磨，祂以如此奇妙的迅速缩短了祂所宣布的三天延迟，将第一天的末尾和第三天的开头与第二天整天连在一起，从而截去了一部分时间，却没有减少天数。因此，救主的复活并没有长久地将祂的灵魂留在\OriginalTerm{阴府}{Hades}，也没有将祂的肉身留在坟墓里；祂无朽坏肉身的复活如此迅速，与其说像死亡，不如说像睡眠，因为天主性并未离开祂所摄取的人性之任何一部分，而是以其大能重新结合了其大能所分离的。\par

% structure: p0007 source=h2 target=subsection reason=relative-heading-rank; base=h2

\subsection{三、基督复活后的显现显示祂的\OriginalTerm{位格}{Person}本质上与先前相同}

随后有许多证据，使要传遍普世的信仰权威得以建立。虽然石头的滚开、空坟墓、裹头巾的排列以及亲自叙述全部事迹的天使，完全确立了主复活的真理，但祂仍多次清楚地向妇女们和宗徒们的眼睛显现，不仅与他们谈话，还留下来与他们一同吃饭，并允许那些被疑虑困扰的人热心好奇的手触摸祂。为此，祂在门徒们所在之处，门都关着时进入，并向他们嘘气赐予圣神，在赐予他们领悟之光后，开启了圣经的奥秘，又亲自向他们显示肋旁的伤口、钉痕以及祂最近受难的所有痕迹，由此可知在祂内，神性和人性的特性不可分割地存留，使我们如此认识圣言不是肉身，同时承认天主的独生子既是\OriginalTerm{圣言}{Word}又是肉身。\par

% structure: p0009 source=h2 target=subsection reason=relative-heading-rank; base=h2

\subsection{四、然而，虽是同一肉身，却也是光荣了的}

外邦人的宗徒保禄，亲爱的诸位，与这信念并不冲突，当他说：\Quote{即使我们曾按肉身认识基督，但如今不再这样认识祂了。}因为主的复活不是肉身的结束，而是改变，祂的实质并未因能力的增加而消灭。性质改变了，但本性并未停止存在：身体成为不可受苦的，而它曾是可能被钉的；它成为不朽坏的，而它曾是可能被伤害的。基督的肉身被称为不再以先前所知的状态被认识，这是恰当的，因为其中不再有任何可受苦的、软弱的，所以它在本质上是相同的，在光荣上却不相同。但若圣保禄关于基督的身体如此主张，有何奇怪呢？当他说到所有属灵的基督徒时：\Quote{因此，从今以后，我们不再按肉身认识谁了。}他说，从今以后，我们开始在基督内经历复活，自从在祂——为众人死的那位——内，我们所有的希望都得到了保证。我们不犹豫不决，不处于不确定的悬念中，而是既已领受了应许的凭证，如今以信德的眼睛看见将要发生的事，并因我们本性的提升而喜乐，我们已经拥有我们所相信的。\par

% structure: p0011 source=h2 target=subsection reason=relative-heading-rank; base=h2

\subsection{五、因望德而得救，我们不可满足肉身的私欲}

因此，我们不要被暂时事物的表象所迷惑，也不要让我们的默想从天上事物转向地上事物。那些大部分尚未发生的事，应视为已经完成；专注于永恒的心灵，应将其渴望固定在那里，因为那里所提供的是永恒的。因为尽管\BibleQuote{我们因望德而得救}〔见：\BibleRef{罗 8:24}〕，并且仍然带着可朽坏可死的肉身，但如果肉欲不支配我们，我们就正确地被称为不在肉体内，并且有理由不再以我们所不随从其意志的事物来称呼。因此，当宗徒说：\BibleQuote{不要放纵肉身的私欲}〔见：\BibleRef{罗 13:14}〕时，我们明白，那些有益于健康、人性软弱所需的事并未被禁止；但因为我们不能满足所有欲望、也不能纵容肉身所贪求的一切，我们认识到，被警告要如此节制：既不要过度，也不要拒绝控制下的肉身所必需的事。因此，同一宗徒在另一处说：\BibleQuote{从来没有人恨过自己的肉身，反而养育它、保护它}〔见：\BibleRef{厄 5:29}〕；当然，就养育和保护而言，不是耽于恶习和奢侈，而是为了其正当的功能，使本性得以恢复并保持应有的秩序：低下的部分不错误地、卑劣地压倒较高的部分，较高的也不屈服于较低的，以免若恶习压倒心灵，本应统治的地方反成为奴役。\par

% structure: p0013 source=h2 target=subsection reason=relative-heading-rank; base=h2

\subsection{六、我们虔诚的决心必须全年持续，而不只限于复活节}

愿天主的子民因此认识到，他们在基督内是新的受造物，并要警醒地理解：他们是被谁所收纳，又收纳了谁。不要让已经更新的事物退回到旧日的不稳定中；不要让那\BibleQuote{手扶着犁}〔见：\BibleRef{路 9:62}〕的人放弃他的工作，而要关注他所播种的，不要回顾他所舍弃的。任何人不要跌回他已从中起来的境地，即使他因身体的软弱仍在某些疾病中受苦，也要迫切渴望被医治和提升。因为这就是健康之路，借着效法在基督内开始的复活；尽管在这易滑的路上，旅人可能遭遇许多意外和跌倒，但他的脚能从泥沼被引导到坚固之地，因为经上记载：\BibleQuote{人的脚步由主引导，祂喜爱他的道路。义人虽跌倒，却不至被弃，因为主要伸手扶持他。}亲爱的诸位，这些思想不仅应留意为复活节庆典，也应为了整个生命的圣化；我们当前的操练应以此为方向，使短暂的体验所带给信徒心灵的喜乐，能成为习惯而持久不变；若有任何过失潜入，当以迅速的补赎消除。因为久病的治疗缓慢而困难，应在伤口尚新时及早施以补救，好使我们从一切跌倒中不断兴起，堪当获得我们荣耀化肉身的不朽复活，在基督耶稣我们的主内——祂与圣父及圣神，永生永王。阿们。\par

\SourceRecord{Translated by Charles Lett Feltoe.；Nicene and Post-Nicene Fathers, Second Series；Vol. 12.；Edited by Philip Schaff and Henry Wace.；Buffalo, NY: Christian Literature Publishing Co.,；1895.；<http://www.newadvent.org/fathers/360371.htm>.}

% structure: p0001 source=h1 target=section reason=page-title

\section{讲道集 72}

\Emph{（论主的复活，二）}\par

% structure: p0003 source=h2 target=subsection reason=relative-heading-rank; base=h2

\subsection{一、十字架不仅是救恩的奥秘，也是跟随的榜样}

逾越奥迹的全部，亲爱的诸位，已在福音叙述中呈现在我们面前，而心灵的耳朵通过肉体的耳朵被触及，以致你们没有人不能对事件有一幅画面：因为神圣默感的故事文本已清楚地显示了主耶稣基督被出卖的背信弃义、定祂罪的审判、钉死祂的残暴，以及祂复活的荣耀。但一篇讲道仍需我们提供，使司铎的劝勉得以添加到圣经的庄严诵读中，正如我相信你们正以虔诚的期待向我们索取你们惯常的职责。因此，因为在信友的耳朵中没有无知的位置，圣言的种子，即福音宣讲，应当在你们心田的土壤中生长，以致当令人窒息的荆棘和蒺藜被去除后，圣善思想的植物和正当愿望的芽苞可自由地结出果实。因为基督的十字架，为凡人的救恩而竖立，既是奥秘又是榜样：是圣事，藉此天主的能力生效；是榜样，藉此人的热诚被激发；因为对于从囚犯的轭下被救出的人，救赎还进一步赐予他们通过效仿跟随十字架道路的能力。因为如果世界的智慧在其错误中如此自傲，以致每个人都追随他所选择的领袖的意见、习惯和整个生活方式，那么除了与祂不可分离地结合，我们又将如何分享基督的名？祂就是祂自己所宣称的\Quote{道路、真理和生命？} \BibleRef{若 14:6}圣善生活的道路、神圣道理的真理和永恒福乐的生命。\par

% structure: p0005 source=h2 target=subsection reason=relative-heading-rank; base=h2

\subsection{二、基督为了我们的救恩取了我们的本性}

因为当整个人类在始祖中堕落时，仁慈的天主意欲如此救助祂按自己肖像所造的受造物，藉祂的独生子耶稣基督，使我们的本性的恢复不脱离它而实现，并且我们的新状态应是对原始地位的进步。有福的，如果我们没有从天主造我们的状态堕落；但更有福的，如果我们保持祂重新造我们的状态。从基督领受形式是大的；但在基督内拥有实体更大。因为我们被那本性接纳到它自身之中（那本性屈尊接受慈爱所规定的那些限制，且并未遭受任何变化。我们被那本性接纳），那本性并没有毁灭我们的本性中属于祂的东西，也没有毁灭祂的本性中属于我们的东西；它使神性与人性的位格在自身中如此合一，以致通过软弱与能力的协调，肉身不能因神性而变得不可侵犯，神性也不能因肉身而变得可受苦。我们被那本性接纳，它没有从我们族类的共同根上折断枝条，却排除了已传于所有人的罪的全部污染。也就是说，软弱和必死性，不是罪，而是罪的惩罚，被世界的救赎者以惩罚的方式承受，以便它们可被算作救赎的代价。因此，在我们所有人中是谴责的遗产的，在基督内是\Quote{敬虔的奥秘}。因为祂无债，祂将自己交予那最残酷的债主，并容忍犹太人的手成为魔鬼折磨祂无玷肉身的工具。祂愿意这肉身服从死亡，甚至直到祂（迅速的）复活，为达此目的：使信祂的人既不觉得迫害不可忍受，也不觉得死亡可怕，因记得他们分享祂的荣耀与祂分享他们的本性一样没有疑问。\par

% structure: p0007 source=h2 target=subsection reason=relative-heading-rank; base=h2

\subsection{三、复活并升天的主仍与我们同在}

因此，亲爱的诸位，如果我们毫不犹疑地用心相信我们用口所宣认的，那么在基督内，我们被钉死，我们死亡，我们被埋葬；也在第三日，我们复活。因此宗徒说：\Quote{你们既然与基督一同复活了，就该寻求天上的事，那里有基督，坐在天主的右边：你们该思念天上的事，不该思念地上的事。因为你们已经死了，你们的生命与基督一同藏在天主内。当基督，你们的生命，显现时，你们也要与祂一同在光荣中显现。} \BibleRef{哥 3:1} 但是为使信友的心知道他们拥有用以弃绝世界的欲望并被提升至天上智慧的凭借，主应许我们祂的临在，说：\Quote{看！我同你们天天在一起，直到今世的终结。} \BibleRef{玛 28:20} 因为圣神藉依撒意亚所说并非徒然：\Quote{看！一位贞女将怀孕生子，人要称祂的名字为厄玛奴耳，意译就是：天主与我们同在。} 因此，耶稣实现了祂名字的恰当含义，并在升天时不弃舍祂的义兄弟，虽然\Quote{祂坐在圣父的右边}，却仍居住在整个身体内，并从上面亲自坚强他们耐心等待，同时祂召唤他们上升至祂的荣耀。\par

% structure: p0009 source=h2 target=subsection reason=relative-heading-rank; base=h2

\subsection{四、我们必须怀有在基督耶稣内同样的心怀}

因此，我们不可在虚妄的追求中沉溺于愚昧，也不可在逆境中屈服于恐惧。一方面，我们无疑被欺骗所谄媚；另一方面，又被烦恼所重压。然而，因为\Quote{大地充满了主的仁慈}，基督的胜利确实属于我们，为要成就祂所说的：\Quote{不要害怕，因为我已战胜了世界}\BibleRef{若 16:33}。那么，无论我们是对抗世界的野心，还是肉体的情欲，或是异端的毒箭，让我们总是以主的十字架武装自己。因为如果我们从旧恶的酵母中（在真理的真诚中）戒绝，我们的逾越节将永不止息。因为在这充满各样苦难的生命的变迁中，我们应当记住宗徒的劝勉，他教导我们说：\Quote{你们该怀有基督耶稣所怀有的心情：祂虽具有天主的形体，却没有以自己与天主同等为应当把持不舍的，反而空虚了自己，取了奴仆的形体，成了人的模样；形态上完全像人一样。祂自谦自卑，听命至死，且死在十字架上。因此，天主极其举扬祂，赐给祂一个名字超越其它一切名字，使天上、地上和地下的一切，一听耶稣的名，无不屈膝叩拜；一切唇舌无不明认耶稣基督是主，以光荣天主圣父}\BibleRef{斐 2:5}。如果，他说，你们明白\Quote{伟大虔敬的奥秘}，并记住天主的独生子为拯救人类所行的一切，\Quote{你们该怀有基督耶稣所怀有的心情}，祂的谦卑不应被任何富人藐视，也不应被任何出身高贵者视为羞耻。因为没有任何人类的幸福能达到如此高度，竟将自身引以为耻：天主，虽处于天主的形体中，却不认为取奴仆的形体对祂有失身份。\par

% structure: p0011 source=h2 target=subsection reason=relative-heading-rank; base=h2

\subsection{五、惟有持守降生成人真理者，方能妥善庆祝逾越节}

效法祂所行的事，爱祂所爱的事物；你们既在自身中寻得天主的恩宠，也当在祂内爱慕你们的本性作为回应。因为祂在贫贱中未失去富有，在谦卑中未减损荣耀，在死亡中未丧失永恒；同样，你们也要步武祂的足迹，轻看尘世之事，以获享天上的事物。因为背起十字架意味着杀死情欲，灭绝恶习，远离虚荣，弃绝一切错谬。因为，虽然主的逾越节任何不洁、放荡、骄傲或悭吝之人都无法庆祝，然而最远离此庆节的却是异端者，尤其是那些对圣言的降生成人持有错误观点的人：他们或轻视属于天主性的一切，或将属于肉身的视为不真实。因为天主子是真天主，从圣父那里拥有圣父的一切，没有时间上的开始，不受任何变化的影响，与独一天主不可分割，与全能者无异，是永恒圣父的永恒独生子；如此，那忠信的理智，在同一个独一天主性的本质中相信圣父、圣子和圣神，既不因引入尊位的等级而分割一体，也不因混淆位格而混乱天主圣三。但是，仅仅认识天主子在圣父的本性中是不够的，除非我们在祂自己的本有属性未被收回的情况下，也在我们的人性中承认祂。因为祂为了人的恢复而承受的自我空虚，是慈悲的安排，而非能力的丧失。因为，根据天主永恒的计划，\Quote{除祂以外，天下没有赐下别的名字，使我们赖以得救}\BibleRef{宗 4:12}，那不可见者使祂的本体成为可见，那永恒者成为暂时的，那不能受苦者成为可受苦的：不是为使能力陷于软弱，而是为使软弱转化为不可摧毁的能力。\par

% structure: p0013 source=h2 target=subsection reason=relative-heading-rank; base=h2

\subsection{六、论\Quote{逾越}一词之奥秘应用}

因此，那被我们称为\Quote{巴斯卦}（\OriginalTerm{巴斯卦}{Pascha}）的庆节，在希伯来人中被称为\Quote{逾越}（\OriginalTerm{逾越}{Phase}），即\Quote{逾越节}，正如福音作者所证实的说：\Quote{在逾越节庆日前，耶稣知道祂离此世归父的时辰已到。}然而，祂这样离开时所凭借的本性若不是我们的，又是什么呢？因为圣父在圣子内，圣子也在圣父内，是不可分离的。但是，因为圣言与肉身是一个位格，被取的人性不与取它的神性分离；而受提拔的尊荣被认为是归于那提拔者，正如宗徒在我们已引用过的段落中说：\Quote{因此，天主极其举扬祂，赐给祂一个名字超越其它一切名字。}这里无疑说的是祂所取的人性被举扬，因此，祂在受难时天主性保持不可分，同样也在天主性的荣耀中与圣父同永恒。为分享这不可言传的恩赐，主亲自为祂的忠实者预备了一个有福的\Quote{逾越}，当祂在受难的门槛上时，不仅为祂的宗徒和门徒，也为整个教会祈求说：\Quote{我不但为他们祈求，也为那些因他们的话而信从我的人祈求，使他们都合而为一，就如祢、父，在我内，我在祢内，使他们也在我们内合而为一}\BibleRef{若 17:20}。\par

% structure: p0015 source=h2 target=subsection reason=relative-heading-rank; base=h2

\subsection{七、惟有真信徒方能庆祝逾越节}

在此合一之中，那些否认在真天主的天主圣子内人的本性常存的人，无法有分；他们攻击这救恩的奥迹，将自己排除于复活节之外。因为他们既背离福音，又否认信经，便不能与我们一同庆祝；尽管他们胆敢自取基督徒的名号，却为每个以基督为元首的受造物所摒弃。你们在这神圣的时节欢欣踊跃，虔敬喜乐，是理所当然的，因为你们不容许任何虚假混入真理，对基督按肉体的诞生、祂的受难与死亡、以及祂身体的复活毫无怀疑；因为你们承认的基督，毫无天主性的分离，祂真是由童贞玛利亚的子宫诞生，真是被悬于十字架的木材上，真是被安放在世间的坟墓里，真正在光荣中复活，真正被安置在圣父尊威的右边。\Quote{由此，}如宗徒所说，\BibleQuote{我们也等候救主，主耶稣基督。祂将改变我们卑微的身体，使之相似祂光荣的身体。}\BibleRef{斐 3:20}。祂永生永王，等等。\par

\SourceRecord{Translated by Charles Lett Feltoe.；Nicene and Post-Nicene Fathers, Second Series；Vol. 12.；Edited by Philip Schaff and Henry Wace.；Buffalo, NY: Christian Literature Publishing Co.,；1895.；<http://www.newadvent.org/fathers/360372.htm>.}

% structure: p0001 source=h1 target=section reason=page-title

\section{讲道第七十三篇}

\Emph{（论主的升天，第一讲）}\par

% structure: p0003 source=h2 target=subsection reason=relative-heading-rank; base=h2

\subsection{一、复活后所记录的事件旨在使我们确信其真实性}

自我们的主耶稣基督蒙福而光荣的复活以来，——藉此，天主的能力在三天内重建了被犹太人的邪恶所推翻的天主的真正圣殿——至圣的四十天，亲爱的诸位，今日结束了；这些日子藉至圣的安排，用于我们最有益的教导，好使在此期间，主延长祂身体临在的逗留，我们的信仰得以藉必需的证据而坚固。因为基督的死曾使门徒们的心大为不安，一种不信的麻木因祂在十字架上的酷刑、祂的断气、祂无生命的身体被埋葬，而潜入他们悲伤的心灵。因为当圣妇们按福音所传，报告石头从坟墓滚开、坟墓中不见尸体、天使为生活的主作证时，宗徒和其他门徒觉得她们的话像是胡言乱语。这种由人性软弱产生的疑惑，真理之神最肯定不会容许存在于祂自己宣讲者心中，除非他们颤抖的焦虑和谨慎的犹豫为我们的信德奠定了基础。是为我们的困惑和危险，在宗徒们身上作了预备；是我们在这些人身上被教导如何面对不敬者的挑剔和世俗智慧的辩驳。\Emph{我们}藉他们的观看而受教，\Emph{我们}藉他们的听闻而学习，\Emph{我们}藉他们的触摸而确信。让我们感谢上智的安排和圣父。必要的缓慢相信。其他人怀疑，好使我们不再怀疑。\par

% structure: p0005 source=h2 target=subsection reason=relative-heading-rank; base=h2

\subsection{二、因此它们具有最高的教导意义}

因此，亲爱的诸位，那介于主的复活与升天之间的日子，并非在无所事事中度过，而是在其中确认了伟大的奥迹，显明了深奥的真理。在此期间，对死亡的可怕恐惧被移除，不仅灵魂的不朽，而且肉身的不朽也被确立。在此期间，主向门徒们嘘气，圣神倾注于所有宗徒，并且除了天国的钥匙之外，主的羊群被委托给蒙福的宗徒伯多禄。那时，主在路上与两位门徒同行，为了扫除我们一切不确定的乌云，祂责备他们胆怯的心之迟钝。他们被光照的心燃起信德的火焰，当他们还是冷淡时，在主解释圣经时变得火热。在擘饼时，他们与主同席时眼睛被打开：他们的眼睛被打开是多么有福，因为他们本性的荣耀被启示，远胜于我们的原祖父母，在他们身上落下了过犯的灾难性后果。\par

% structure: p0007 source=h2 target=subsection reason=relative-heading-rank; base=h2

\subsection{三、它们证实了肉身的复活}

在这些奇迹和其他奇迹的过程中，当门徒们被困惑的思绪困扰，主出现在他们中间并说：\Quote{愿你们平安}。祂为了不让他们心中所想的成为他们的固定意见（因为他们以为看见的是神体，而非血肉），祂驳斥他们与真理如此不一致的想法，向怀疑者提供仍在祂手脚上的十字架印记，邀请他们仔细触摸查验，因为钉痕和枪痕被保留下来，为医治不信之心的创伤，好使他们不是以动摇的信德，而是以最坚定的知识领悟：那曾躺在坟墓中的本性，将坐在天主圣父的宝座上。\par

% structure: p0009 source=h2 target=subsection reason=relative-heading-rank; base=h2

\subsection{四、基督的升天赐予我们比魔鬼夺走的更大的特权和喜乐}

因此，亲爱的诸位，在这介于主的复活与升天之间的整个时期，天主的眷顾有此目的：教导并印刻在祂子民的眼目和心中，使主耶稣基督被承认真实复活了，正如祂真实地诞生、受难、死亡一样。因此，最蒙福的宗徒和所有门徒，他们曾在祂十字架上的死亡中困惑，在相信祂复活上迟钝，如今被真理的清晰所坚强，当主升入高天时，他们不仅没有悲伤，反而充满极大的喜乐。他们的喜乐确实是伟大而不可言喻的，因为在圣众眼前，人性超越了所有天上受造物的尊威，上升，超越天使的等级，超过总领天使的高度，其上升不受任何高度的限制，直到被接去与永恒之父同坐，与祂的荣耀共享宝座，这荣耀与祂在圣子内所合一的本性同在。既然基督的升天是我们的提升，而身体的希望也上升，到了头的光荣所先到之处，让我们以相称的喜乐和忠诚的感谢而欢跃，亲爱的诸位。因为今天不仅我们被确认为乐园的拥有者，而且藉着基督已穿透诸天的高处，并藉着基督不可言喻的恩宠获得了比我们因魔鬼的恶意所失去的更伟大的事物。对于我们，我们恶毒的仇敌曾将我们从原初居所的福乐中驱逐出来，天主之子却使我们成为祂自己的肢体，并安置在圣父的右边，祂与圣父和圣神同享生命，并为天主，至于无穷之世。阿们。\par

\SourceRecord{Translated by Charles Lett Feltoe.；Nicene and Post-Nicene Fathers, Second Series；Vol. 12.；Edited by Philip Schaff and Henry Wace.；Buffalo, NY: Christian Literature Publishing Co.,；1895.；<http://www.newadvent.org/fathers/360373.htm>.}

% structure: p0001 source=h1 target=section reason=page-title

\section{第七十四篇讲道}

\Emph{（论主的升天，第二篇）}\par

% structure: p0003 source=h2 target=subsection reason=relative-heading-rank; base=h2

\subsection{一、升天完成我们对他——既是天主也是人——的信德}

我们救恩的奥秘，亲爱的诸位，宇宙的创造主以自己宝血为代价所珍视的，如今已在屈辱的条件下，从他肉身诞生之日起直到他受难的终结，得以完成。并且，即使在\Quote{奴隶的形态}中，许多天主性的标记已经闪耀出来，然而那整个时期的事件特别显示了他所取人性的真实性。但在受难之后，当死亡的锁链被打破时（死亡曾因攻击那不知罪的祂而暴露了自己的力量），软弱变成了能力，必死变成了永生，侮辱变成了光荣——主耶稣基督以许多清晰的证据在众人眼前显明了这一切，直到他将战胜死亡所获得的胜利凯旋地甚至带入天上。因此，正如在逾越节纪念中，主的复活是我们喜乐的原因；同样，我们现今欢庆的主题是他的升天，因为我们纪念并恰当地敬礼那一天，我们的卑微本性在基督内被高举于天上万军之上，超过所有天使的品级，超越一切掌权者的高度，与天主圣父同坐。我们正是建立并建造在这上智安排的事件次序上，为的是天主的恩宠变得更加奇妙，因为尽管从人的视线中移除了那理应令人敬畏的事物，信德并未失败，望德并未动摇，爱德并未冷却。因为这是伟大心智的力量和坚定信德灵魂的光明：毫不迟疑地相信肉眼所未见的事物，并将情感寄托于你看不见的地方。而这份虔诚从何在我们心中产生，或人如何因信德而成义，倘若我们的救恩仅依赖于那些在我们眼目之下的事物？因此，我们的主对那似乎怀疑基督复活，直到他用视觉和触觉在他自己的肉身中检验了受难的痕迹的人说：\BibleQuote{因为你看见了我，才相信；那些没有看见而相信的，才是有福的。\BibleRef{若 20:29}}\par

% structure: p0005 source=h2 target=subsection reason=relative-heading-rank; base=h2

\subsection{二、升天使我们的信德更卓越、更坚强}

因此，亲爱的诸位，为使我们有分于此福乐，当所有关于福音宣讲和新约奥迹的事都已成就时，我们的主耶稣基督在复活后第四十天，在门徒面前被提升天，终止了他与我们在肉身上的临在，以坐在圣父的右边，直到为增加教会子女而神圣预定的时期完成，那时他要在上升时所同有的肉身中，来审判生者死者。因此，那时我们救赎主可见的临在，转变为一种圣事性的临在；并且为使信德更卓越、更坚强，视觉让位于教义，这教义的权威要由被从上而来光芒照亮的信德心灵所接受。\par

% structure: p0007 source=h2 target=subsection reason=relative-heading-rank; base=h2

\subsection{三、这信德在众人身上的奇妙效果}

这信德，因主的升天而增长，又因圣神的恩赐而坚固，没有被锁链、监禁、流放、饥饿、火刑、野兽攻击、残酷迫害者的精巧折磨所吓倒。因为这种信德，普世的人不仅男子，甚至妇女，不仅未成年的男孩，甚至柔弱的少女，都奋战至流血。这信德驱逐邪灵、驱除疾病、复活死者；藉着它，诸位有福的宗徒自己，虽曾被如此多的奇迹所证实、被如此多的训诲所教导，却在主受难的恐怖中惊慌失措，且未毫不犹豫地接受他复活的真理，但在主的升天之后取得了如此进步，以至于先前所有使他们充满恐惧的事都变成了喜乐。因为他们已将整个心灵的默观高举到那坐在圣父右边者的天主性上，不再受肉体视觉的障碍所阻挠，而将其心灵的目光投向那从未离开圣父的那一位，祂下降尘世时从未离开圣父，上升天庭时亦未舍弃门徒。\par

% structure: p0009 source=h2 target=subsection reason=relative-heading-rank; base=h2

\subsection{四、他的升天锤炼我们的信德：天使对他的事奉显示了他权柄的广大}

因此，亲爱的诸位，人子与天主之子在回到圣父尊威的光荣中时，获得了更卓越、更神圣的名声，并且以不可言喻的方式，就其天主性而言，开始更亲近圣父，而就其人性而言，则变得更远。受过更好教育的信德开始更接近对圣子与圣父平等的认识，无需触摸基督内可朽的身体——基督借此而小于圣父，因为，当荣耀的身体的本性仍然存留时，信徒的信德被召叫不是用血肉的手，而是用属灵的悟性去触摸那与圣父平等的独生子。因此，主在复活后，当代表教会的玛利亚玛达肋纳急忙前来触摸祂时，祂说：\BibleQuote{你别碰我，因为我还没有升到我的父那里去\BibleRef{若 20:17}}：意思是，我不愿意你如同来到一个身体前那样接近我，也不愿你凭肉身的知觉认识我：我将你延后，为的是更高的事，我为你预备更伟大的事：当我升到我的父那里时，那时你将更完美、更真实地触摸我，因为你将把握你不能触摸的，相信你不能看见的。但是，当门徒们以热切惊异的目光仰望着升天的主时，有两个身着耀眼光芒衣裳的天使站在他们旁边，他们说：\BibleQuote{加里肋亚人，你们为什么站着望天呢？这位从你们中被接到天上去的耶稣，你们看见祂怎样升天，祂也要怎样降来\BibleRef{宗 1:11}。} 借着这些话，教会所有的子女都被教导要相信耶稣基督将来会以祂升天时所具有的同一肉身可见地来临，并且不要怀疑万事都服从于那从祂诞生之初就有天使服侍的主。因为，正如天使向蒙福的童贞女报告基督将由圣神受孕一样，天上的众生也向牧人们歌唱祂由童贞女诞生。正如来自高天的使者首先证实了祂从死者中复活，同样，天使的服务也被用来预告祂以真实肉身审判世界的来临，好让我们明白当审判时将有怎样伟大的权能伴随祂而来，因为这样伟大的使者甚至在祂受审时服侍了祂。\par

% structure: p0011 source=h2 target=subsection reason=relative-heading-rank; base=h2

\subsection{V. 我们必须轻视地上的事物，提升到上面的事物，尤其是通过积极的慈悲和爱德工作}

所以，亲爱的诸位，让我们以属灵的喜乐欢欣，让我们愉快地向上主献上相称的感恩，并畅通无阻地将我们内心的眼睛举向那些基督所在的高处。蒙召被提升的心灵不可被地上的情感所压制，预定归于永恒的事物不可专注于那朽坏的事物，已走上真理之路的人不可陷入诡诈的罗网，信友必须如此经过这些短暂的事物，以致他们记得自己是在这世界的幽谷中作客旅，即使遇到一些吸引人的东西，也不可罪恶地拥抱它们，而要勇敢地经过它们。因为蒙福的宗徒伯多禄激励我们行此虔诚，并以那因三次宣认爱主而生的牧养基督羊群的热切恳求对我们说：\BibleQuote{亲爱的，我恳求你们，作为客旅和寄居的，要戒绝那些与灵魂作战的肉欲\BibleRef{伯前 2:11}。} 但肉体的欢愉为谁作战呢，如果不是为魔鬼，它的快乐就是用可朽坏之善物的诱惑束缚那些追求上面事物的灵魂，并把它们从它自己已被驱逐的那些居所中拉开？每一个信友都必须警惕它的阴谋，以便从攻击的侧面粉碎它的敌人。亲爱的诸位，没有比仁慈的怜悯和慷慨的爱德更强大的武器来对抗魔鬼的诡计，借着它们，每一项罪要么被逃避，要么被征服。但这崇高的能力只有在对立之物被推倒后才能达到。有什么比贪婪更敌对怜悯和爱德的工作呢？而贪婪是一切的根源。除非它因缺乏滋养而被摧毁，否则在那颗心中——这邪恶的野草已扎根——必然会长出荆棘和蒺藜般的恶习，而不是任何真正良善的种子。因此，亲爱的诸位，让我们抵抗这致祸的邪恶，并\Quote{追求爱德}，没有它任何德行都不能兴盛，好使我们借着这条爱德之路——基督曾顺着它降到我们中间——我们也能上升到祂那里，与天主圣父和圣神一起，享受尊荣和光荣，至于无穷之世。阿们。\par

\SourceRecord{Translated by Charles Lett Feltoe.；Nicene and Post-Nicene Fathers, Second Series；Vol. 12.；Edited by Philip Schaff and Henry Wace.；Buffalo, NY: Christian Literature Publishing Co.,；1895.；<http://www.newadvent.org/fathers/360374.htm>.}

% structure: p0001 source=h1 target=section reason=page-title

\section{第七十五篇讲道}

\Emph{（圣神降临节，第1篇）}\par

% structure: p0003 source=h2 target=subsection reason=relative-heading-rank; base=h2

\subsection{I. 梅瑟颁布法律为圣神的倾注预备了道路}

亲爱的各位，所有天主教徒的心都认识到，今天的隆重节日应被视为主要庆节之一，并且毫无疑问，这一天应受到极大的尊敬，因为圣神以祂极卓越恩赐的奇迹祝圣了它。因为自从主升上所有天穹之上，坐在天主父的右边以来，这是闪耀的第十天，也是祂复活以来的第五十天，正是它开始的那一天，并且它本身包含新旧奥迹的伟大启示，由此最清楚地显示：恩宠通过法律被预告，而法律通过恩宠得以满全。因为正如古时，当希伯来民族从埃及人手中被释放时，在祭杀羔羊后的第五十天，法律在西乃山被颁布；同样，在基督受难之后，其中真正的天主的羔羊被宰杀，在祂复活后的第五十天，圣神降临在宗徒和信众身上，以致热心的基督徒可以轻易看出，旧约的开端为福音的开端作了预备，而第二个盟约是由同一圣神所建立，祂也曾设立了第一个盟约。\par

% structure: p0005 source=h2 target=subsection reason=relative-heading-rank; base=h2

\subsection{II. 那\Quote{不同方言}的恩赐是何等奇妙！}

因为正如宗徒的记载所证明：\BibleQuote{当五旬节的日子满了，众门徒都聚集在同一地方，忽然从天上来了一阵响声，好像暴风刮来，充满了他们所在的整座房屋。又有火舌似的分开显现，停在各人头上。他们都充满了圣神，就用其他语言说起话来，正如圣神赐予他们所说的一样。\BibleRef{宗 2:1}。}噢！智慧的话语何其迅速，当天主为师傅时，所教导的多么快就被学会。理解无需解释，使用无需练习，学习无需时间；但真理之神随意吹拂，各民族特有的语言在教会的口中成为共同的财富。因此，从那天起，福音宣讲的号角已响亮吹起；从那天起，恩赐的甘霖、祝福的河流，浇灌了每一片荒漠和所有干涸之地，因为为了更新地面，天主之神\BibleQuote{运行在水面上，}又为了驱散旧日的黑暗，新光的光芒闪耀，当那忙碌的舌头的火焰点燃了主光明的圣言和炽热的雄辩，其中既有启发悟性的能力，也有烧毁罪恶的力量。\par

% structure: p0007 source=h2 target=subsection reason=relative-heading-rank; base=h2

\subsection{III. 天主圣三的三个位格在一切事上完全平等}

但是，亲爱的，所行之事的实际形式虽然极其奇妙，并且无疑在所有人语言的欢庆合唱中，圣神的威严临在，但谁也不可认为，祂的神性本体出现在肉眼所见的形式中。因为祂的本性，是无形可见的，且与圣父和圣子共同所有，以祂所喜悦的外在标记显示了祂恩赐和工作的特性，但将祂本质的属性保留在祂自己的天主性内：因为人的视力无法看见圣神，正如无法看见圣父或圣子一样。因为在神圣的天主圣三中，没有什么是不同或不等的，凡是关于其本体所能想到的，在能力、光荣或永恒上都不承认有差异。并且在每个位格的特性上，圣父是一位，圣子是另一位，圣神是另一位，然而天主性并非分离和不同；因为圣子是圣父的独生子，圣神是圣父和圣子的圣神，并非像每个受造物是圣父和圣子的受造物那样，而是作为与两者同生活和同有能力，并且永恒地从那即是父和子者所出。因此，当主在祂受难之日以前，向祂的门徒应许圣神的来临时，祂说：\BibleQuote{我还有许多事要告诉你们，然而你们现在不能担负。但当祂，真理之神来到时，祂要引导你们进入一切真理。因为祂不凭自己说话，而是凡祂所听到的，祂要说话，并要把将来的事告诉你们。凡父所有的，都是我的；因此我说，祂要从我那里领受，并要告诉你们。\BibleRef{若 16:12}。}因此，并非有一些事物属于父，另一些属于子，另一些属于圣神：而是凡父所有的，子也有，圣神也有：并且从未有一个时候这种共融不存在，因为对于他们，拥有一切就是永远存在。在他们内，不要想象任何时间、等级或差异，并且，如果没有人能解释关于天主的真理，那么谁也不要胆敢断定不真实的事。因为不完全陈述关于祂不可言喻的本性，比制定一个实际上错误的定义更容易被原谅。因此，无论热诚的心能想象父的永恒不变的光荣，让他们同时理解子与圣神也是如此，没有任何分离或差异。因为我们承认这有福的天主圣三是独一的天主，原因就在于，在这三个位格中，无论是本体、能力、意志，还是行动，都没有差异。\par

% structure: p0009 source=h2 target=subsection reason=relative-heading-rank; base=h2

\subsection{IV. 马其顿异端与亚略异端同样亵渎}

因此，我们憎恶亚略派，他们主张圣父与圣子之间有差异；同样，我们也憎恶马其顿派，他们虽然承认圣父与圣子平等，却认为圣神是较低的本性，没有意识到他们因此陷入了那种无论在今世或是在来世审判中都无法得到赦免的亵渎之罪，如主所说：\Quote{凡说话干犯人子的，可得赦免；但说话干犯圣神的，无论是在今世或是在来世，都不得赦免 \BibleRef{玛 12:32}}。因此，坚持这种不虔敬乃是不可赦免的，因为它使人脱离了能使他认罪的那一位；那没有辩护者为他代求的人，永远得不到医治的宽恕。因为对圣父的呼求来自祂，忏悔的眼泪来自祂，恳求者的呻吟来自祂，并且\Quote{除非受圣神感动，也没有人能说\InnerQuote{耶稣是主} \BibleRef{格前 12:3}}，祂的同等权能和唯一神性，与圣父圣子相同，宗徒极其明确地宣告，说：\Quote{恩宠虽有区别，却是同一的圣神；服务的职务虽有区别，却是同一的主；功效虽有区别，却是同一的天主，在一切人身上行一切事 \BibleRef{格前 12:3}}。\par

% structure: p0011 source=h2 target=subsection reason=relative-heading-rank; base=h2

\subsection{五、圣神的工作仍在教会中继续}

亲爱的诸位，借着这些以及无数其他天主的圣言权威所光照的证据，让我们同心合意地被激发去尊敬圣神降临节，欢欣鼓舞，光荣圣神；借着祂，整个教会得以圣化，每一个理性的灵魂得以苏醒；祂是信德的启发者，知识的导师，爱德的泉源，贞洁的封印，以及一切权能的原因。愿信友们的心灵喜乐，因为普世万民以各种语言宣认唯一的天主，圣父、圣子及圣神，并且祂临在的标志——以火的形象显现——仍在祂的工作和恩赐中延续。因为真理之神亲自以祂光明的光辉照耀祂荣耀的居所，不愿祂的殿宇中有任何黑暗或冷淡。并且，也是借着祂的援助和教导，我们之中建立了守斋与施舍的净化。因为这个可敬的日子之后，有一个最有益的习惯，所有圣人一直认为对他们最有益，我们以牧者的关怀劝勉你们勤行此习惯，以便在过去的日子里若因疏忽大意而沾染了任何污点，能通过守斋的纪律得到补赎，并通过虔敬的热忱得以纠正。因此，让我们在星期三和星期五守斋，并在星期六照惯例为此目的守夜，借着我们的主耶稣基督，祂与圣父及圣神，永生永王，万世无穷。阿们。\par

\SourceRecord{Translated by Charles Lett Feltoe.；Nicene and Post-Nicene Fathers, Second Series；Vol. 12.；Edited by Philip Schaff and Henry Wace.；Buffalo, NY: Christian Literature Publishing Co.,；1895.；<http://www.newadvent.org/fathers/360375.htm>.}

% structure: p0001 source=h1 target=section reason=page-title

\section{讲道第七十七篇}

\Emph{论五旬节，第三讲}\par

% structure: p0003 source=h2 target=subsection reason=relative-heading-rank; base=h2

\subsection{一、圣神的工作并非始于五旬节，而是延续至今，因为天主圣三在行动与意志上是一体的}

今天这普世尊崇的庆节，亲爱的诸位，因圣神的降临而成为神圣。这降临发生在主复活后第五十天，正如所盼望的，降在宗徒及信众身上。这盼望之所以存在，因为主耶稣曾应许祂要来，并非此时才首次成为圣人内心的居停主人，而是要激励已奉献给祂的心更趋炽热，并以更丰盈的恩赐充满它们；祂是增加而非开始祂的恩赐，并非因更丰富的恩惠而初次运作。因为圣神的尊威从不与圣父及圣子的全能分离；天主的治理在安排万物时所完成的一切，都出自整个圣三的眷顾。其中，慈悲与仁爱、审判与公义是同一的；在意志无分歧之处，行动亦无分裂。因此，圣父所光照的，圣子也光照，圣神也光照；尽管受差遣者、差遣者与应许者各有位格，但合一与圣三同时启示给我们，使我们明白：那拥有平等、不容孤寂的本质，属于同一实体，而非同一位置。\par

% structure: p0005 source=h2 target=subsection reason=relative-heading-rank; base=h2

\subsection{二、圣三的每一位都参与了我们的救赎}

因此，在不可分割的天主性完美合作下，某些事由圣父完成，某些由圣子，某些由圣神，这特别属于我们救赎的安排和得救的方式。因为人若按天主的肖像与模样受造，保有了其本性的尊严，没有被魔鬼的诡计欺骗而因私欲违犯所定的律法，那么世界的造物主就不会成为受造物，永恒者就不会进入时间领域，与天主圣父平等的天主圣子也不会取了奴仆的形象和罪恶肉身的相似。但正如经上说：\Quote{因魔鬼的嫉妒，死亡进入了世界\BibleRef{智 2:24}}，而被俘虏的人类，除非那一位为我们辩护，即那无损其尊威而成为真人、且唯独无罪者，否则无法被释放。因此，圣三的慈悲将我们恢复的工作分给自身，使圣父被平息，圣子作平息，圣神点燃。因为那些要得救的人也须在自身方面做些什么，藉将心转向救赎主而脱离仇敌的统治，正如宗徒所说：\Quote{天主派遣了祂圣子的圣神到我们心中，呼喊：\InnerQuote{阿爸，父呀！}\BibleRef{迦 4:6}}，\Quote{哪里主的圣神在那里，那里就有自由\BibleRef{格后 3:17}}，\Quote{除非在圣神内，没有人能称耶稣为主\BibleRef{格前 12:3}}。\par

% structure: p0007 source=h2 target=subsection reason=relative-heading-rank; base=h2

\subsection{三、这种职能的分配并不损害圣三的合一}

因此，亲爱的诸位，在恩宠的引导下，如果我们忠诚而明智地理解圣父、圣子、圣神的特定工作，以及在我们恢复中三者共同的工作，就必无疑会如此接受那藉贬抑和肉身为我们所成就的事，以致不会认为圣三的唯一而同等的荣耀有任何不相称之处。因为尽管没有心智足以思想天主，没有口舌足以谈论天主，但无论人的理智在天主圣父的本性上领悟了什么，除非对祂的独生子和圣神也持有同一不变的真理，否则我们的沉思便是不忠的，为肉身的侵扰所蒙蔽；并且那似乎对圣父的正确结论也会丧失，因为若合一不被持守，整个圣三就被遗弃；任何因不平等而不同者，都不能真正称为一。\par

% structure: p0009 source=h2 target=subsection reason=relative-heading-rank; base=h2

\subsection{四、思想天主时，必须摒弃一切物质概念}

因此，当我们定意明认圣父、圣子及圣神时，务要远离有形之物的形状、受造在时间中的年龄、以及一切物质的身体和地方。让那在空间中伸展、被界限围住、以及不总是处处且完整的东西，从心中驱逐出去。对圣三一体的概念必须摒弃间隔或等级的想法；如果一个人对天主有任何相称的思念，切莫不敢将之归于其中的任何一位，仿佛更荣耀地归于圣父，而不归于圣子和圣神。真正的虔诚并非将圣父置于独生子之前：对圣子的侮辱就是对圣父的侮辱；对一位的贬损就是对两者的贬损。因为既然祂们的永恒与天主性都是共同的，圣父便不被认为是全能而不可改变的，如果祂生了比自己小的，或藉拥有一个从前没有的而获益。\par

% structure: p0011 source=h2 target=subsection reason=relative-heading-rank; base=h2

\subsection{五、基督作为人小于圣父，作为天主则同等}

主耶稣确实向祂的门徒们说，正如福音选读中读到的那样：\Quote{如果你们爱我，你们就必定欢喜，因为我往父那里去，因为父比我大；}但那些耳朵，常常听见\Quote{我与父原是一体}和\Quote{看见我，就是看见了父}这些话，便接受了这一说法，并不以为天主性有何差异，也不将它理解为那与圣父同永恒、同本质的本体。因此，人因圣言降生成人而得的提升，也交托给了圣宗徒们；他们因主将要离开他们而忧伤，却因自身尊严的提升而被激励走向永恒的喜乐；祂说：\Quote{如果你们爱我，}\Quote{你们就必定欢喜，因为我往父那里去：}这就是说，如果你们以全备的知识，看出因着以下事实而赐予你们的荣耀：我出于天主圣父，却也从一位人类母亲而生；我是不可见的，却使自己成为可见的；我是永恒的，却取了\Quote{天主的形象}，接受了\Quote{奴仆的形象}，——那么，\Quote{你们就会欢喜，因为我往父那里去。}因为这次升天是为你们而提供的，你们的卑微在我内被提升到诸天之上的地方，坐在圣父的右边。但我，与圣父同是那圣父之所是，与我的父不可分离地同在；我从祂那里到你们这里来，并没有离开祂；同样，我从你们那里回到祂那里，也不会舍弃你们。所以，你们应当欢喜，\Quote{因为我往父那里去，因为父比我大。}因为我已使你们与我自己结合，并成了人子，好使你们有能力成为天主的子女。因此，虽然我在两种形态中是同一的，但在那使我与你们相同的方面，我比圣父小；而在那使我与圣父不可分的方面，我甚至比我自己大。所以，让那比圣父小的本性往圣父那里去，好让肉身能到圣言常在那里，并使天主教会唯一的信德相信：祂——作为人，我们并不否认祂较小——作为天主，却与圣父同等。\par

% structure: p0013 source=h2 target=subsection reason=relative-heading-rank; base=h2

\subsection{六、圣子与圣父所有的这种平等，圣神也拥有}

因此，亲爱的诸位，让我们鄙视那些不虔敬的异端者虚妄而盲目的诡诈，他们以对这句话的歪曲解释自夸，当主说：\Quote{凡父所有的一切，都是我的，\BibleRef{若 16:15}}，他们并不明白，凡敢否认圣子的，也就是从圣父那里夺去；他们在人间事务上竟如此愚昧，以致认为，因为圣子承担了我们的本性，所以圣父的所有物就不再属于祂的独生子了。天主身上的仁慈并不削弱能力，祂对所爱之受造物的和好也不减损永恒的光荣。圣父所有的，圣子也有；圣父和圣子所有的，圣神也有，因为整个天主圣三是唯一的天主。但这信德不是属世智慧的发现，也不是人意见的确认：是独生子亲自教导了它，圣神亲自建立了它；关于圣神，不得形成不同于关于圣父和圣子的概念。因为尽管祂不是圣父，也不是圣子，但祂与圣父和圣子不可分离：正如祂在天主圣三中有自己的位格，同样祂在天主性中与圣父和圣子有同一实体，充满万有，包含万有，并与圣父和圣子统御万有；愿尊荣和光荣归于祂，直到永远。阿们。\par

\SourceRecord{Translated by Charles Lett Feltoe.；Nicene and Post-Nicene Fathers, Second Series；Vol. 12.；Edited by Philip Schaff and Henry Wace.；Buffalo, NY: Christian Literature Publishing Co.,；1895.；<http://www.newadvent.org/fathers/360377.htm>.}

% structure: p0001 source=h1 target=section reason=page-title

\section{第七十八篇讲道}

\Emph{（论圣神降临期斋戒，第一篇）}\par

% structure: p0003 source=h2 target=subsection reason=relative-heading-rank; base=h2

\subsection{一、从宗徒时代至今，自我克制是抵御魔鬼攻击的最佳防御}

亲爱的诸位，今日的庆节因圣神的降临而被祝圣，随后，如你们所知，是一个庄严的斋戒，这斋戒作为医治灵魂和身体的救恩性制度，我们必须以虔诚的遵守来持守。因为当宗徒们被充满所应许的能力，真理之神进入他们的心时，我们毫不怀疑，在其他天堂教义的奥秘中，这种属灵自我克制的纪律首先是在\OriginalTerm{护慰者}{Paraclete}的推动下被想到的，以便因斋戒而圣化的心灵更适合领受将要赐予他们的\OriginalTerm{圣油}{chrism}。基督的门徒有全能援助的保护，初生教会的首领们藉着圣神的临在，受圣父与圣子的整个神性的守护。但是，面对迫害者威胁的攻击，面对不敬神者的可怕呼喊，他们不能用肉体的力量或娇纵的肉身来战斗，因为取悦外在的事物最伤害内在的人，一个人的肉体越受制伏，理性灵魂就越纯洁。\par

% structure: p0005 source=h2 target=subsection reason=relative-heading-rank; base=h2

\subsection{二、诱惑者在那些学会这些策略的人面前攻击失败}

因此，那些以他们的榜样和传统教导了教会所有儿女的导师们，以神圣的斋戒开始了基督徒战争的初步训练，以便在必须与属灵的邪恶作战时，他们可以披上克己的铠甲，用以杀死恶习的诱因。因为无形的仇敌和非物质的敌人，如果我们不被任何肉体的贪欲所缠绕，就不会对我们有任何力量。伤害我们的欲望确实常常在诱惑者内活跃，但如果他在我们内找不到任何进攻的有利地势，他就会被解除武装，变得无力。但是，谁被这脆弱的肉身所包围，置身于这死亡的身体之中，即使他取得了很大确定的进步，现在能如此确定自己的安全，以至于相信自己免于所有诱惑的危险呢？虽然天主的恩宠每日赐予祂的圣人们胜利，但祂并未除去战斗的机会，因为这正是我们护佑者仁慈的一部分，祂总是设计让我们不断变化的本质可以赢得一些东西，免得它因战斗的结束而自夸。\par

% structure: p0007 source=h2 target=subsection reason=relative-heading-rank; base=h2

\subsection{三、因此，这斋戒在圣神降临节后非常适时}

因此，在神圣喜乐的日子之后——我们曾将这些日子用于纪念主从死者中复活然后升天的荣耀——并在领受圣神的恩赐之后，规定了斋戒作为有益且必要的实践，以便倘若在节日的喜乐中，因疏忽或混乱而出现了过分的放纵，就可以以严格克己的补救来纠正，这必须更加谨慎地实行，以便在这一天神圣地赐予教会的恩典常在我们内。因为我们成了圣神的宫殿，比以往更多地被神圣的活水浇灌，我们不应被任何贪欲征服，也不应被任何恶习占据，以使神圣大能的居所不被任何污染玷污。\par

% structure: p0009 source=h2 target=subsection reason=relative-heading-rank; base=h2

\subsection{四、通过善用这斋戒，我们将赢得天主的恩宠}

而且，这确实对所有人都是可能的，只要有天主的帮助和引导，只要我们通过斋戒的净化和仁慈的慷慨，努力从罪恶的污秽中解脱，并在爱德的果实上富足。因为无论花费在喂养穷人、治愈病人、赎回囚犯或任何其他虔敬行为上的一切，都不会减少，反而会增加；信友的仁爱所花费的，在天主眼中绝不会失去，因为一个人施舍救济的，他为自己积存了奖赏。因为\Quote{怜悯人的人是有福的，因为天主将怜悯他们}\BibleRef{玛 5:7}；缺失将不被记念，在真实宗教的临在被证实的地方。因此，在星期三和星期五，让我们斋戒，并在星期六在至福宗徒伯多禄面前守夜，藉他的祈祷，我们确信将从精神的仇敌和身体的敌人中获得解脱；藉着我们的主耶稣基督，祂与圣父及圣神，永生永王，世世无尽。阿们。\par

\SourceRecord{Translated by Charles Lett Feltoe.；Nicene and Post-Nicene Fathers, Second Series；Vol. 12.；Edited by Philip Schaff and Henry Wace.；Buffalo, NY: Christian Literature Publishing Co.,；1895.；<http://www.newadvent.org/fathers/360378.htm>.}

% structure: p0001 source=h1 target=section reason=page-title

\section{\WorkTitle{第八十二篇讲道}}

\Emph{在宗徒伯多禄与保禄庆日（六月二十九日）}\par

% structure: p0003 source=h2 target=subsection reason=relative-heading-rank; base=h2

\subsection{一、罗马城之所以崇高，全赖这两位宗徒}

极可爱的诸位，全世界确实参与所有神圣的周年纪念，对同一信仰的忠诚要求：凡是为众人救恩所记载的事，都应普遍欢庆。但是，除了今天这庆节从普世获得的尊崇之外，它在我们城中更应受到特殊和独特的欢欣，以便在他们殉道之日，在宗徒之长完成光荣结局的地方，喜乐得以满盈。因为正是这些人，基督福音之光藉着他们照耀了你，啊，罗马；也藉着他们，你原是谬误的导师，却成了真理的门徒。他们是你的圣父和真正牧者，使你有资格被列入天上的国度，并在远比那以热心奠定你最初城墙的人更美好、更吉祥的征兆下建造了你；其中一位以他的名字给你命名的人，却用他兄弟的血玷污了你。正是这些人将你提升到如此荣耀，使你成为圣洁的邦国、特选的子民、王家的司祭民族\BibleRef{伯前 2:9}，并通过真福伯多禄的圣座成为世界的首脑，你因崇拜天主而获得的统治权比世俗政府更为广阔。因为尽管你因许多胜利而壮大，并将你的统治扩展到陆地和海上，但你的战争劳作所征服的，不及基督的和平所征服的。\par

% structure: p0005 source=h2 target=subsection reason=relative-heading-rank; base=h2

\subsection{二、罗马帝国的扩张是神圣计划的一部分}

因为良善、公义、全能与永不向人类吝惜祂仁慈的天主，并且一直以最丰盛的恩惠教导所有人认识祂自己，却以更隐秘的旨意和更深厚的爱，怜悯了流浪者自愿的盲目和向恶的倾向，派遣了祂\OriginalTerm{同等同永的圣言}{co-equal and co-eternal Word}。这圣言降生成人，如此将天主性与人性结合，以至于祂将自己的本性降至极致，从而将我们的本性提升到最高。然而，为使这不可言喻的恩宠的结果传遍全世界，天主的眷顾预备了罗马帝国，其发展达到如此范围，以致所有民族都紧密相连。因为天主计划的工作特别要求许多王国在同一个帝国下联合起来，以便对世界的宣讲能够迅速到达所有人民，当他们被置于一个国家的统治之下时。然而，那个国家，虽然统治了几乎所有民族，却不知其壮大的根源，反而被所有民族的错误所奴役，并自以为极大地培育了宗教，因为它不拒绝任何谬误。因此，它藉着基督获得解放就更为奇妙，因为它曾被撒殚如此牢牢捆绑。\par

% structure: p0007 source=h2 target=subsection reason=relative-heading-rank; base=h2

\subsection{三、关于十二宗徒分散时，圣伯多禄被派往罗马}

因为当十二宗徒藉着圣神领受了说万国方言的能力之后，将世界分派给自己，承担起以福音教导世界的任务时，真福伯多禄，宗徒团的首领，被任命到罗马帝国的堡垒，以便那为万民的救恩而彰显的真理之光，能从头本身更有效地散布到整个世界的身体。哪一个民族当时没有代表住在这城中？哪一个民族不知道罗马所学到的？就在这里，必须粉碎哲学的教条，驱散世俗智慧的愚妄，驳斥对魔鬼的崇拜，根除所有偶像崇拜的亵渎，因为最顽固的迷信在此汇集了各处所发明的一切错误。\par

% structure: p0009 source=h2 target=subsection reason=relative-heading-rank; base=h2

\subsection{四、圣伯多禄的爱战胜了前来罗马的恐惧}

真福宗徒伯多禄，你毫不畏惧来到这城中；而与你分享荣耀的宗徒保禄，仍在忙于管理其他教会时，你进入了这猛兽咆哮的森林、这深邃风暴的海洋，比你在海面行走时更加勇敢。你曾在盖法家中被大司祭的使女吓到，却不惧怕世界的主母罗马。难道克劳狄的权力比比拉多的判决或犹太人的野蛮愤怒更小？尼禄的残暴更少吗？因此，是爱的力量战胜了恐惧的理由：你不认为那些你承诺去爱的人是可惧的。但你这种无畏的爱之情感，想必在那时已经孕育，当时你对主的爱的宣认，藉着\OriginalTerm{三次询问的奥秘}{the mystery of the thrice-repeated question}得到了确认。对你的这份热诚所要求的，无非是将你自身所富足的食物，施予喂养你所爱的祂的羊群。\par

% structure: p0011 source=h2 target=subsection reason=relative-heading-rank; base=h2

\subsection{五、圣伯多禄蒙天主眷顾，为伟大使命作了准备}

你的信心也因许多奇迹征兆、许多恩宠的恩赐、许多能力的明证而增加。你已教导了从受割礼者中相信的百姓；你已在安提约基雅建立了教会，基督之名号首次在那里兴起；你已在\Place{本都}、\Place{迦拉达}、\Place{卡帕多细亚}、\Place{亚细亚}和\Place{比提尼亚}，教导了福音讯息的规诫；并且，对工作的成功毫无疑虑，深知你生命的短暂，你将基督十字架的战利品带入罗马的堡垒，藉着天主的预定，大能的荣耀和许多苦难的光荣与你同行。\par

% structure: p0013 source=h2 target=subsection reason=relative-heading-rank; base=h2

\subsection{六、从圣伯多禄和圣保禄的血中诞生了许多高贵殉道者}

也来到那里的，是你的蒙福的兄弟宗徒\Person{保禄}，\Quote{蒙拣选的器皿 \BibleRef{宗 9:15}，} 外邦人的特殊导师，在\Person{尼禄}统治下，一切无辜、一切谦逊、一切自由都处于危险之时，他与你联合。\Person{尼禄}的狂怒因各种恶习的过度而燃起，将他抛入如此疯狂的烈火熔炉中，使他成为第一个以全面迫害攻击基督信徒名号的人，仿佛天主的恩宠能因圣人的死亡而熄灭，而圣人最大的收益正是通过轻看这短暂的尘世生命而获得永恒的福乐。\Quote{珍贵的，} 因此，\Quote{在主的眼中，圣人的逝世是珍贵的：} 任何程度的残忍都无法摧毁建立在基督十字架奥迹之上的宗教。迫害不会减少教会，反而增加教会；主的田地被披上越发丰盛的庄稼，而单独落下的谷粒发芽并百倍增产。因此，从这两颗天播的种子产生了多么庞大的后裔，这由成千上万蒙福的殉道者所证明，他们与宗徒的胜利媲美，身穿紫袍、闪耀红光，遍布全城，好像用无数宝石编成的一顶冠冕加冕了它。\par

% structure: p0015 source=h2 target=subsection reason=relative-heading-rank; base=h2

\subsection{七、二者功劳不可分}

亲爱的诸位，对于天主为我们树立作为忍耐的榜样和信德的确证的这个群体，在纪念诸圣时理应处处欢欣，但对于这两位宗徒的卓越，我们更应以更大的喜乐正确地夸耀，因为天主的恩宠在教会肢体中将他们提升到如此高的地位，使他们如同身体的双眼之明光，而身体的头是基督。关于他们的功绩和美德，非言语所能尽述，我们不应加以区别，因为他们在蒙召上平等，在劳苦上相似，在死亡上无别。但是，正如我们自己经验所证明，我们的祖先所坚持的，我们相信并确知，在这尘世的一切劳苦中，我们必须时常藉着特殊代祷者的祈祷获得天主的仁慈，使我们因宗徒的功德而被提升，正如我们因自己的罪而被压垮。因我们的主\Person{耶稣基督}，等等。\par

\SourceRecord{Translated by Charles Lett Feltoe.；Nicene and Post-Nicene Fathers, Second Series；Vol. 12.；Edited by Philip Schaff and Henry Wace.；Buffalo, NY: Christian Literature Publishing Co.,；1895.；<http://www.newadvent.org/fathers/360382.htm>.}

% structure: p0001 source=h1 target=section reason=page-title

\section{讲道第八十四篇}

\Emph{论对纪念的忽视}\par

% structure: p0003 source=h2 target=subsection reason=relative-heading-rank; base=h2

\subsection{一、罗马的教会人士有忘记过去的审判与仁慈、对天主忘恩负义的危险}

在场人数之少本身已表明，亲爱的诸位，那曾使全体信众为纪念我们受惩罚与得释放之日而聚集一起感谢天主的虔诚热心，在此最近一次几乎完全被忽视了；这令我心中极为悲伤与恐惧。因为人很容易对天主忘恩负义，因忘记祂的恩惠而对惩罚不感悲伤，对释放也不感喜悦。因此我害怕，亲爱的诸位，恐怕先知的那句话斥责这样的人，说：\Quote{你鞭打了他们，他们却不悲伤；你责罚了他们，他们却拒绝接受纠正。\BibleRef{耶 5:3}}因为在那些对天主的侍奉如此冷漠的人身上，能看到什么改善呢？说来惭愧，但不可缄默：花在魔鬼身上的比花在宗徒身上的更多，疯狂的表演比受祝福的殉道更吸引人。是谁使这座城市恢复安全？是谁从被掳中拯救了它？是看赛马的人的娱乐，还是圣人们的关怀？诚然，是藉圣人们的祈祷，天主的不悦的判决被转移了，使我们这些本应受义怒的人，得以被保留下来获得宽恕。\par

% structure: p0005 source=h2 target=subsection reason=relative-heading-rank; base=h2

\subsection{二、让他们及时利用天主的忍耐，归向祂}

我恳求你们，亲爱的，让救主的那些话触动你们的心，祂曾以祂仁慈的大能洁净了十个癞病人，说他们中只有一个回来感谢祂\BibleRef{路 17:18}：这无疑意味着，虽然那些忘恩负义者获得了身体的健康，但他们在这一虔敬本分上的缺失，却源于内心的不虔敬。因此，亲爱的诸位，为使这不感恩的印记不致落在你们身上，你们应归向主，纪念祂曾屈尊在我们中所行的奇事；并将我们的释放归因于——并非如不敬者所猜想的星象影响，而是全能天主的不可言喻的仁慈，祂曾屈尊软化了狂暴野蛮人的心——你们当以信德的一切精力，投身于如此宏恩的纪念之中。严重的忽视必须以更大的悔改标记来补偿。让我们使用那饶恕了我们的祂的仁慈，以改善我们自己，使圣伯多禄及所有圣人——他们在诸多患难中常与我们相近——屈尊为你们向仁慈的天主恳求，藉我们的主耶稣基督。阿们。\par

\SourceRecord{Translated by Charles Lett Feltoe.；Nicene and Post-Nicene Fathers, Second Series；Vol. 12.；Edited by Philip Schaff and Henry Wace.；Buffalo, NY: Christian Literature Publishing Co.,；1895.；<http://www.newadvent.org/fathers/360384.htm>.}

% structure: p0001 source=h1 target=section reason=page-title

\section{讲道集第八十五篇}

\Emph{圣老楞佐殉道者庆日（八月十日）讲道}\par

% structure: p0003 source=h2 target=subsection reason=relative-heading-rank; base=h2

\subsection{殉道者的榜样至为宝贵}

诸位至爱的弟兄，一切德行的巅峰与一切义德的圆满，皆生于爱天主及爱近人之爱；然而，在这爱中，没有比在真福殉道者身上更显明、更光辉的了。他们因效法主的爱情，并因受苦的相似，而亲近为众人死的主耶稣基督。诚然，主救赎我们所用的爱，无人能以善意相匹，因为一个注定一日必死的人为义人而死是一回事，而那位不受罪债束缚者为不义之人舍命是另一回事\BibleRef{罗 5:7}。然而，殉道者也对众人大有裨益，因为赐予他们勇气的主，借此表明死亡的刑罚与十字架的痛苦，对祂的任何追随者皆不足畏惧，且可为多人所效法。因此，若无一善人独善其身，亦无一智人独益其智，而真德之性如此，即使其光照耀之处众多人脱离黑暗谬误，则教导天主子民，没有比殉道者的楷模更有用了。辩才可使祈求易达，推理可有效劝服；然而榜样胜于言语，实践比训诲更具教育意义。\par

% structure: p0005 source=h2 target=subsection reason=relative-heading-rank; base=h2

\subsection{圣人殉道事迹述略}

真福殉道者圣老楞佐在这至为卓越的教导方式中何其光荣坚强，今日因他的苦难而成为纪念日，甚至他的迫害者也能感受到，当他们发现他那奇妙勇气——主要生于对基督的爱——不仅自身不屈服，还以他的坚忍榜样坚固他人。原来，当外邦统治者的狂怒正猛烈攻击基督最精选的肢体，并特别攻击那些属于司铎品级者时，邪恶迫害者的怒火倾泻于执事老楞佐身上。他不仅在举行神圣礼仪上卓越，而且在管理教会财产上也出众。迫害者妄想从一人的捕获中获得双倍掠物：因为若迫使他交出圣宝藏，也就将他逐出真宗教的范围。于是，这个如此贪财又如此仇恨真理的人，以双重武器武装自己：以贪婪掠取黄金，以不敬俘获基督。他要求那位无伪的圣所守护者，将他贪心所觊觎的教会财富献给他。但圣执事指给他看：他储存在何处？他指着许多贫穷圣人的队伍，在供食供衣中，他储存了一笔不能失落的财富，并且因金钱已用于如此神圣的事业而更加安全。\par

% structure: p0007 source=h2 target=subsection reason=relative-heading-rank; base=h2

\subsection{苦难描述续篇}

因此，挫败的掠夺者恼怒，并燃起对那如此使用财富的宗教的仇恨，决定劫掠一个更大的宝藏，即夺去他所富有的神圣寄托，因他查不出他手中存有固体的钱财。他命令老楞佐背弃基督，并准备以可怕的酷刑来考验执事的坚强勇气；当最初的酷刑一无所获时，更残酷的刑罚接踵而至。他的四肢被许多割伤撕裂，然后被命令放在铁架子上用火烤，那铁架子本身已足够热以烧他，四肢还时常翻动，以使痛苦更甚，死亡更缓。\par

% structure: p0009 source=h2 target=subsection reason=relative-heading-rank; base=h2

\subsection{圣老楞佐已战胜其迫害者}

你一无所获，你一事无成，残酷的残暴者。他的肉身从你的计谋中释放，当老楞佐离世升天时，你被击败了。基督的爱火不能被你的火焰战胜，外面燃烧的火不及内在燃烧的火炽烈。迫害者，你只是在愤怒中为殉道者效劳；你只是在增加痛苦的同时加大了奖赏。因为你的狡计所设想的，哪一样不是归荣耀于胜利者？甚至刑具也被视为凯旋的一部分。\newline{} 因此，诸位至爱的弟兄，让我们以属灵的喜乐欢欣，并因这杰出的人在主内的幸福结局而夸耀。主\Quote{在祂的圣者中显为奇妙}，在他们身上祂赐予我们支持和榜样，并将他的荣耀广传于普世，以致从日出之地到日落之处，他执事之光的照耀遍及世界，罗马因老楞佐而闻名，如同耶路撒冷因斯德望而尊贵。我们信赖时常藉他的祈祷和转求得到助佑，使凡愿在基督内虔敬生活的，都如宗徒所说：\Quote{凡愿在基督内虔敬生活的，必将遭受迫害}\BibleRef{弟后 3:12}，我们得以被爱德的精神所坚强，并由坚定信德的恒忍所巩固，以克服一切诱惑。藉我们的主耶稣基督，等等。\par

\SourceRecord{Translated by Charles Lett Feltoe.；Nicene and Post-Nicene Fathers, Second Series；Vol. 12.；Edited by Philip Schaff and Henry Wace.；Buffalo, NY: Christian Literature Publishing Co.,；1895.；<http://www.newadvent.org/fathers/360385.htm>.}

% structure: p0001 source=h1 target=section reason=page-title

\section{讲道 88}

\Emph{论七月斋戒，第三篇。}\par

% structure: p0003 source=h2 target=subsection reason=relative-heading-rank; base=h2

\subsection{古代先知所宣布的斋戒，今日仍然必要}

亲爱的诸位，宗教斋戒在赢得天主的仁慈和更新人类软弱之命运方面有何益处，我们从圣先知们的陈述中得知；他们宣告，天主的公义——以色列子民因自己的邪恶屡次招致其惩罚——只能藉斋戒平息。因此，岳厄尔先知警告他们说：\Quote{上主你们的天主这样说：你们应全心归向我，禁食、哭泣、悲哀；应撕裂你们的心，而不是你们的衣服，归向上主你们的天主，因为祂是宽仁慈悲的，缓于发怒，富于慈爱，且极其慈悲。}又说：\Quote{制定斋戒，宣告治疗，召集人民，圣化教会。}这项劝勉在我们这个时代也必须遵行，因为这些治疗的方法必须由我们同样宣告，以便在遵守古代圣化的过程中，基督徒的虔诚能获得犹太人的悖逆所失去的。\par

% structure: p0005 source=h2 target=subsection reason=relative-heading-rank; base=h2

\subsection{公共的敬礼比私人的具有更高品格}

但是，对天主法令的尊重总会带来特别的祝福，无论我们自愿服务的范围如何；因此，公开宣告的庆典比那些建立在私人制度上的庆典具有更高的品格。因为个人凭自己判断所施行的克制，只涉及教会某一特定部分的益处；而整个教会所守的斋戒，则无人被排除在普遍净化之外；当天主子民的心在一种共同的圣洁服从之举中汇聚在一起时，他们便变得最为强大；当基督徒军队的营垒中处处是同样的战斗和防御准备时，尽管残忍的敌人暴怒不安，到处布下隐藏的陷阱，但若他发现无人疏于戒备、无人沉溺于懒惰、无人不积极于虔敬之事，他便不能捕捉或伤害任何人。\par

% structure: p0007 source=h2 target=subsection reason=relative-heading-rank; base=h2

\subsection{九月的斋戒以这种公开的方式召唤我们自我改正}

因此，亲爱的诸位，我们如今正是被第七个月的隆重斋戒邀请到这种不可战胜的团结力量中，以便我们摆脱世俗的忧虑和地上的挂虑，将灵魂提升到主面前。由于这种努力虽然始终必要，我们却不能永久持守；且因人性的软弱，我们常常从崇高的事物跌回尘世之事；至少在这些为了我们的改正而最有益地制定的日子里，让我们从世俗事务中抽身，为增进我们永恒的福祉而偷得一点时间。\Quote{我们众人在许多事上都有过失\BibleRef{雅 3:2}。}尽管藉着天主每日的恩赐，我们从各种污秽中被洗净，但仍有不慎的灵魂常常沾染更深的污点，需要更大的关注来洗净，更强的努力来消除。而罪过的完全消除，是在整个教会献上同一祈祷和同一宣认时获得的。因为如果主应许对两三个人的圣洁虔诚祈求，必予成全\BibleRef{玛 18:19}，那么对成千上万同心合意、异口同声祈祷的子民，又有什么会拒绝呢？\par

% structure: p0009 source=h2 target=subsection reason=relative-heading-rank; base=h2

\subsection{在天主眼中，财物与行为的共有万分宝贵}

亲爱的诸位，在主的眼中，当基督的全体子民共同履行同一种义务，所有男女各阶层以同一心意合作时，这是一件伟大而十分宝贵的事；当同样的目的激励所有人远离邪恶、行善；当天主在祂仆人的工作中受光荣，一切虔诚的源头因不断的感恩而受赞美时，饥饿者得饱食，赤身者得衣穿，患病者得探望；人寻求的不是自己的利益，而是\Quote{别人的利益}；在救济他人痛苦时，每个人都尽量善用自己的资财；并且容易找到\Quote{乐捐的人}，在那里人的作为只受其能力范围的限制。藉着天主的恩宠，\Quote{祂在一切内工作一切}，信友们的利益和功德都共同享有。因为那些收入不同的人，却能思想相同；当一个人为别人的慷慨而欢欣时，他的情感使他与那拥有不同消费能力的人处于同等地位。在这样的团体中，没有混乱或分歧，因为整个身体的所有肢体都同意于一个坚强的虔诚目的，那以他人财富为荣的人并不因自己的贫穷而羞愧。因为每一部分的卓越，是整个身体的光荣；当我们都被天主的圣神引导时，不仅我们自己所做的事属于我们，而且我们为之欢喜的别人所做的事也属于我们。\par

% structure: p0011 source=h2 target=subsection reason=relative-heading-rank; base=h2

\subsection{那么，让我们善用这个机会}

那么，亲爱的诸位，让我们以良善意志的和谐一致的主意，持守这至圣的合一于其蒙福的完整，并投入这严肃的斋戒。没有人被要求做任何艰难、苛刻的事，也没有人被迫超过自己的力量，无论是在禁食的规律上还是在施舍的数量上。各人知道他能做什么、不能做什么：让每个人都按公正和合理的标准评估自己，付出他的份额，使慈悲的祭献不因忧愁而献上，也不被视为损失。用在虔敬工作上的支出，应多到足以使心灵成义、洗净良心，简言之，使施者和受者都得益。那灵魂真是有福的，且确实值得钦佩，它在行善的爱中不担心资源的匮乏，并不怀疑那曾赐予他以往所花费的，仍会给他金钱使用。但因为很少人有这般伟大的心胸，而各人不放弃对自己家人的照顾确是一件虔敬的事，所以我们不贬低更成全的人，而为你们立下这条普遍规则，劝你们按自己能力的大小去遵行天主的命令。因为喜乐与慈善的人相称，他应以如此方式管理他的慷慨，使穷人因所得的帮助而欢乐，同时家庭的需要也不至受损。\Quote{那赐种子给播种者，也赐粮食作食物的，必使你们的种子增多，并增加你们正义的果实。\BibleRef{格后 9:10}}因此，我们应在星期三和星期五守斋；并在星期六一起在至福的宗徒伯多禄面前守夜，藉他的功劳和祈祷，我们确信天主的仁慈将藉我们的主耶稣基督，赐予我们一切所需，祂永生永王，直到永远。阿们。\par

\SourceRecord{Translated by Charles Lett Feltoe.；Nicene and Post-Nicene Fathers, Second Series；Vol. 12.；Edited by Philip Schaff and Henry Wace.；Buffalo, NY: Christian Literature Publishing Co.,；1895.；<http://www.newadvent.org/fathers/360388.htm>.}

% structure: p0001 source=h1 target=section reason=page-title

\section{讲道集 第九十篇}

（论七月斋期，第五篇）\par

% structure: p0003 source=h2 target=subsection reason=relative-heading-rank; base=h2

\subsection{一、我们必须不断寻求宽恕，因为我们时刻可能犯罪}

至爱的诸位，我们宣布神圣的七月斋期，是为操练共同的虔敬，并以慈父般的劝勉激励你们，以你们的遵守使那原本属于犹太人的斋期成为基督徒的。因为无论在何时，寻求天主的仁慈都应伴随着身心两方面的克己，这符合新约与旧约，因为没有任何事比一个人自己审判自己，并因知道自己从未无罪而不断恳求宽恕，更能获得天主的应允。人性本身有其缺陷，并非由造物主植入，而是由犯过者所招致，并通过生育的律法传给后代，以致从可朽坏的身体生出那也能败坏灵魂的东西。因此，虽然内在的人如今已在基督内重生，并脱离了捆绑的锁链，它却与肉体不断争斗，在试图克制虚妄欲望时忍受抵抗。在这斗争中，不易获得如此完全的胜利，以至于那些必须断绝的习惯不再妨碍我们，那些必须被消灭的恶习不再伤害我们。无论心灵以何种智慧和谨慎作为外在感官的审判官，即使在统治和限制肉体贪欲的努力中，诱惑总是近在咫尺。因为谁能如此超脱肉体的快乐或痛苦，以至于不被外在所喜悦或折磨的事物影响？喜悦与哀伤与人不可分离：人没有一部分免于愤怒的点燃、快感的压倒、痛苦的挫败。当统治者和被统治者同样容易受到相同激情的影响时，哪里会有对罪的回避？主正确地喊道：\Quote{心神固然切愿，但肉体却软弱。}\par

% structure: p0005 source=h2 target=subsection reason=relative-heading-rank; base=h2

\subsection{二、基督亲自成为道路，祂命令我们行走其上}

免得我们因绝望而陷入完全的不作为，祂应许说，那因人的无能为力而对人不可能的事，将因天主的德能而成为可能：\BibleQuote{因为通往生命的门是窄的，路是狭的} \BibleRef{玛 7:14} ，没有人能踏上它，没有人能前进一步，除非基督亲自成为道路，解开了接近的困难：如此，旅途的安排者成为使我们能够完成旅程的方法，因为祂不仅施加劳苦，也带领我们到达安息的港湾。因此，在祂内我们找到忍耐的榜样，在祂内我们拥有永生的希望；因为\BibleQuote{如果我们与祂一同受苦，也必与祂一同受光荣} \BibleRef{弟后 2:12} ，因为正如宗徒所说：\BibleQuote{那说自己住在祂内的，就应当照祂所行的去行} \BibleRef{若一 2:6} 。否则我们便是徒有虚名和外表，若不跟随祂的足迹，我们却以祂的名夸耀；那些足迹，若我们只爱祂命令我们爱的，就不会令我们厌烦，反而会解救我们脱离一切危险。\par

% structure: p0007 source=h2 target=subsection reason=relative-heading-rank; base=h2

\subsection{三、对天主的爱与对世俗之爱的对比}

因为有两种爱，一切愿望由此而生，它们在性质上不同，正如它们的来源不同。有理性的灵魂不能没有爱，它要么是爱天主，要么是爱世俗。在天主的爱中没有过度，但在世俗的爱中一切都是有害的。因此，我们必须不可分离地依附于永恒的宝藏，而对暂时的事物则要像过客一样使用，为使我们这些急于返回故乡的旅人，在这世上所遇见的一切美好事物都成为旅途的助力，而非羁留的罗网。因此，蒙福的宗徒宣告：\BibleQuote{时间是短促的：从此以后，有妻子的要像没有一样；哭泣的要像不哭泣的；欢乐的要像不欢乐的；置买的要像一无所得的；享用这世界的要像不享用的，因为这世界的局面正在逝去。} \BibleRef{格前 7:29} 但是，由于世界以其外貌、丰盛和多样吸引我们，不容易离开它，除非在可见事物的美中，爱造物主胜过爱受造物；因为祂说：\BibleQuote{你当全心、全灵、全意、全力爱上主你的天主} \BibleRef{玛 22:37} ，祂愿意我们在注视时从祂爱的束缚中解脱出来。当祂也将爱近人联于这条诫命时，祂命令我们效法祂的仁慈，使我们爱祂所爱的，做祂所做的。因为虽然我们是\Quote{天主的田地，天主的建筑物}，并且\Quote{栽种的算不得什么，浇灌的算不得什么，只在那使之生长的天主}，但祂在一切事上要求我们的服务和事奉，愿意我们作祂恩宠的管家，使那带着天主肖像的人能承行天主的旨意。为此缘故，我们在主祷文中虔诚地祈求：\Quote{愿祢的国来临，愿祢的旨意奉行在人间，如同在天上。}在这话中我们祈求的，岂不是愿天主征服那些祂尚未征服的人，并如祂在天上使天使成为祂旨意的仆役，同样在地上也使人为仆役吗？在寻求这事时，我们爱天主，也爱近人；我们内的爱只有一个对象，因为我们愿仆人服侍，愿主统治。\par

% structure: p0009 source=h2 target=subsection reason=relative-heading-rank; base=h2

\subsection{四、对天主的爱藉善工而培养}

因此，诸位亲爱的弟兄，这种排除属世之情的心境，因行善的习惯而得以坚固，因为良心必然因善行而喜悦，并乐意做它因做了而欢喜的事。如此，守斋得以实行，施舍得以慷慨，正义得以维护，祈祷得以频仍，而个人的愿望成为众人的共同愿望。劳作培育忍耐，温柔熄灭愤怒，仁爱践踏仇恨，不洁的欲望被神圣的渴望所杀死，贪吝被慷慨所驱逐，而累赘的财富成为善行的工具。但是，因为魔鬼的网罗即使在这样的事态中也未止息，极其合宜地，在一年中的某些时节，我们活力的更新被预备好了：尤其现在，当那贪图现世利益者可能因天气温和与土地肥沃而自夸，并将收成存入大仓房时，可能对自己的灵魂说：\BibleQuote{你拥有许多财物，吃吧，喝吧，} 愿他留意天主声音的斥责，并听到它说：\BibleQuote{愚昧的人啊，今夜他们要向你索取你的灵魂，而你预备的一切，它们将归给谁呢？}\BibleRef{路 12:19} 这应是智者最为焦虑的考虑，以便——既然此生的日子短少，其期限不定——死亡永不突然临到他，并且他明知自己必死，能完全准备好迎接结局。因此，为使这既能圣化我们的身体，又能更新我们的灵魂，让我们在星期三与星期五守斋，并在星期六与最可敬的宗徒\Person{伯多禄}一起守夜，他的祈祷将帮助我们藉着我们的主基督获得我们神圣愿望的实现，祂与圣父及圣神永生永王，至于无穷之世。阿们。\par

\SourceRecord{Translated by Charles Lett Feltoe.；Nicene and Post-Nicene Fathers, Second Series；Vol. 12.；Edited by Philip Schaff and Henry Wace.；Buffalo, NY: Christian Literature Publishing Co.,；1895.；<http://www.newadvent.org/fathers/360390.htm>.}

% structure: p0001 source=h1 target=section reason=page-title

\section{讲道集 91}

\Emph{论七月斋期，第六篇}\par

% structure: p0003 source=h2 target=subsection reason=relative-heading-rank; base=h2

\subsection{一、克己必须包括灵魂与身体的纪律}

亲爱的诸位，天主眷顾无不在一切事上扶助信徒的虔诚。因为连世界的元素也为圣洁的身心修炼服务，因为昼夜和月份的明确多样变化为我们开启了诫命的不同篇章，因而季节在某种意义上也向我们言说着神圣制度所命定的事。故此，既然一年的循环已将第七个月带回到我们中间，我确信你们的心灵已受到激发，要持守这庄严的斋戒；因为你们已从经验中得知，这预备如何净化人的外在和内在，使人在禁绝合法之事时，更能抗拒不合法之事。但亲爱的诸位，不要将你们的克己计划局限于身体的苦修或仅仅减少食物。因为这德行的更大益处属于灵魂的贞洁，它不仅摧毁肉体的情欲，还轻视世俗智慧的虚妄，正如宗徒所说：\Quote{你们要谨慎，免得有人用哲学和虚空的诡诈，按人的传统，把你们掳去\BibleRef{哥 2:8}。}\par

% structure: p0005 source=h2 target=subsection reason=relative-heading-rank; base=h2

\subsection{二、尤其要禁绝异端，无论是欧提克斯的异端还是聂斯多略的异端}

因此，我们必须节制饮食，但更要禁绝谬误，使那不受任何肉体快乐支配的心灵也不被任何谎言所俘获：因为正如过去的日子，在我们这个时代也不乏真理的敌人，他们胆敢在天主教会内挑起内战，诱使无知者附和他们的不敬教义，并因那些被他们从基督身体上分离的人而自夸人数的增长。因为，还有什么是如此反对先知、如此抵触福音、如此违背宗徒的教导，如同宣讲在由玛利亚所生的主耶稣基督内只有一个本性，且与永恒圣父同永恒，不受时间影响呢？如果只承认人的本性，那么拯救的天主性在哪里？如果只有天主性，那么被拯救的人性在哪里？然而公教信德，它抵挡一切谬误，同时也驳斥这些亵渎之言，谴责聂斯多略（\OriginalTerm{聂斯多略}{Nestorius}），他将天主性与人性分开；并斥责欧提克斯（\OriginalTerm{欧提克斯}{Eutyches}），他在天主性中消灭了人性。因为真天主子，他自己是真天主，拥有与圣父和圣神的合一与平等，也屈尊成为真人，并在童贞圣母受孕后，没有与她的肉身和生育分离，如此将人性结合于自己，以致他永远不变地是天主；如此将天主性赋予人，以致不是毁灭人，而是借着荣耀提升人。因为他成了\Quote{奴仆的形体}，并没有停止成为\Quote{天主的形体}，他不是一位与另一位结合，而是在两者中是一位，因此自从\Quote{圣言成了血肉}以来，我们的信德不受任何境遇变迁的困扰，无论在大能的奇迹中，还是在受苦的屈辱中，我们都相信他既是天主，也是人，且既是人，也是天主。\par

% structure: p0007 source=h2 target=subsection reason=relative-heading-rank; base=h2

\subsection{三、降生成人的真理由圣体盛宴和天主设立的施舍制度所证实}

亲爱的诸位，请全心全意宣认这信仰，并弃绝异端者的邪恶谎言，使你们的斋戒和施舍不因任何谬误的感染而被玷污：因为我们的祭献之奉献和慈悲之施舍之所以洁净，乃在于行这些事的人明白他们所行之事。因为主说：\BibleQuote{你们若不吃人子的肉，不喝他的血，在你们内便没有生命\BibleRef{若 6:53}}，你们应当如此参与圣餐桌，以致对基督身体和血的真实性（\OriginalTerm{基督的身体和血}{Christ's Body and Blood}）毫不怀疑。因为口中所领受的正是信德所相信的，那些对所领受之物有争议的人，他们的\Quote{阿们}是虚妄的。但当先知说：\Quote{眷顾贫穷和贫乏的人是有福的}时，那在穷人中间分发衣服和食物的人值得称赞，因为他知道自己在穷人身上穿衣服和喂养基督：因为主亲自说：\BibleQuote{你们对我这些最小兄弟中的一个所做的，就是对我做的\BibleRef{玛 25:40}}。因此，基督是唯一的，真天主和真人（\OriginalTerm{真天主和真人}{True God and True Man}），在自己方面是富有的，在我们方面是贫穷的，接受恩赐也分施恩赐，与必朽者分享，并是死者的赋予生命者（\OriginalTerm{赋予生命者}{Quickener of the dead}），以致\Quote{因耶稣之名，一切在天上的、地上的和地底下的，都要屈膝，并一切唇舌都要宣认：主耶稣基督是在天主圣父的光荣中\BibleRef{斐 2:10}}，他与圣神一起，永生永王，直到永远。阿们。\par

\SourceRecord{Translated by Charles Lett Feltoe.；Nicene and Post-Nicene Fathers, Second Series；Vol. 12.；Edited by Philip Schaff and Henry Wace.；Buffalo, NY: Christian Literature Publishing Co.,；1895.；<http://www.newadvent.org/fathers/360391.htm>.}

% structure: p0001 source=h1 target=section reason=page-title

\section{讲道 95}

\Emph{真福八端讲道，圣} \BibleRef{玛 5:1-9}\par

% structure: p0003 source=h2 target=subsection reason=relative-heading-rank; base=h2

\subsection{一、主题介绍}

当我们的主耶稣基督，亲爱的，宣讲天国的福音，并在整个\Place{加里肋亚}治愈各种疾病时，祂大能作为的名声已传遍整个\Place{叙利亚}：也有大批从\Place{犹太}全地来的群众涌向天上的医生。\BibleRef{玛 4:23}因为人的无知，对未见之事相信迟缓，对未知之事希望迟延，那些将要受教于神圣学问的人，需要被身体的好处和可见的奇迹所唤醒：这样当他们实际体验到祂慈善的能力时，就不会对祂教导的健全性有怀疑。因此，为叫主能用外在的治愈作为内在疗法的引子，并在医治身体之后，在灵魂中施行治愈，祂从周围的人群中分开，登上邻近一座山的僻静处，在那里召唤祂的宗徒们到祂跟前，为能从那个奥秘座位的高处，教导他们更崇高的教义，从那个地方和行动的本质中表明，祂就是那曾与\Person{梅瑟}说话、尊荣他的人：那时确实带着更可怕的正义，但现在带着更神圣的仁慈，为使所应许的得以应验，当先知\Person{耶肋米亚}说：\BibleQuote{看，日子将到，我要为以色列家和犹大家订立一个新的盟约。那些日子以后，主说，我要将我的法律放在他们的理智中，并写在他们的心上。}因此，那曾与\Person{梅瑟}说话的，也向宗徒们说话，\OriginalTerm{圣言}{Word}敏捷的手写下并存放了新盟约的奥秘在门徒们的心内：没有像古时那样的密云包围祂，也没有可怕的声音和闪电吓跑人们不敢靠近山，而是安静地、自由地，祂的讲论传到那些站在旁边的人的耳中：为使法律的严苛让位于\OriginalTerm{恩宠}{grace}的温柔，而\Quote{义子的精神}能驱散奴役的恐惧。\par

% structure: p0005 source=h2 target=subsection reason=relative-heading-rank; base=h2

\subsection{二、讨论谦虚的真福}

基督教导的本质，由祂自己的神圣言论所证实：为叫那些希望达到永恒幸福的人，能理解通往那崇高幸福的上升步骤。他说：\Quote{神贫的人是有福的，因为天国是他们的。\BibleRef{玛 5:3}}如果说\Quote{贫穷的人是有福的}而没有加上任何说明，那么也许会让人怀疑祂所说的是哪一种穷人：这样一来，这贫穷本身似乎就足以赢得天国，而许多穷人却因艰难沉重的需要而遭受这贫穷。但当祂说\Quote{\Emph{神贫的人}是有福的}时，祂表明天国必须归于那些因心灵的谦卑而非财物的匮乏而被推荐的人。然而，不能怀疑，这种谦卑的拥有对穷人来说比富人更容易获得：因为顺从是缺乏者的伴侣，而心灵的傲慢与财富同在。尽管如此，即使在许多富人身上，也能找到那种精神，它使用自己的富足，不是为了增加骄傲，而是为了慈善工作，并认为在帮助他人困苦时所花费的才是最大的收益。各种类型和等级的人都能分享这种德行，因为人在意愿上可以平等，尽管在命运上不平等：那些在属灵财富上相等的人，在现世财物上有多大差异并不重要。因此，不贪恋现世事物、不寻求以世俗财富增加自己、而热心积累天上财富的贫穷，是有福的。\par

% structure: p0007 source=h2 target=subsection reason=relative-heading-rank; base=h2

\subsection{三、圣经中的谦虚榜样}

这种崇高心灵的谦卑，首先由宗徒们（继主之后）给我们立下了榜样，他们一听到天上师傅的声音，就毫无差别地舍弃了所有的一切，从捕鱼迅速转变为捕人的渔夫，并通过效法他们的信德，使许多人变得像他们一样，那时，与教会那些首生的儿女一起，\Quote{众信徒都是一心一意。\BibleRef{宗 4:32}}因为他们舍弃了全部财物和所有物，通过最虔诚的神贫，用永恒的财富充实了自己，并依照宗徒们的宣讲，乐于一无所有，在基督内拥有一切。因此，蒙福的宗徒伯多禄上圣殿时，被一个跛足者求施舍，他说：\Quote{银子和金子，我没有；但把我所有的给你：因纳匝肋人耶稣基督的名字，你起来行走罢！\BibleRef{宗 3:6}}还有什么比这谦卑更崇高？还有什么比这贫穷更富有？他没有金钱的库存，但他有本性的恩赐。那个母亲从胎中生出就跛足的人，被伯多禄用一句话治好了；那个没有给凯撒的一枚钱币上刻有凯撒像的人，却在人身上恢复了基督的肖像。借着这宝藏的财富，不仅是那个恢复行走能力的人得到了帮助，而且还有五千人，当时因这治愈的奇迹，在宗徒的劝勉下相信了。那个穷人，没有东西可以给求他的人，却赐予了如此丰厚的圣宠恩赐，以至于，正如他使一个人双脚直立，同样他治愈了成千上万相信者的内心，使他们\Quote{跳跃如鹿}，在基督内，他找到了在犹太人的不信中跛行的人。\par

% structure: p0009 source=h2 target=subsection reason=relative-heading-rank; base=h2

\subsection{四、讨论哀恸的真福}

在这最幸福的谦逊宣告之后，主接着说道：\Quote{哀恸的人是有福的，因为他们要受安慰。\BibleRef{玛 5:4}} 亲爱的诸位，这种哀恸，即应许有永恒安慰的哀恸，并非此世的苦难；在全人类悲伤中所倾泻的哀哭，并不能使任何人得福。圣洁叹息的理由、有福眼泪的缘由是截然不同的。虔诚的悲伤为罪过——无论是他人的还是自己的——而哀恸；它不为天主正义所施行的事哀恸，而是为人的邪恶所犯下的罪而悲伤，在那里，作恶的人比受恶的人更应被哀悼，因为不义之人的恶行将他投入惩罚，而义人的忍耐却引领他走向光荣。\par

% structure: p0011 source=h2 target=subsection reason=relative-heading-rank; base=h2

\subsection{五、温良之福}

接着主说：\Quote{温良的人是有福的，因为他们要承受土地。} 对温良和柔和、谦卑和克制、并准备忍受一切伤害的人，应许了土地作为他们的产业。这不应被视为微小或廉价的产业，仿佛它与我们天上的居所不同，因为正是这些人才被理解为进入天国的人。因此，应许给温良的人、要赐给柔和的人作为产业的土地，就是圣人的肉身，这肉身因谦逊的赏报将在幸福的复活中改变，披上不朽的光荣，不再与精神相悖，而是与灵魂的意志完全合一和谐。因为那时外在的人将成为内在的人和平无瑕的产业；那时，沉浸于瞻仰天主的理智将不受人性软弱的任何障碍所阻碍，也不再需要说：\Quote{那腐朽的身体压迫灵魂，这地上的帐幕压住那多虑的心思。\BibleRef{智 9:15}} 因为土地不会与它的居住者争斗，也不会对统治者的管辖有任何反抗。因为温良的人将在永久的和平中拥有它，他们的权利将一无所失，\Quote{当这可朽坏的穿上了不可朽坏的，这可死的穿上了不可死的\BibleRef{格前 15:53}} 之时，使他们的危险变为赏报，负担成为荣耀。\par

% structure: p0013 source=h2 target=subsection reason=relative-heading-rank; base=h2

\subsection{六、渴慕正义之福}

此后主接着说：\Quote{饥渴慕义的人是有福的，因为他们要得饱饫。\BibleRef{玛 5:6}} 这饥饿、这饥渴寻求的并非任何属肉体、属世的事物；它渴望以正义的美膳得到饱足，并希望被接纳进入一切最深的奥秘，并以主自己为满足。渴求这食粮、渴望这饮品的灵魂是有福的；它若从未尝过其甘甜，必定不会寻求它。但是，听到先知的精神对他说：\Quote{请你们体验，请你们观看：上主是何等的甘饴；} 它从上主领受了一部分甘甜，便燃烧起对最纯洁愉悦的爱，以致摒弃一切暂时之物，以最大的热忱投入到对正义的饮食中，并领悟那第一条诫命的真理，它说：\Quote{你应全心、全灵、全力爱上主你的天主：} 因为爱天主无非就是爱正义。最后，正如在那段经文中，对近人的关怀与对天主的爱相连，同样，在这里，慈悲之德也与对正义的渴慕相连，于是说道：\par

% structure: p0015 source=h2 target=subsection reason=relative-heading-rank; base=h2

\subsection{七、怜悯之福}

\Quote{怜悯人的人是有福的，因为他们要受怜悯。\BibleRef{玛 5:7}} 基督徒啊，你要认清你智慧的价值，并理解你被召唤得什么赏报，以及你必须通过什么样的训练方法才能达到。怜悯希望你怜悯，正义希望你正义，为的是使造物主在祂的受造物身上被看见，天主的肖像通过模仿的线条反映在人心这面镜子中。行善者的信德无忧无虑：你将拥有一切你所渴望的，并将无止境地获得你所爱的。由于通过你的施舍，一切对你都是纯洁的，因此你也要达到那随之而应许的福，主接着说道：\par

% structure: p0017 source=h2 target=subsection reason=relative-heading-rank; base=h2

\subsection{八、心地纯洁之福}

\Quote{心里洁净的人是有福的，因为他们要看见天主。\BibleRef{玛 5:8}} 亲爱的诸位，为他准备了如此伟大赏报的人，其幸福是多么巨大啊！那么，使心洁净，岂不是竭力追求上面所提到的那些德行吗？看见天主的幸福是何等的大，什么心思能领会，什么口舌能言说呢？然而，这将在人的本性转变后实现，以致不再\Quote{借着镜子}，也不\Quote{借着谜语}，而是\Quote{面对面\BibleRef{格前 13:12}} 地看见天主本身\Quote{祂是怎样的\BibleRef{若一 3:2}}，这是没有人能看见的；并通过永恒默观的不可言喻的喜乐，获得\Quote{眼睛未曾看见，耳朵未曾听见，人心也未曾想到的}。这幸福被应许给心地纯洁的人，是合宜的。因为真光的光辉不能被不洁净的视线看见；那将是明亮洁净之心的幸福的东西，对玷污的心灵却是惩罚。因此，应避开世俗虚幻的迷雾，洁净你内心的眼睛，除去一切邪恶的污秽，使这视野得以自由地饱享天主这伟大的显现。因为我们明白，为了达到这一点，下面的话引导我们。\par

% structure: p0019 source=h2 target=subsection reason=relative-heading-rank; base=h2

\subsection{九、缔造和平之福}

\Quote{缔造和平的人是有福的，因为他们要称为天主的子女\BibleRef{玛 5:9}。} 这种真福，亲爱的，不属于任何种类的一致与和谐，而是属于宗徒所说的那种：\Quote{与天主和睦；}以及达味先知所说的：\Quote{爱慕祢法律的人必享广大的平安，他们毫无绊脚石。}这种平安，即使是最密切的友情和最相似的思想，如果不符合天主的旨意，也不能真正获得。相似恶欲、结伙犯罪、同流合污，都不能挣得这种平安。爱世俗不与爱天主相合，不愿意离开今世之子的人，不能加入天主的子女的行列。相反，那些心思常与天主同在的人，\Quote{竭力保守圣神的合一在和平的纽带中\BibleRef{弗 4:3}，}从不违反永律，并发出信德的祈祷：\Quote{祢的旨意奉行在人间如同在天上\BibleRef{玛 6:10}。}这些是\Quote{缔造和平的人，}他们完全同心合意，全然和谐，并要被称为子女\Quote{属于天主并与基督同为承继者\BibleRef{罗 8:17}，}因为对天主之爱与对邻人之爱的明证将是这样：我们不再遭受灾祸，不惧怕任何冒犯，而是结束一切试炼的争战，在天主最完美的平安中安息，借着我们的主，他与圣父和圣神，永生永王，世世无穷。亚孟。\par

\SourceRecord{Translated by Charles Lett Feltoe.；Nicene and Post-Nicene Fathers, Second Series；Vol. 12.；Edited by Philip Schaff and Henry Wace.；Buffalo, NY: Christian Literature Publishing Co.,；1895.；<http://www.newadvent.org/fathers/360395.htm>.}
